
\chapter{文本/图像到3D生成}
\label{ch:08}

\paragraph*{本章阅读锚点}
本章的锚点是生成先验。原始 3DGS 从真实图像优化一个场景；生成方法借助扩散模型或多视角生成模型补足视角、形状或纹理，但需要区分生成假设和真实观测。


一句话就能生成3D模型？一张照片就能重建完整场景？这曾是计算机图形学的终极幻想，如今正被3D Gaussian Splatting与扩散模型（diffusion model）的联姻快速变为现实。3DGS的显式表示和实时渲染能力为生成式3D建模提供了天然骨架，而文本到图像扩散模型积累的海量视觉先验则为3D生成注入了前所未有的想象力。

两条路线正在交汇。一条以Score Distillation Sampling（SDS，分数蒸馏采样）为代表，将2D扩散模型的知识"蒸馏"到3D表示中。另一条以前馈重建网络为代表，从多视角扩散输出直接推断3D高斯。然而，核心挑战依然严峻。SDS优化的速度瓶颈与过平滑问题、多视角生成中的几何不一致性、从文本/单图到完整3D场景的语义鸿沟，都是尚未完全解决的开放问题。

本章将系统梳理3DGS生成领域的技术全景，从基础的SDS蒸馏方法出发，覆盖文本到3D生成、图像到3D重建、3D生成模型与GAN、4D动态生成等多个方向。

\begin{figure}[htbp]
\centering
\resizebox{\textwidth}{!}{\begin{tikzpicture}[
  route/.style={rectangle, rounded corners=5pt, draw=#1, fill=#1!10,
    text width=3cm, minimum height=0.9cm, align=center, font=\small},
  method/.style={rectangle, rounded corners=3pt, draw=#1, fill=#1!6,
    font=\scriptsize, text width=2.8cm, align=center, minimum height=0.6cm},
  arr/.style={-{Stealth[length=4pt]}, techgray!60},
  >=Stealth
]

  \node[route=techgray, text width=4cm, font=\small\bfseries] (root) at (37.7, 5.6) {文本/图像到3D生成路线};
  \node[route=techblue] (b0) at (12.4, 3.1) {文本到3D生成};
  \node[method=techblue] (b0_l0) at (1.5, 0.6) {SDS蒸馏\\{\tiny DreamGaussian, MicroDreamer, GaussianDreamer, LucidDreamer, ConsDreamer}};
  \node[method=techblue] (b0_l1) at (5.7, 0.6) {前馈快速生成\\{\tiny Turbo3D, GradeADreamer, Dreamer XL}};
  \node[method=techblue] (b0_l2) at (10.0, 0.6) {多视角扩散\\{\tiny CSD, 3D-Adapter, TD-Attn, VIST3A}};
  \node[method=techblue] (b0_l3) at (14.7, 0.6) {大规模/高斯原生\\{\tiny MGM, TexGaussian, DirectTriGS, DreamPolisher}};
  \node[method=techblue] (b0_l4) at (19.4, 0.6) {场景级生成\\{\tiny DreamScene360, HoloDreamer, FastScene, CompGS, LayoutDreamer}};
  \node[route=techgreen] (b1) at (31.7, 3.1) {图像到3D生成};
  \node[method=techgreen] (b1_l0) at (23.2, 0.6) {前馈重建\\{\tiny LeanGaussian, GSD, DiffusionGS, DSplats}};
  \node[method=techgreen] (b1_l1) at (27.1, 0.6) {优化式重建\\{\tiny LiftImage3D, GSV3D, Hi3D, Vista3D}};
  \node[method=techgreen] (b1_l2) at (31.7, 0.6) {大规模基础模型\\{\tiny MVGamba, 3D-Adapter}};
  \node[method=techgreen] (b1_l3) at (36.2, 0.6) {高质量重建\\{\tiny GS-RGBN, PhyCAGE, GaussianZoom}};
  \node[route=techorange] (b2) at (43.7, 3.1) {3D生成模型/GAN};
  \node[method=techorange] (b2_l0) at (40.1, 0.6) {3D GAN\\{\tiny GS Decoder, GSGAN, F3D-Gaus}};
  \node[method=techorange] (b2_l1) at (43.7, 0.6) {扩散模型\\{\tiny DiffGS, DiffSplat, GGS}};
  \node[method=techorange] (b2_l2) at (47.7, 0.6) {自回归/VAE\\{\tiny ARGS, LTM3D, UVGS, Can3Tok}};
  \node[method=techorange] (b2_l3) at (51.9, 0.6) {蒸馏/适配器\\{\tiny DD3G, 3D-Adapter, Zero-1-to-G}};
  \node[route=techpurple] (b3) at (57.7, 3.1) {4D动态生成};
  \node[method=techpurple] (b3_l0) at (55.9, 0.6) {扩散蒸馏\\{\tiny DreamGaussian4D, Diffusion4D, STP4D}};
  \node[method=techpurple] (b3_l1) at (59.5, 0.6) {时空重建\\{\tiny EG4D, PLA4D, 4DSTR}};
  \node[method=techpurple] (b3_l2) at (63.2, 0.6) {角色/化身\\{\tiny CharacterShot, GenMOJO}};
  \draw[arr] (root) -- (b0);
  \draw[arr] (root) -- (b1);
  \draw[arr] (root) -- (b2);
  \draw[arr] (root) -- (b3);
  \draw[arr] (b0) -- (b0_l0);
  \draw[arr] (b0) -- (b0_l1);
  \draw[arr] (b0) -- (b0_l2);
  \draw[arr] (b0) -- (b0_l3);
  \draw[arr] (b0) -- (b0_l4);
  \draw[arr] (b1) -- (b1_l0);
  \draw[arr] (b1) -- (b1_l1);
  \draw[arr] (b1) -- (b1_l2);
  \draw[arr] (b1) -- (b1_l3);
  \draw[arr] (b2) -- (b2_l0);
  \draw[arr] (b2) -- (b2_l1);
  \draw[arr] (b2) -- (b2_l2);
  \draw[arr] (b2) -- (b2_l3);
  \draw[arr] (b3) -- (b3_l0);
  \draw[arr] (b3) -- (b3_l1);
  \draw[arr] (b3) -- (b3_l2);
\end{tikzpicture}}
\caption{文本/图像到3D生成技术路线图}
\label{fig:ch08_gen3d_taxonomy}
\end{figure}

%=====================================================================
\section{文本到3D生成}
\label{sec:08_1}

文本到3D生成是3D内容创建中最具想象力的方向。用户仅需输入一句自然语言描述，系统便自动生成具有合理几何与丰富纹理的3D资产。本节首先介绍SDS及其改进方法，这一优化范式是绝大多数文本到3D方法的数学基础。随后按技术路线组织文本到3D生成的核心方法，涵盖SDS优化路线、前馈/快速生成路线、多视角扩散路线、大规模3D生成路线、高斯原生生成路线以及场景级文本到3D生成。

\subsection{SDS与Score蒸馏}
\label{sec:08_1_1}

本节介绍 Score Distillation Sampling 的数学原理及其核心改进方法。

\subsubsection{Score 蒸馏基础方法}

本路线介绍分数蒸馏采样（SDS）的数学原理及其在3DGS中的核心改进方法，是后续所有文本到3D生成路线的共同基础。

Score Distillation Sampling（SDS）由Poole（2022.09）~\cite{2209_14988}在DreamFusion中提出，是文本到3D生成的基石。其核心思想是：不必训练3D扩散模型，而是利用预训练的2D文本到图像扩散模型为3D表示的优化提供梯度信号。直觉上，SDS做的事情如下。在每一个优化步骤中，先从随机相机视角渲染3D模型得到2D图像，再将图像加噪后送入扩散模型，利用扩散模型预测的噪声与实际添加噪声的差异来更新3D参数，使渲染结果逐渐逼近扩散模型所学习的图像分布。

形式化地，给定3D表示参数$\theta$（在3DGS中即所有高斯基元的参数集合）、渲染函数$g$、预训练扩散模型$\epsilon_\phi$，SDS的梯度为：

\begin{equation}
\label{eq:sds}
\nabla_\theta \mathcal{L}_{\text{SDS}}(\theta) = \mathbb{E}_{t,\epsilon}\left[w(t)\left(\hat{\epsilon}_\phi\left(g(\theta);y,t\right) - \epsilon\right)\frac{\partial g(\theta)}{\partial \theta}\right]
\end{equation}

其中$t$是噪声水平，$\epsilon \sim \mathcal{N}(0,I)$是添加的高斯噪声，$\hat{\epsilon}_\phi$是扩散模型预测的噪声，$y$是文本嵌入，$w(t)$是时间步权重。白话注解：这个梯度指向的方向是"让渲染图像看起来更像是扩散模型会生成的图像"。当扩散模型认为渲染图像中的噪声模式与文本描述不匹配时，梯度就会推动3D参数朝正确的方向更新。

SDS虽然是文本到3D的开创性方法，但存在两个突出问题。（1）优化效率低。每次参数更新都需要一次扩散模型前向传播，整个优化过程可能需要数万步。（2）生成质量受限。过平滑的纹理、过饱和的颜色以及"Janus问题"（物体多个面出现同一特征，即模型无法区分正面和背面）是常见的视觉缺陷。围绕这些问题，研究者提出了多种改进方案。

\paragraph{SDS 分数蒸馏与 3DGS 结合}

DreamGaussian（2023.09）~\cite{2309_16653}是首个将SDS优化与3DGS表示结合的方法，将3DGS从重建任务推广到生成任务。本方法修改了3DGS的优化目标与致密化策略：用SDS蒸馏损失替代传统的重建损失驱动高斯参数更新，而3DGS的体渲染管线（前到后深度排序、$\alpha$-混合）保持不变。此前，3DGS仅用于静态场景重建：以SfM（Structure from Motion，运动恢复结构）点云为初始化，用多视角ground-truth图像的$L_1$+$\text{D-SSIM}$损失监督参数优化。DreamGaussian将这一范式推广到生成任务，不依赖任何多视角观测，仅凭单张图像或文本提示即可从零生成3D资产。

\textbf{设计动机。}DreamFusion已证明SDS+NeRF可以实现文本到3D生成，但其效率瓶颈严重。NeRF的体渲染需要对每条射线采样数百个点并前向传播MLP，单次优化需要数万步，总耗时约10分钟。DreamGaussian的核心洞察是：3DGS的基于tile的光栅化渲染比NeRF的射线采样快一个数量级（推理速度约是NeRF的$\sim$100倍），因此用3DGS替换NeRF作为SDS优化的载体，可以从渲染效率层面大幅加速生成过程。然而，直接将重建流程中的3DGS参数配置搬到生成场景并不奏效。重建任务有确定性的ground-truth梯度信号来驱动致密化，而SDS的随机性监督使标准密度控制失效。因此，必须重新设计致密化策略以适配生成场景。

\textbf{第一阶段：SDS优化3D高斯。}在此阶段，DreamGaussian遵循公式~\eqref{eq:sds}的SDS梯度框架，但针对3DGS表示做了三项关键适配。

第一项是致密化策略调整。在标准3DGS重建管线中，高斯基元由SfM点云初始化，致密化由累积位置梯度幅值触发。但在生成设置下，SDS的随机噪声预测每步都产生不同方向的梯度，导致累积梯度的统计意义失效。为此，DreamGaussian采用"少基元、高频致密化"策略——image-to-3D任务初始化$5000$个随机粒子（分布在半径$0.5$的球体内，初始不透明度$0.1$、灰色），每$100$步致密化一次；text-to-3D任务仅用$1000$个粒子，每$50$步致密化一次。梯度阈值从重建中的$0.5$降低至$0.01$（text-to-3D），以鼓励更积极的致密化来弥补初始化粒子的不足。直觉上，这好比用更少的种子起步，但更频繁地播种——因为SDS的"雨"来得不如ground-truth那样稳定。

第二项是简化颜色表示。标准3DGS使用球谐函数（spherical harmonics, SH）来建模视角相关的颜色效果。SH系数在高阶时可达$16$维，参数空间大，在SDS这种每步目标不一致的优化中容易陷入局部最优。因此DreamGaussian禁用了SH，仅保留每个高斯基元的简单RGB颜色$c \in \mathbb{R}^3$。直觉上，SH如同一个精细的调色板，需要稳定的多视角信号才能准确学习。SDS的随机监督给不了这种稳定性，所以退而求其次只用最基础的颜色。

第三项是时间步线性退火。DreamGaussian遵循DreamTime（2023）~\cite{huang2023dreamtime}的建议，将SDS中的噪声时间步$t$按线性策略从高到低退火。训练初期$t$较高，SDS添加的噪声大，扩散模型提供的是大尺度结构信号（整体形状、布局）。训练后期$t$逐渐降低，SDS聚焦于局部细节（纹理、边缘）。直觉上，这种"先粗后细"的时间步调度类似于雕刻：先用粗刀大刀阔斧，再用细刀精雕细琢。作者在消融实验中验证了这一设计的必要性。没有时间步退火时，相同迭代次数下模型收敛到更差的形状。

对于image-to-3D任务，除了SDS损失外，DreamGaussian还引入了参考视角对齐损失来约束生成结果与输入图像的一致性：

\begin{equation}
\label{eq:dg_ref}
\mathcal{L}_{\text{ref}} = \lambda_{\text{rgb}} \|I^r_\theta - I^r_{\text{gt}}\|_2^2 + \lambda_{\alpha} \|A^r_\theta - A^r_{\text{gt}}\|_2^2
\end{equation}

其中$I^r_\theta$和$A^r_\theta$分别是参考视角$r$的渲染RGB和透明度，$I^r_{\text{gt}}$和$A^r_{\text{gt}}$是输入图像及其前景mask，$\lambda_{\text{rgb}}$和$\lambda_{\alpha}$在训练过程中从$0$线性增长到$10^4$和$10^3$。白话注解：这个损失确保模型在参考视角的渲染结果与输入图像像素对齐，线性增长的权重使得SDS先自由探索3D结构，然后逐步收紧对参考视角的约束。

第一阶段的总训练步数为$500$步。渲染分辨率从$64$逐步提升到$512$，与致密化进程配合——低分辨率时高斯数量少、覆盖粗粒度几何，随着分辨率提升致密化填充细节。相机位姿在固定半径$2$（image-to-3D）或$2.5$（text-to-3D）的球面上随机采样，方位角$[-180^\circ, 180^\circ]$，仰角$[-30^\circ, 30^\circ]$，y轴FOV为$49^\circ$。

\textbf{第二阶段：纹理提取与UV细化。}经过第一阶段SDS优化后，3D高斯已经编码了几何和粗略纹理信息，但直接输出往往模糊且缺乏细节。这是因为SDS的歧义性（每步优化提供不一致的3D引导）使得高斯基元难以精确收敛。DreamGaussian观察到，延长第一阶段训练也无法解决模糊问题。

因此第二阶段将3D高斯转换为带纹理的三角网格，然后在UV空间细化纹理。网格提取的关键步骤是局部密度查询：将3D空间$(-1,1)^3$划分为$16^3$个重叠块，利用3DGS中大高斯会被分裂或剪枝的特性做块级剔除（类似tile-based culling），在每个块内查询$8^3$密度网格，最终组合为$128^3$的密度场：

\begin{equation}
\label{eq:dg_density}
d(\mathbf{x}) = \sum_i \alpha_i \exp\left(-\frac{1}{2}(\mathbf{x} - \mathbf{x}_i)^T \Sigma_i^{-1}(\mathbf{x} - \mathbf{x}_i)\right)
\end{equation}

其中$\alpha_i$是第$i$个高斯的不透明度，$\Sigma_i$是由缩放$s_i$和旋转$q_i$构建的协方差矩阵。白话注解：这个公式就是在3D空间中计算每个位置被所有高斯基元"覆盖"的密度。每个高斯是一个3D椭球，距离中心越近密度越高，不透明度越大贡献越强。达到经验阈值$1.0$的位置被Marching Cubes（行进立方体）算法提取为网格表面。

纹理初始化通过颜色反向投影（color back-projection）完成：在8个方位角$\times$3个仰角加上顶部和底部视角渲染RGB图像，每个像素通过UV坐标映射回纹理图集。排除相机空间z方向法线较小的像素以避免网格边界处的不稳定投影。

纹理细化阶段借鉴了SDEdit（2022）~\cite{2021sdedit}的图像到图像合成思路。直觉上，已有粗纹理相当于一幅模糊的草图，SDEdit式的做法是先对草图加噪（但噪声强度要有限，不能完全破坏原有内容），然后用扩散模型去噪——去噪过程在恢复图像的同时会自然地补全细节。形式化地，从任意相机视角$p$渲染粗纹理图像$I^p_\text{coarse}$，添加噪声后经扩散模型多步去噪得到细化图像：

\begin{equation}
\label{eq:dg_uv}
I^p_\text{refined} = f_\phi(I^p_\text{coarse} + \epsilon(t_s); t_s, c)
\end{equation}

其中$\epsilon(t_s)$是时间步$t_s$处的随机噪声，$c$是条件输入（image-to-3D中为相对相机位姿$\Delta p$，text-to-3D中为文本嵌入$e$），$t_s$被精心选择以限制噪声强度。细化后的图像通过像素级MSE损失$\mathcal{L}_\text{tex} = \|I^p_\theta - I^p_\text{refined}\|_2^2$优化纹理图集。这里不直接用SDS损失细化UV纹理，是因为UV纹理使用mipmap采样技术，SDS的模糊梯度传播到各mipmap层级会导致过饱和的色块伪影。第二阶段仅需约$50$步即可获得高质量纹理。

\textbf{效率与效果。}在效率方面，DreamGaussian从单张图像生成带纹理网格仅需约2分钟（image-to-3D两阶段各约1分钟），text-to-3D因Stable Diffusion模型分辨率$512 \times 512$较大需约2分钟。相比DreamFusion约10分钟/样本的优化时间，实现了约10倍的加速。所有实验在NVIDIA V100（16GB）GPU上完成，峰值显存使用低于8GB。速度提升的根本来源是3DGS的光栅化渲染效率——无需射线采样和MLP前向传播，基于tile的并行化使得单步渲染比NeRF快约两个数量级。

\paragraph{基于分数的迭代重建}

如果说 DreamGaussian 证明了 3DGS 与 SDS 的可行性，那么 MicroDreamer（2024.04）~\cite{2404_19525} 则追问了一个更本质的问题：SDS 本身能否被重写为更高效的范式？\textbf{动机。}SDS 每次参数更新都需要一次扩散模型前向传播（NFE 过高），且损失在潜在空间计算需要反向传播通过 VAE 编码器，这两者共同导致优化速度缓慢；MicroDreamer 将"蒸馏"改写为"重建"范式来根本性地降低 NFE。MicroDreamer 的核心改动是将"蒸馏"范式改写为"重建"范式，以大幅降低 NFE（Number of Function Evaluations，函数评估次数）。该方法与 DreamGaussian 共享 3DGS 表示，但优化策略完全不同。方法指出 SDS 慢的两个根因：（1）每次3D参数更新都需要一次扩散模型前向传播（即NFE过高），整个优化需数万步扩散推理；（2）SDS损失在潜在空间计算，需要反向传播通过VAE编码器，进一步拖慢速度。直觉上，SDS像是"每走一步都要问一次老师"，而MicroDreamer改为"问一次老师后自己练习很多步"。

为此，MicroDreamer提出Score-based Iterative Reconstruction（SIR，基于分数的迭代重建）。

\textbf{动机。}SDS 的根本瓶颈在于每次 3D 参数更新都需调用扩散模型（NFE 过高）；SIR 将扩散推理与 3D 参数更新解耦，每次扩散推理服务多次 3D 参数更新，将 NFE 从数万降至约 160。SIR模仿可微分3D重建过程：在每次外循环迭代中，先用扩散模型精炼出一组多视角图像作为"伪真值"，然后在保持相机位姿不变的情况下对3D参数做多步内循环优化：

\begin{equation}
\label{eq:sir}
\mathcal{J}_{\text{SIR}}(\theta; c^{(k)}, \hat{x}^{(k)}) = \|g(\theta, c^{(k)}) - \hat{x}^{(k)}\|_1
\end{equation}

其中$c^{(k)}$是第$k$次迭代的相机位姿集合，$\hat{x}^{(k)}$是经扩散模型精炼的多视角图像。白话注解：SIR将"蒸馏"改写为"重建"。先让扩散模型生成一组一致的多视角图像，再用标准的$L_1$重建损失去拟合它们。关键优势在于：重建损失在像素空间直接计算，无需反向传播通过VAE；且每次扩散推理可服务多次3D参数更新，NFE从数万降至约160。MicroDreamer在单张A100 GPU上从3DGS生成网格仅需约20秒，速度是DreamGaussian的2倍。

\paragraph{一致性蒸馏连接分数蒸馏}

上述方法从优化策略角度改进 SDS 的效率，而 GCS（2024.07）~\cite{2407_13584} 则从理论层面重新审视蒸馏的一致性约束。

\textbf{核心洞察。}SDS 的随机采样与一致性模型的确定性采样，本质上都在追求同一个目标——让梯度始终指向同一个确定的目标，只是实现路径不同。

\textbf{动机。}SDS 的随机采样导致每步梯度指向不同方向，方差极大，优化过程在多个局部最优间振荡；一致性模型的确定性采样可消除这种随机性。GCS（Guided Consistency Sampling，引导一致性采样）核心改动是用一致性模型的确定性采样替代SDS的随机采样以减少梯度方差。

方法首先建立了理论桥梁。一致性模型（Consistency Model, CM）追求自一致性，即同一条ODE（常微分方程）轨迹上的所有点映射到同一个起点。SDS及其变体（如VSD）追求跨一致性，即渲染图像的扩散轨迹与文本条件的去噪轨迹对齐。直觉上，SDS的梯度每步指向不同的随机方向（因为每次采样的噪声不同），而GCS希望梯度始终指向同一个确定的目标。基于这一分析，GCS利用DPM-Solver-1作为一致性函数，提出紧致一致性（CC）损失来确保子轨迹内的自一致性，并配合条件引导（CG）损失来强化跨一致性，两者共同减少梯度方向的波动。

此外，方法观察到3DGS优化中渲染视角持续出现亮度过度饱和的问题。这是3DGS在SDS梯度驱动下的独特缺陷，因为高斯基元的不透明度会沿梯度方向累积。为此提出亮度均衡生成（Brightness-Equalized Generation, BEG）方案，通过对亮度信号进行均衡化处理来抑制不希望的亮度累积。

\paragraph{生成蒸馏中的自适应密度控制}

前面GCS从梯度采样策略层面减少了方差，但SDS优化中的另一类问题源于3DGS自身的密度控制机制。

\textbf{动机。}在重建任务中，位置梯度随优化收敛而衰减，可可靠识别几何欠拟合区域；但在 SDS 优化中，扩散模型每步的伪目标都在变化，导致梯度幅值不收敛且方向剧烈振荡，标准密度控制完全失效。CAdam（Context-Adaptive Moment Estimation，上下文自适应矩估计）（2026.05）~\cite{2605_20872}聚焦于SDS优化中3DGS密度控制的根本缺陷：标准致密化在随机监督下完全失效。

问题根因在于：在重建任务中，位置梯度会随优化收敛而衰减（$\|\nabla L\| \to 0$），因此累积梯度幅值能可靠地识别几何欠拟合区域。但在生成式蒸馏中，扩散模型每步的伪目标$\hat{x}_0^t$都在变化，导致两种异常动态。（1）不收敛：即使几何已视觉合理，梯度幅值仍居高不下。（2）方向方差：噪声预测$\hat{\epsilon}_\phi$的方差使梯度方向在步间剧烈振荡（称为zigzag现象）。直觉上，标准密度控制像是在风暴中用风速计判断风向，风的随机阵风会被误判为持续的信号。

CAdam将密度控制重新定义为统计信号验证问题。将每步梯度$g_t$建模为确定性几何漂移$d_t$与随机噪声$n_t$的叠加（$g_t = d_t + n_t$, $\mathbb{E}[n_t]=0$），利用梯度的一阶矩（均值）实现"干涉原理"：随机波动通过破坏性干涉相互抵消，而一致的几何漂移通过建设性干涉积累，从而将真实几何信号从生成噪声中分离。CAdam将高斯数量减少85\%--97\%，同时保持可比较的感知质量。

\subsection{文本到3D生成}
\label{sec:08_1_2}

SDS为文本到3D提供了优化框架，但生成质量、效率和可控性仍有巨大的提升空间。本节按技术路线组织53篇论文，深入分析各路线的核心创新与权衡。本节讨论6条技术路线：（1）SDS优化路线；（2）前馈/快速生成路线；（3）多视角扩散路线；（4）大规模3D生成路线；（5）高斯原生生成路线；（6）场景级文本到3D生成。

\subsubsection{SDS优化路线}

这条路线以SDS为基础优化范式，通过改进蒸馏策略、初始化方式和正则化来提升生成质量。

\paragraph{桥接 2D/3D 扩散模型}

GaussianDreamer（2023.10）~\cite{2310_08529}修改了DreamGaussian/SDS的初始化环节。它用3D扩散模型的输出替代随机初始化作为3DGS的起点，而SDS优化框架本身保持不变。

\textbf{动机。}纯随机初始化在 SDS 优化中容易陷入几何不一致的局部最优（如 Janus 问题），而 3D 扩散模型虽质量有限但天然具备 3D 一致性，可作为更优的初始化起点。设计动机源于2D与3D扩散模型各自的互补缺陷。3D扩散模型（如Shap-E（2022.12）~\cite{2212_02624}）以3D数据训练，天然具备3D一致性，但训练数据稀缺导致生成质量有限、泛化能力弱。2D扩散模型（如Stable Diffusion）拥有海量训练数据和强大泛化能力，却无法保证多视角的3D一致性。GaussianDreamer的核心思路是用3D扩散模型提供"粗骨架"，再用2D扩散模型"精雕细琢"，以3DGS作为两者之间的桥接表示。

\textbf{第一阶段：3D扩散模型初始化。}与DreamGaussian在球体内随机采样初始点不同，GaussianDreamer利用3D扩散模型$F_{3D}$（如Shap-E或文本到运动的扩散模型）根据文本提示$y$生成三角网格$m = F_{3D}(y)$，再将网格顶点和纹理颜色转化为点云$\mathcal{P}_m(\mathbf{x}_m, \mathbf{c}_m)$。然而，直接使用这些点云存在两个问题：点云位置过于稀疏（仅网格顶点），且颜色过于简单（来自3D扩散模型的粗糙纹理），不足以支撑后续SDS优化。为此，方法提出两个增强操作——\textbf{噪声点生长}（noisy point growing）和\textbf{颜色扰动}（color perturbation）。

噪声点生长的目标是在原始点云周围合理地补充密度：先计算点云$\mathcal{P}_m$的包围盒（bounding box），在包围盒内均匀采样大量候选点$\mathbf{x}_u$，然后利用KDTree搜索每个候选点到原始点云的最近邻距离——仅保留归一化距离小于$0.01$的点，确保新增点紧贴原始表面而非偏离到空间中随意位置。颜色扰动则为新增点赋予与最近邻点相近但有随机偏差的颜色，避免初始化颜色过于单一：

\begin{equation}
\mathbf{c}_r = \mathbf{c}_m + \boldsymbol{\delta}, \quad \boldsymbol{\delta} \in [0, 0.2]
\end{equation}

其中$\mathbf{c}_m$是最近邻原始点的颜色，$\boldsymbol{\delta}$是$[0, 0.2]$范围内的随机扰动向量——直觉上，噪声点生长像是在已有的稀疏骨架上"长肉"，在表面附近补充足够的点密度，而颜色扰动让这些"新肉"的颜色稍有变化，为后续优化提供更大的搜索空间。最终，原始点云与生长点云合并为$\mathcal{P}_f = \mathcal{P}_m \cup \mathcal{P}_r$，用于初始化3D高斯的位置$\boldsymbol{\mu}_b$和颜色$c_b$；不透明度$\alpha_b$初始化为$0.1$，协方差$\Sigma_b$根据最近邻点间距计算。

\textbf{第二阶段：2D扩散模型SDS优化。}在3D扩散模型提供的一致性初始化之上，方法采用标准的SDS损失（式~\ref{eq:sds}）对3D高斯参数$\boldsymbol{\theta}_i$进行优化，用预训练的2D扩散模型$F_{2D}$（Stable Diffusion）逐步丰富几何和外观细节。直觉上，第一阶段搭建了一个"大致形状正确"的3D结构，第二阶段则让2D扩散模型在保持3D一致性的前提下为其添加高质量纹理。这种两阶段设计在单张GPU上仅需约15分钟即可生成高质量的3D实例或3D虚拟形象，相比DreamFusion实现了数量级的速度提升，同时避免了从零开始优化时的3D不一致问题。

\paragraph{无领域限制的场景生成}

LucidDreamer（2023.11）~\cite{2311_13384}与GaussianDreamer同为SDS优化方法，但核心差异在于：GaussianDreamer聚焦物体级生成并依赖3D扩散模型初始化，LucidDreamer则面向场景级生成且完全不需要3D训练数据。

\textbf{动机。}场景级生成比物体级更具挑战性——场景范围大、内容多样，无法用简单几何初始化；LucidDreamer 通过"边走边看"的渐进扩展策略，利用扩散模型的 inpainting 能力逐步填充未见区域。方法提出一种域无关的管线，交替执行两个步骤。（1）Dreaming：以已有3D点云为几何引导，通过扩散模型（Stable Diffusion）的inpainting从新的相机视角生成多视角图像，再用单目深度估计（如ZoeDepth）将像素反投影到3D空间，形成新的点云。直觉上，这像是"边走边看"，从一个初始视角出发，逐步向未观测方向扩展场景。（2）Alignment：将新生成的3D点云与已有场景进行对齐融合，通过相机轨迹规划确保扩展方向的连贯性。重复Dreaming-Alignment循环直到场景覆盖足够大的范围后，将完整3D点云作为3DGS优化的初始化，通过3DGS的$\alpha$-混合渲染来填充点云中因深度估计误差产生的空隙。LucidDreamer生成的3DGS场景细节丰富且无领域约束。

\paragraph{绑定几何的高质量生成}

GaussianDreamerPro（2024.06）~\cite{2406_18462}继承GaussianDreamer的3D扩散模型初始化思路，但核心改动是将自由增长的3D高斯绑定到逐步演化的几何结构上。

\textbf{动机。}在生成任务中，3D 高斯倾向于无控制地向各个方向增长（因为 SDS 的随机监督使梯度方向不一致），导致生成质量远低于重建质量；将高斯绑定到几何结构可约束其运动范围。设计动机源于一个关键观察：在生成任务中，3D高斯倾向于无控制地向各个方向增长（因为SDS的随机监督使梯度方向不一致），导致生成质量远低于重建质量。在重建中确定性梯度指向唯一最优，而在生成中随机梯度使高斯四处扩散。为解决此问题，方法提出几何引导的三阶段框架。

第一阶段（粗几何生成）：继承GaussianDreamer，用3D扩散模型（如Shap-E）生成初始网格，但采用2D高斯（2DGS）而非3D高斯作为表示。2D高斯能更好地贴合表面，提供更准确的几何。第二阶段（基础纹理生成）：在2D高斯上用2D扩散模型通过SDS优化外观和几何，得到基本合格的3D资产。第三阶段（高质量纹理增强）：将2D高斯导出为带颜色顶点的网格，然后在网格表面初始化3D高斯并绑定到网格上。这些"几何约束"高斯只能在网格表面附近移动，从而在SDS优化中保持结构一致性。直觉上，前两阶段像是"搭骨架"，第三阶段是在骨架上"贴皮肤"并精修细节。最终输出绑定到网格的3D高斯资产，还可无缝集成到动画、组合和物理模拟等下游管线中。

\paragraph{基于整流流的蒸馏}

FlowDreamer（2024.08）~\cite{2408_05008}与GaussianDreamer共享SDS优化框架，但核心改动是替换扩散先验本身。它用整流流模型（rectified flow，如Stable Diffusion 3）替代标准扩散模型。

\textbf{动机。}标准扩散的轨迹是弯曲的，分数函数$\epsilon_\phi(x_t, t)$与时间$t$强相关，导致在不同$t$上采样时产生不同方向的梯度，造成过平滑；整流流定义直线路径，理论上不同$t$处的梯度方向一致，收敛更快。设计动机在于：标准扩散的轨迹是弯曲的，分数函数$\epsilon_\phi(x_t, t)$与时间$t$强相关，导致在不同$t$上采样时产生不同方向的梯度，造成过平滑。整流流则定义直线路径$x_t = t\epsilon + (1-t)x_0$，学习时间无关的向量场$v_\theta$，理论上不同$t$处的梯度方向一致，收敛更快。直觉上，扩散模型的去噪像在弯曲的山路上绕行，而整流流走的是直线隧道。

方法首先将SDS适配到整流流框架，推导向量场蒸馏采样（VFDS）损失：$\nabla_\theta L_{\text{VFDS}} = \mathbb{E}_{\epsilon,t}[w(t)(v_\theta(x_t,t) - (\epsilon - x))\frac{\partial x}{\partial \theta}]$。但实验发现VFDS仍导致过平滑。根因是VFDS随机采样噪声$\epsilon$，使同一图像对应多条不同的ODE轨迹，梯度方向不一致。为此提出Unique Couple Matching（UCM）损失：利用整流流的两个关键性质（可逆性和耦合性），为每个渲染图像搜索唯一的配对噪声，引导3D模型沿同一条ODE轨迹优化。FlowDreamer在3DGS和NeRF两种设置下均提升了文本细节和收敛速度，CLIP（Contrastive Language-Image Pre-training）相似度在3DGS设置下达到30.70（ViT-B-32），超越LucidDreamer的28.81。

\paragraph{变分分布映射}

DreamMapping（2024.09）~\cite{2409_05099}同样改进SDS的分布建模，但与FlowDreamer替换先验不同，核心改动是重新定义渲染图像与扩散生成图像之间的关系。

\textbf{动机。}SDS 过平滑的根因在于其无条件项将 3D 参数推向模式平均；先前方法通过构建变分分布缓解，但需要计算 UNet 雅可比矩阵，耗时 5-10 小时；DreamMapping 用轻量级网络建模渲染图像到生成图像的退化关系，将变分分布的构建从雅可比计算转化为高效网络训练。方法首先系统分析了SDS的过平滑根因：SDS的无条件项会将3D参数推向$p_\theta(x_t|y)$的模式平均，导致所有时间步的模式被平均化，产生过平滑。先前工作（如VSD、ProlificDreamer）通过构建渲染图像的变分分布来缓解此问题，但需要额外计算UNet的雅可比矩阵，耗时5-10小时。DreamMapping提出变分分布映射（VDM）——关键洞察是将渲染图像视为扩散生成图像的"退化版本"：渲染图像$x = g(\theta, c)$质量低于扩散生成图像，但两者之间存在可建模的退化关系。

VDM用一个轻量级可学习网络$M_\omega$（仅约1M参数的U-Lite）建模这个退化过程，从而将变分分布的构建从昂贵的扩散模型雅可比计算转化为高效的神经网络训练。直觉上，VDM不再问"渲染图像在扩散分布的什么位置"，而是问"渲染图像是扩散图像经怎样的退化得到的"。此外，方法发现渲染图像与生成图像的分布相关性随时间步$t$减小而降低（$t=0$时完全独立），据此提出分布系数退火（DCA）策略——以$t=300$为分界点，在$t > 300$时施加强分布约束，$t \leq 300$时减弱约束，避免低噪声水平下的过度正则化。DreamMapping使用3DGS（初始化约50000个点）作为3D表示，5000次迭代即可生成高保真3D资产。

\paragraph{缓解先验视角偏差}

ConsDreamer（2025.04）~\cite{2504_02316}不从SDS损失函数本身入手，而是聚焦于一个被忽视的根因——文本到图像（T2I）模型中的先验视角偏差（prior view bias），这一偏差是导致Janus问题的真正来源。

\textbf{动机。}此前方法多从 SDS 损失函数或优化策略入手缓解 Janus 问题，但忽视了真正的根源在于 T2I 模型本身的视角偏见；只有从条件项和无条件项两方面分别纠正视角偏差，才能从根本上解决 Janus 问题。直觉上，T2I模型在训练时看到的大多数图像都是物体的"正面"，因此即使用户提示词中没有指定视角，模型也会默认生成正面视角——当3D优化从不同角度渲染时，每个角度都被推向"正面"，导致物体出现多个正面（即Janus问题）。方法首先从数学上分析3D表示$\theta$的概率模型：$p_{3D}(\theta) = \prod_{n=1}^N p_{2D}(z_0 | c, v_n)$，其中每个视角的2D分布$p_{2D}$由T2I模型决定，而T2I模型即使在没有视角信息的提示词$c$中也会隐含视角偏好。

针对这一问题，ConsDreamer从SDS的条件项和无条件项两方面分别解决：（1）针对条件项——视角解耦模块（View Disentanglement Module, VDM）从条件提示中去除与视角无关的语义成分，注入精确的目标视角控制$v_n$，使T2I模型按指定视角而非默认视角生成；（2）针对无条件项——基于相似性的偏序损失（similarity-based partial order loss）利用物理先验约束：两个视角的方位角越接近，其渲染图像的余弦相似度应越高。形式化地，若$|\Delta\phi_{ij}| < |\Delta\phi_{kl}|$，则$\cos(\mathbf{I}_i, \mathbf{I}_j) > \cos(\mathbf{I}_k, \mathbf{I}_l)$应成立。ConsDreamer可无缝集成到多种3D表示和分数蒸馏范式中。

\paragraph{分段一致性轨迹蒸馏}

SegmentDreamer（2025.07）~\cite{2507_05256}继承GCS的一致性蒸馏路线，但核心改动是将全局一致性分解为分段一致性以解决CFG条件引导不当的问题。

\textbf{动机。}全局一致性要求 ODE 轨迹从起点到终点全程一致，对长轨迹来说约束太强；将轨迹划分为多个子分段，在每个分段内独立确保一致性，可收紧蒸馏误差上界。方法首先分析了现有CD-based方法的不平衡问题：一致性蒸馏采样（CDS）损失$L_{CDS}(\theta) = \mathbb{E}_{t,s}[c(t)\|F_\theta(z_t,t,y) - F_\theta(\hat{z}_s,s,y)\|_2^2]$侧重自一致性（同一条ODE轨迹上的映射一致），但缺少分类器分数（classifier score），无法生成语义一致的内容；GCS虽然加入了条件引导损失$L_{CG}$，但由于其使用的$\epsilon_\phi(z_t,t,y)$缺少目标时间步参数，条件引导本身存在缺陷，且蒸馏误差上界为$O(\delta_t) \cdot T$（$\delta_t = |t-s|$），在文本到3D中因$|t-s|$通常较大而效果有限。直觉上，现有方法要求ODE轨迹从起点到终点全程一致，这对长轨迹来说约束太强、难以满足。SegmentDreamer提出分段一致性轨迹蒸馏（SCTD）：将概率流ODE轨迹$[0,T]$划分为多个子轨迹$[t_0, t_1], [t_1, t_2], \ldots$，在每个子分段内独立确保一致性，从而将蒸馏误差上界从$O(\delta_t) \cdot T$收紧到$O(\delta_t) \cdot T/K$（$K$为分段数）。理论上，SCTD比全局一致性方法提供了更紧的蒸馏误差上界。

\subsubsection{前馈/快速生成路线}

本路线介绍通过训练大规模前馈重建模型实现亚秒级文本到3D生成的方法，将传统SDS优化范式的生成时间从分钟级压缩到秒级。

SDS优化需要逐样本迭代，生成时间以分钟计。前馈方法通过训练大规模重建模型，在单次前向推理中输出3DGS，将生成时间压缩到秒级甚至亚秒级。

\paragraph{亚秒级文本到 3D 生成}

本方法修改了SDS优化的\textbf{整个生成范式}——用\textbf{训练好的前馈重建模型}替代\textbf{逐样本迭代优化}——最终的3DGS输出表示保持不变。

\textbf{动机。}SDS 优化需要对每个文本提示单独运行数千步迭代，单样本生成时间以分钟计；前馈路线通过在大规模 3D 数据上离线训练重建模型，推理时只需一次前向传播即可输出完整 3D 资产。Turbo3D（2024.12）~\cite{2412_04470}在不到1秒内（具体为0.35秒）生成高质量3DGS资产，相比SDS优化路线的数分钟乃至数十分钟，加速了两个数量级以上。

\textbf{Baseline做法与前馈路线的动机：}SDS优化路线（如DreamFusion、ProlificDreamer）需要对每个文本提示单独运行扩散模型引导的梯度下降。优化一个3DGS场景通常需要数千到上万步SDS迭代，每步都需调用扩散模型的去噪网络计算分数梯度，单样本生成时间以分钟计。前馈路线的根本思路是：与其对每个提示\textit{在线}迭代优化，不如在\textit{离线}阶段用大规模3D数据训练一个重建模型，使其学会从多视角图像（或其潜在编码）直接推理出3DGS参数，推理时只需一次前向传播即可输出完整3D资产。

\textbf{整体架构：}Turbo3D包含两个串联的前馈组件，两者均在潜在空间中运行：（1）少步多视角扩散生成器——接收文本提示，仅用4步去噪即输出4个视角的潜在编码（而非完整像素图像）；（2）潜在空间高斯重建器（Latent GS-LRM）——直接将这些潜在编码重建为3DGS参数。整个管线无需任何迭代优化。

\textbf{从多步到少步：蒸馏中的复合模式坍缩问题：}多视角扩散模型（如MVDream）本身需要数十步去噪采样才能生成高质量多视角图像。为了将采样步数从数十步压缩到4步，Turbo3D采用分布匹配蒸馏（Distribution Matching Distillation, DMD）。其核心思想是训练一个少步学生生成器，使其输出分布逼近多步教师模型的输出分布。然而，直接对MV教师进行蒸馏会导致\textbf{复合模式坍缩}（compounding mode collapse）。MV教师本身已在Objaverse合成数据上微调，天然偏向合成、卡通化的视觉风格。蒸馏过程进一步牺牲多样性换取效率，将学生模型的输出``锁定''在Objaverse数据模式中，远离照片真实感。这是一个逐层放大的偏差：微调带来一次偏移，蒸馏带来第二次偏移，两者叠加导致生成质量严重退化。

\textbf{Dual-Teacher蒸馏：}为解决复合模式坍缩，Turbo3D引入双教师蒸馏策略——核心直觉是：让两个``专家''各司其职，MV教师负责教导``四个视角之间如何保持一致''，SV（单视角）教师负责教导``每个视角如何看起来像自然照片''。SV教师训练于大规模高质量美学图像，其数据分布远比Objaverse更加丰富和真实。损失函数为：

\begin{equation}
\mathcal{L}_{\text{DMD}}^{\text{Dual}} = D_{\text{KL}}(p_\theta \| p^{\text{MV}}) + \frac{1}{K}\sum_{i=1}^{K} D_{\text{KL}}(p_\theta(x_t^i) \| p^{\text{SV}}(x_t^i))
\end{equation}

直觉上，第一项KL散度确保四个视角联合分布与MV教师一致（维持视角间几何一致性），第二项对每个视角独立计算KL散度，将其拉向SV教师所代表的自然图像分布（恢复照片真实感）。其中$K=4$为视角数，平衡权重$\lambda=1$（实验中发现两个教师权重相当即可取得最佳效果）。消融实验表明，仅使用MV教师时蒸馏质量大幅下降（CLIP分数显著降低），加入SV教师后接近多步MV教师的生成质量。

\textbf{潜在空间GS-LRM：}传统的多视角到3D重建管线需要先将潜在编码解码为像素图像，再送入像素空间的重建模型。但VAE解码器中的Conv2D算子在高分辨率下效率低下，成为重建速度的瓶颈。Turbo3D将高斯重建器（GS-LRM，一种基于Transformer的前馈重建网络）的输入从像素空间移到潜在空间，直接使用4步生成器输出的多视角潜在编码而非解码后的图像。这一设计有双重收益：消除了VAE解码的额外耗时和内存开销，同时将Transformer的输入序列长度减半（潜在空间分辨率仅为像素空间的$1/8$）。训练时仍用像素空间的新视角渲染损失（$\mathcal{L}_2$ + 感知损失）监督重建的3DGS，确保重建质量不受影响。实验表明，潜在空间GS-LRM的重建质量与像素空间GS-LRM持平，但为整个管线节省约0.34秒。

\textbf{训练规模：}整个训练分三个阶段：（1）在Objaverse数据集（约730K对象）上微调DiT架构的多步MV扩散模型，30K迭代，32张A100 GPU，批量128，学习率$3 \times 10^{-5}$；（2）双教师蒸馏为4步生成器，10K迭代，学习率$5 \times 10^{-6}$；（3）从头训练潜在GS-LRM，80K迭代，批量256，学习率$4 \times 10^{-4}$。最终系统在256分辨率下总推理时间仅0.35秒，在512分辨率下也仅需约0.7秒，同时CLIP分数达到27.61、VQA分数达到0.76，均大幅领先同期的LGM、Instant3D等方法。

\paragraph{多视角扩散先验三阶段管线}

GradeADreamer（2024.06）~\cite{2406_09850}通过三阶段训练管线缓解Janus问题并提高生成质量。

\textbf{动机。}Janus 问题是文本到 3D 生成中的顽疾，源于 2D 扩散模型的正面视角偏见；GradeADreamer 利用 MVDream 的多视角一致性作为先验约束，从根本上抑制多面重叠。方法遵循Fantasia3D的几何-纹理解耦思路，但无需手工网格先验。

\textbf{与Turbo3D的差异：}Turbo3D通过蒸馏实现单次前馈推理，而GradeADreamer保留三阶段迭代优化，以换取更高的生成质量。

\textbf{三阶段管线：}第一阶段使用MVDream（256$\times$256）从随机高斯中优化出具有粗糙细节的高斯先验，利用MVDream的多视角一致性有效减少Janus问题。第二阶段切换至StableDiffusion 2.1（512$\times$512）精炼几何，获得更高分辨率的形状细节。第三阶段从精炼后的高斯中提取网格，优化空间变化的PBR材质模型——包含漫反射$k_d \in \mathbb{R}^3$、粗糙度与金属度$k_{rm} \in \mathbb{R}^2$、法线扰动$k_n \in \mathbb{R}^3$。全部三个阶段均使用SDS作为分数蒸馏目标（作者实验表明SDS已足够，VSD虽质量相当但耗时更长）。整个管线在单张RTX 3090上总生成时间不超过30分钟。

\paragraph{轨迹分数匹配}

Dreamer XL（2024.05）~\cite{2405_11252}针对区间分数匹配（ISM）中DDIM反演导致的伪真值不一致问题，提出轨迹分数匹配（TSM）方法。

\textbf{与Turbo3D的差异：}Turbo3D绕过SDS优化采用前馈重建，而Dreamer XL仍在SDS框架内改进蒸馏质量。

\textbf{动机：}ISM通过DDIM反演获取噪声轨迹上的$x_s$和$x_t$来计算蒸馏梯度，但反演过程在条件模型中引入累积误差，产生不一致的伪真值，影响优化质量。

\textbf{TSM核心思想：}从同一噪声潜在编码$x_s$出发，沿DDIM反演路径同时前推得到$x_\mu$和$x_t$（其中$s < \mu < t$），构建双路径比较。由于两条路径共享起点$x_s$，其累积误差$x_t - x_\mu$小于ISM中的$x_t - x_s$。TSM梯度为：

\begin{equation}
\nabla_\theta \mathcal{L}_{\text{TSM}} = \mathbb{E}_{t,c,\epsilon}\left[\omega(t)(\hat{\epsilon}_\phi(x_t, t, y_\text{tgt}) - \hat{\epsilon}_\phi(x_\mu, \mu, \varnothing))\frac{\partial g(\theta, c)}{\partial \theta}\right]
\end{equation}

理论上ISM是TSM在偏移率$\eta=0$时的特例。

\textbf{SDXL适配与逐像素梯度裁剪：}方法进一步采用Stable Diffusion XL（支持$1024\times 1024$分辨率）作为引导模型。SDXL的VAE不稳定，导致3DGS优化中梯度异常，触发过度的自适应致密化。为此提出逐像素梯度裁剪——对深度梯度$g_d$独立计算缩放因子并按像素裁剪：$g_d' = g_d \cdot \text{clip}(|g_d|, c)/|g_d|$，确保梯度方向不变的同时抑制异常区域的致密化。

\subsubsection{多视角扩散路线}

多视角扩散模型（如MVDream、Zero123++）同时生成多个视角的图像，天然提供比单视角扩散更强的3D一致性。这条路线的核心逻辑是"先生成一致的2D多视角图像，再用3DGS重建将其提升为可渲染的3D资产"——与前面单视角蒸馏路线的"边优化3D边咨询扩散模型"形成互补。

\paragraph{耦合分数蒸馏}

本方法修改了SDS优化的\textbf{梯度计算环节}，用\textbf{多视角联合分布先验}替代\textbf{单视角独立优化}，3DGS作为可学习3D表示保持不变。

\textbf{动机。}SDS 的逐视角独立采样导致扩散模型缺乏视角间关联信息，Janus 率高达 36.67\%；纯多视角微调虽可消除 Janus 但牺牲多样性（不同种子产出近乎相同结果）；CSD 通过概率乘积规则将两者耦合，同时保证一致性与多样性。CSD（Coupled Score Distillation，耦合分数蒸馏）（2025.05）~\cite{2505_04262}的核心贡献是将SDS的逐视角独立蒸馏重构为多视角联合优化问题，从根本上消除Janus多面伪影，同时保持生成的多样性。

\textbf{Baseline做法与问题：}SDS优化路线在每个优化步骤中仅采样单个视角$v_i$，用预训练的文本到图像扩散模型计算该视角的分数梯度来更新3D表示$\theta$。由于视角之间完全独立采样，扩散模型缺乏视角间的关联信息。直觉上，这就像用多只手各自拉扯同一块橡皮泥，力量之间互不协调，最终形状必然扭曲。具体表现为Janus问题：生成的3D资产在不同视角出现重复面部、多余肢体或内容漂移。实验数据显示，SDS的Janus率高达36.67\%，ProlificDreamer甚至达到56.67\%。

\textbf{设计动机——为什么需要多视角联合优化：}一条直观的解决思路是直接训练一个多视角联合扩散模型$p_t(\tilde{\mathbf{x}}_t|\tilde{v}, y)$（其中$\tilde{\mathbf{x}}_t$为多视角噪声图像集合），但这受限于有限的多视角训练数据，容易过拟合。论文从理论上揭示了三类方法在多样性-一致性权衡中的互补关系。单视角优化（SDS/VSD）利用大规模2D扩散模型获得丰富多样性，但缺乏视角约束导致Janus问题。纯多视角微调方法（如MVDream直接用SDS优化）通过有限3D数据建立视角关联，但过拟合导致不同种子初始化产出近乎相同的结果。CSD则将单视角优化的多样性与多视角先验的几何一致性进行耦合，使优化梯度同时包含两个信号。

\textbf{改了什么——从单视角到耦合多视角：}CSD的核心数学推导基于概率乘积规则（probability product rule），将多视角联合分布分解为单视角边际分布与条件联合分布的乘积：$q_t(\tilde{\mathbf{x}}_t|\tilde{v}, y) = q_t(\tilde{\mathbf{x}}_t|\tilde{v}, y, \mathbf{x}_t^i) \cdot q_t(\mathbf{x}_t^i|v_i, y)$。这一分解将原本的多视角联合优化转化为两项之和——一项负责单视角多样性，另一项负责多视角几何一致性。直觉上，这相当于让一位艺术家先自由创作（保证多样性），再由一位质检员同时检查四个正交视角的协调性（保证一致性），两者的合力塑造最终3D形态。

\textbf{实现原理——最终的CSD梯度公式：}在每次优化迭代中，CSD随机采样四个互相正交的相机视角$\tilde{v}' = (v_i, v_{i+90}, v_{i+180}, v_{i+270})$，渲染出对应图像后计算两个梯度分量：

\begin{equation}
\nabla_\theta \mathcal{L}_{\text{CSD}} \triangleq \mathbb{E}_{t, \tilde{v}', \epsilon}\left[\omega(t)\left(\underbrace{(\hat{\epsilon}_\phi(\mathbf{x}_t^i, t, y_{v_i}) - \epsilon_\varphi(\mathbf{x}_t^i, t, v_i, y))}_{\text{单视角项：保多样性}}\frac{\partial \mathbf{x}_0^i}{\partial \theta} + \lambda\underbrace{(\epsilon_M(\tilde{\mathbf{x}}_t', t, \tilde{v}', y) - \epsilon)}_{\text{多视角项：保一致性}}\frac{\partial \tilde{\mathbf{x}}_0'}{\partial \theta}\right)\right]
\end{equation}

其中单视角项由预训练SD2.1的CFG（Classifier-Free Guidance，无分类器引导）输出$\hat{\epsilon}_\phi$减去可学习LoRA（Low-Rank Adaptation，低秩适配）的噪声预测$\epsilon_\varphi$构成（受ProlificDreamer的VSD启发），负责捕捉纹理细节和保持多样性；多视角项由微调后的MVDream模型$\epsilon_M$减去添加的噪声$\epsilon$构成，负责施加跨视角几何一致性约束。$\lambda \in [0.1, 1.0]$为正则化系数，平衡多视角先验与单视角优化的相对强度。论文在3DGS优化阶段动态调整$\lambda$，在网格精调阶段固定为0.1。

\textbf{具体参数与实现细节：}CSD直接从随机初始化的3DGS开始优化，高斯点初始位置随机分布在半径0.5的球体内，颜色设为灰色，不透明度初始化为0.1。优化分两阶段进行：第一阶段对3DGS优化4000步，渲染分辨率从64逐步提升至1024，使用Adam优化器（$\beta_1=0.9, \beta_2=0.99$）；第二阶段将优化好的3DGS提取为可变形四面体网格（DMTet），分别精调几何和纹理。LoRA每$k=1$步更新一次，在优化过程中持续适配当前渲染结果。实验中，CSD在30个文本提示的基准测试上将Janus率降至3.33\%（最低），同时CLIPScore达到30.23\%（最高），生成时间约50分钟（3DGS阶段）加1.5小时（网格阶段），在两块RTX 3090上运行。

\paragraph{即插即用 3D 几何感知模块}

3D-Adapter（2024.10）~\cite{2410_18974}提出通用插件模块，为预训练的图像扩散模型注入3D几何感知能力。

\textbf{动机。}单纯的 I/O 同步（先去噪再重建）无法利用 3D 一致性反馈来纠正去噪过程中的几何错误——去噪和重建是割裂的两个阶段；3D-Adapter 在每个去噪步骤中构建"去噪-重建-渲染-反馈"闭环，使 3D 一致性信息实时回传。

\textbf{与CSD的差异：}CSD修改SDS的梯度公式使其耦合多视角先验，而3D-Adapter直接在去噪网络内部嵌入3D反馈机制，不改变蒸馏目标。

\textbf{3D反馈增强：}核心思想是在每个去噪步骤中构建``去噪$\to$重建$\to$渲染$\to$反馈''的闭环。具体而言，基础U-Net编码器处理噪声视角$\mathbf{x}_t$后，先用一份U-Net解码器输出中间去噪视角$\hat{\mathbf{x}}_t$；再通过3D重建模型（GRM，Gaussian Reconstruction Model，高斯重建模型）将其提升为一致的3D表示，渲染出RGBD视角$\mathbf{s}_t$；最后通过ControlNet编码器将渲染特征加法注入基础编码器特征，由基础解码器产出最终输出：

\begin{equation}
\mathbf{x}_0 = D_\theta(\mathbf{x}_t, \mathbf{c}, t, \{R(\hat{D}(\mathbf{x}_t, \mathbf{c}, t)^{\hat{\mathbf{x}}_t})\} \to \mathbf{s}_t)
\end{equation}

其中$R$表示3D重建与渲染，$D_\theta$表示增强后的U-Net。

\textbf{两个变体：}基于GRM的前馈版本（快速，需两阶段训练：先微调GRM适应低质量中间视角，再微调反馈ControlNet）和基于NeRF/DMTet优化的通用版本（无需训练，利用Stable Diffusion的现成tile/depth ControlNet）。推理时采用类似CFG的引导策略：$\mathbf{x}_0^g = \alpha(D_\theta(\mathbf{x}_t,\mathbf{c},t,\mathbf{s}_t) - D_\theta(\mathbf{x}_t,\mathbf{c},t,\mathbf{0})) + \beta D(\mathbf{x}_t,\mathbf{c},t) + (1-\beta)D(\mathbf{x}_t,\mathbf{0},t)$，消除ControlNet的数据集偏差。

\paragraph{3D 注意力去偏}

TD-Attn（2025.12）~\cite{2512_07345}从注意力机制层面分析多视角不一致的根因，并提出两个互补模块加以解决。

\textbf{动机。}CSD 等方法在 SDS 梯度层面耦合多视角先验，但未触及扩散模型内部产生视角偏差的根源——交叉注意力中主体词始终激活先验视角特征；TD-Attn 直接干预扩散模型的交叉注意力计算，从源头上消除视角偏差。

\textbf{与CSD的差异：}CSD在SDS梯度层面耦合多视角先验，而TD-Attn直接干预扩散模型内部的交叉注意力计算，从源头上消除视角偏差。

\textbf{先验视角偏差的理论分析：}方法发现T2I模型中存在``先验视角偏差''——训练数据中视角分布不均且标注稀疏，导致模型在处理不含视角标注的文本时隐式偏向先验视角。形式化地，视角后验分布可分解为$p(v|Y) = (1-\gamma)\delta_{v,v^*} + \gamma p_D(v|y_\text{obj})$，其中$\gamma \in (0,1)$为偏差系数。当目标视角$v^*$远离先验视角$v_\text{prior}$时，概率比$R = p_D(v_\text{prior}|y_\text{obj}) / [(1-\gamma) + \gamma p_D(v^*|y_\text{obj})]$远大于1，先验视角压倒目标条件，导致主体词在交叉注意力中始终激活先验视角特征。

\textbf{3D-AAG（3D感知注意力引导）：}在每个去噪步骤中，将2D交叉注意力图反向映射到3DGS表示上，累积多视角的注意力权重构建视角一致的3D注意力高斯：$w_i = \sum_v \sum_p I(\hat{S}_{2D}^v)[o_i(p) T_i^v(p) I(S(p)_{2D}^v)]$，其中$o_i$为高斯不透明度，$T_i^v$为透射率。3D注意力高斯融合了多视角信息，再用KL散度损失$L_\text{attn}$引导各视角2D注意力图与3D注意力高斯对齐，稀释先验视角偏差。

\textbf{HAM（分层注意力调制）：}利用语义引导树定位UNet中对视角高度敏感的注意力层，通过增强目标视角语义来降低偏好比$R$，确保后验概率$p(v|Y)$趋向目标视角$v^*$。

\paragraph{视频生成器与 3D 重建拼接}

VIST3A（2025.10）~\cite{2510_13454}提出通用框架，将预训练的文本到视频模型作为``生成器''，前馈3D重建系统作为``解码器''进行拼接。

\textbf{动机。}视频生成器与 3D 重建模型各自预训练于不同领域，直接串联存在潜在表示不对齐的问题——生成器输出的潜在编码可能无法被 3D 解码器正确解码；VIST3A 通过模型拼接和奖励微调实现跨域对齐。

\textbf{与CSD的差异：}CSD修改蒸馏梯度公式，而VIST3A不改变扩散模型架构，而是通过模型拼接构建端到端的3D VAE。

\textbf{模型拼接（Model Stitching）：}在前馈3D模型$F_{1:l}$中扫描所有候选层$k$，对每个层拟合线性拼接层$S_k^* = \arg\min_{S_k}\|BS_k - A_k\|_F^2$（其中$B$为编码器潜在编码，$A_k$为第$k$层激活），选择误差最小的层$k^*$进行拼接：$M_\text{stitched} = F_{k^*+1:l} \circ S \circ E$。拼接后仅用少量数据自监督微调拼接层$S$和下游部分$F_{k^*+1:l}$（使用LoRA防止偏离预训练权重），无需标签。

\textbf{直接奖励微调：}生成器的去噪循环产生的潜在编码需与3D VAE解码器对齐。总损失$L_\text{total} = L_\text{gen} - r(z_0(\theta, c, z_T), c)$，其中奖励函数$r$包含三项：（1）多视角图像质量（CLIP + HPSv2）；（2）3D表示质量（渲染3D场景后重新评估）；（3）3D一致性（比较视频解码器的多视角图像与3D渲染图像的$L_1$+LPIPS（Learned Perceptual Image Patch Similarity，学习的感知图像块相似度））。这种端到端对齐确保生成器产生可被3D解码器正确解码的潜在编码。

\subsubsection{大规模3D生成路线}

本路线介绍在海量3D数据上训练的大规模生成模型（即"3D基础模型"），能够一次性从文本或图像生成高质量3DGS资产。

\paragraph{大材质高斯模型}

MGM（Large Material Gaussian Model）（2025.09）~\cite{2509_22112}修改了前述3DGS生成方法的输出表示——用5通道PBR（physically based rendering，物理基础渲染）材质属性替代球谐系数编码的RGB颜色——3DGS的几何参数（位置、不透明度、缩放、旋转）保持不变。

\textbf{动机。}前述方法输出的 3DGS 用球谐系数编码"带光照的外观"，光照与材质纠缠在一起，无法更换环境光照重新渲染；游戏、影视等工业管线要求 3D 资产支持重光照，必须显式分离材质属性与光照。简言之，MGM不只是生成``看起来对''的3D模型，而是生成``物理上对''的3D资产：每个高斯基元携带反照率（albedo，物体本身的颜色，不含光照）、粗糙度和金属度，使得生成结果可以被任意环境光照重新照亮。

\textbf{Baseline做法与设计动机：}前述SDS优化路线、前馈重建路线和多视角扩散路线的生成结果有一个共同局限。它们输出的3DGS用球谐系数编码颜色，而球谐系数编码的是``带光照的外观''。这意味着渲染结果中光照和材质是纠缠在一起的：你无法将一个在日光下生成的物体放到夜景中重新照亮，也无法单独编辑材质属性。在游戏、影视、AR等实际工业管线中，3D资产必须支持``重光照''（relighting，即更换环境光照条件后重新渲染），这要求显式分离材质属性（物体本身长什么样）与光照（环境如何照亮它）。MGM正是面向这一需求设计的：在大规模3D数据上训练一个完整的``多视角PBR扩散 $\to$ 材质高斯重建''管线，一步到位地生成可重光照的3D资产。

\textbf{多视角PBR扩散模型：}核心思想是：将原本生成``着色图像''（shaded image，即带光照的RGB图）的多视角扩散模型，重定向到``材质属性空间''——让扩散模型直接输出反照率、粗糙度和金属度图。

方法微调Stable Diffusion或MVDream，训练两个子模型：一个生成反照率（3通道，对应RGB三色），另一个生成粗糙度与金属度（2通道）。对于后者，由于标准扩散模型的输入输出块只支持3或4通道，需要扩展输入输出块的外层（inflate outermost blocks），并以参数复制（parameter copy）方式初始化新增分支，确保训练从预训练权重平滑过渡。以深度图$I_d$和法线图$I_n$作为ControlNet条件输入，从$F=4$个视角生成$C=5$通道的PBR图像。直觉上，深度图和法线图告诉扩散模型``物体长什么样''（几何形状），扩散模型则在此基础上预测``物体表面是什么材质''。

训练使用Objaverse等3D数据集渲染的多视角PBR数据集$\mathcal{X}_\text{mv-PBR}$。直觉上，扩散模型的训练目标与标准多视角扩散模型一致——只是从预测RGB噪声变为预测PBR属性噪声。扩散损失定义为：

\begin{equation}
\mathcal{L}(\theta, \mathcal{X}_\text{mv-PBR}) = \mathbb{E}_{\mathbf{x}, \boldsymbol{\epsilon}, c, t}\left[\|\boldsymbol{\epsilon} - \boldsymbol{\epsilon}_\theta(\mathbf{x}_t, \boldsymbol{\tau}, I_c, c, t)\|_2^2\right]
\end{equation}

其中$\mathbf{x}_t \in \mathbb{R}^{F \times H \times W \times C}$是加噪后的PBR图像（$F=4$视角，$C=5$通道），$\boldsymbol{\epsilon}$是真实噪声，$\boldsymbol{\tau}$是文本嵌入，$I_c = \{I_d, I_n\}$是深度/法线条件，$c \in \mathbb{R}^{F \times 16}$是相机外参。白话注解：这个损失的含义是——扩散模型需要学会从噪声中逐步恢复出正确的PBR材质属性，而不是恢复带光照的RGB图。

推理时，先用CraftsMan等预训练3D网格生成模型获取代理几何（proxy geometry），渲染出多视角的深度图和法线图作为条件，再由PBR扩散模型生成各视角的材质图像。

\textbf{材质高斯表示：}一个直觉是：既然PBR属性（反照率、粗糙度、金属度）是``视角无关''的——物体表面某点的粗糙度不会因为你换个角度看就改变——那么就不需要用高阶球谐函数来编码视角依赖的外观。因此，MGM保留3DGS的所有几何参数不变，将球谐系数替换为与视角无关的5通道材质属性：3通道反照率 + 1通道粗糙度 + 1通道金属度。由于反照率不含光照分量，仅使用0阶SH系数（即一个常数值）。完整表示为$\theta = \{\mathbf{x}, \alpha, \mathbf{s}, \mathbf{q}, c_a, c_r, c_m\}$，其中$c_a \in \mathbb{R}^3$为反照率，$c_r, c_m \in \mathbb{R}^1$分别为粗糙度和金属度。

\textbf{材质高斯重建：}一个朴素方案是分别训练三个重建模型来预测反照率、粗糙度和金属度，但这会重建出三个几何不一致的点云，无法融合。为此，MGM设计统一的材质高斯表示，训练单一模型同时预测所有属性。重建阶段利用基于Transformer的体解码器（来自LaRa），接收多视角图像特征$H_f = (h_a, h_r, h_m)$和可学习查询嵌入$H_e$，预测3D高斯体积：$H_g = V_g(\text{cat}(H_f), \text{cat}(H_e))$，再通过轻量级参数解码网络生成最终高斯基元。其中，图像特征包含三类编码器的输出：反照率编码器$F_a$、深度/法线编码器$F_g$、粗糙度/金属度编码器$F_{r/m}$，各自与相机嵌入拼接后送入体解码器。

\textbf{关键发现——几何监督的必要性：}PBR图像相比着色图像缺少光照线索（例如阴影、高光），直接用PBR数据训练重建模型会导致几何质量显著下降。此外，重建出的几何必须与PBR扩散模型使用的代理几何保持一致。因此，MGM在重建阶段引入深度/法线图的双重作用：（1）作为图像编码器的输入特征，增强模型对空间位置的理解；（2）作为几何正则化监督$L_\text{Geometry} = L_2(\hat{I}_d, I_d) + L_2(\hat{I}_n, I_n)$，确保从2DGS渲染出的深度/法线与代理几何对齐。

\textbf{两阶段训练策略：}由于高质量粗糙度/金属度训练数据远少于反照率数据，直接同时训练所有通道效果不佳。MGM采用两阶段策略：第一阶段冻结粗糙度/金属度组件，仅训练反照率；第二阶段解冻所有材质组件联合训练。这种由简到难的课程学习（curriculum learning）方式提高了训练质量和收敛速度。

\textbf{可重光照渲染：}MGM生成的材质高斯天然支持重光照——这正是引入PBR表示的核心回报。具体而言，渲染时从2DGS光栅化获得各像素的反照率$c_a$、粗糙度$c_r$、金属度$c_m$和法线$\mathbf{n}$，再将其代入经典渲染方程和Cook-Torrance微面元BRDF模型计算最终颜色：

\begin{equation}
L_o(\mathbf{p}, \mathbf{v}) = \int_\Omega L_i(\mathbf{p}, \boldsymbol{\omega}_i)\, f_r(\boldsymbol{\omega}_i, \mathbf{v})\, (\mathbf{n} \cdot \boldsymbol{\omega}_i)\, d\boldsymbol{\omega}_i
\end{equation}

其中$f_r$是Cook-Torrance BRDF（Bidirectional Reflectance Distribution Function，双向反射分布函数），$\mathbf{p}$是表面点，$\mathbf{v}$是视角方向，$L_i$是入射辐射度。白话注解：给定一个环境光照贴图（environment light map），渲染方程计算每个表面点从环境中接收到的所有光照经材质反射后到达相机的颜色。更换不同的环境光照贴图即可实现任意重光照，而无需重新生成3D模型。

\paragraph{基于八叉树的 PBR 材质生成}

TexGaussian（2024.11）~\cite{2411_19654}专注于为3D网格快速生成PBR材质。

\textbf{动机。}MGM 生成完整 3DGS 资产（含 PBR），但实际工业管线中已有大量 3D 网格只需材质贴图；TexGaussian 假设输入网格已知，仅生成其表面的 PBR 材质，且通过八叉树对齐确保高斯自然贴合网格表面。

\textbf{与MGM的差异：}MGM从文本生成完整的3DGS资产（含PBR），而TexGaussian假设已有输入网格，仅生成其表面材质贴图。

\textbf{八叉树对齐的高斯泼溅：}在输入网格表面采样$N=100{,}000$个3D点构建八叉树，将3D高斯放置在每个最细叶节点中心（不额外调整位置），确保高斯自然贴合网格表面。每个高斯包含13个参数：1维不透明度、3维缩放、4维旋转、3维反照率、1维粗糙度、1维金属度（不含球谐系数，因PBR图像无光照变化）。训练时采用softplus激活并将缩放乘以0.01，确保初始高斯贴合物体轮廓而非向外扩张。

\textbf{基于八叉树的3D U-Net：}使用八叉树卷积神经网络预测高斯参数，采用回归目标而非扩散去噪——输入为八叉树叶节点的几何特征（法线、局部位移）和CLIP文本特征，输出直接为13通道高斯参数。损失函数包含多视角渲染的MSE + LPIPS（反照率/粗糙度/金属度）、alpha图MSE以及3D参数MSE：$L_\text{total} = L_\text{albedo} + L_R + L_M + L_\alpha + L_\text{3D}$。推理时将渲染的多视角图像烘焙到UV空间生成最终贴图。

\subsubsection{高斯原生生成路线}

与前述方法不同，这类方法不使用SDS优化或中间的2D表示，而是直接在3DGS表示上构建生成模型。

\paragraph{三平面高斯场表示}

本方法修改了SDS蒸馏路线的\textbf{整体生成范式}——用「三平面场表示 + 潜在扩散直接生成」替代「逐轮SDS梯度反传优化」——3DGS作为最终输出表示保持不变，但生成速度从数十分钟缩短至一分钟以内。

\textbf{动机。}SDS 的间接监督信号需要逐轮蒸馏才能驱动 3D 表示演化，耗时且不稳定；直接在 3DGS 的中间表示（三平面）上构建扩散模型可一次前向推理完成生成，且与 2D 卷积天然兼容。

\textbf{Baseline做法（SDS蒸馏路线）：}前述SDS类方法（如DreamFusion、ProlificDreamer）在每一轮优化中，先用2D扩散模型对当前渲染图像计算Score Distillation梯度，再将梯度回传至3DGS参数。这种方式需要反复迭代数百至数千轮，单次生成耗时10分钟以上，且容易受Janus问题（多面重叠）和不稳定梯度的困扰。核心瓶颈在于：2D扩散模型的监督信号是间接的，需要逐轮蒸馏才能驱动3D表示演化。

\textbf{设计动机——为什么不直接生成3DGS：}3DGS的本质是一个离散点云，每个高斯点包含位置$\mathbf{p}$、缩放$\mathbf{s}$、旋转四元数$\mathbf{q}$、不透明度$\alpha$和球谐系数SH等多通道参数（完整定义为$\mathcal{G}_{\mathbf{p}}:=\{s, \mathbf{q}, \mathrm{SH}, \alpha\}$）。这种数据结构有三个不利于生成的特性：（1）\textbf{无序性}——点之间没有固定顺序，标准卷积/注意力无法直接处理；（2）\textbf{变长性}——不同物体的高斯点数量不同，无法对齐到统一张量；（3）\textbf{多通道异质性}——位置、旋转、外观等参数的数值范围和分布差异极大。因此，需要一种中间表示将离散点云转化为连续的、固定大小的、可被卷积网络处理的格式。

\textbf{三平面表示——直觉与定义：}三平面（triplane）的核心思想是把3D空间沿三个坐标平面``拍扁''为三张2D特征图——就像从正上方、正前方、正侧方分别拍摄一张X光片，三张片叠起来就能描述整个3D物体。具体而言，triplane由$xy$、$xz$、$yz$三个平面组成，每个平面是一张$128\times 128$的特征图、16个通道。任意3D点$\mathbf{p}$通过正交投影在三个平面上查询特征（双线性插值），将三组特征拼接即为该点的完整表示：$F_{\mathbf{p}}=(F_{xy}, F_{xz}, F_{yz})$。这种表示有两重优势：（1）内存效率远高于密集体素（$3\times 128^2 \times 16$ vs.\ $128^3 \times C$）；（2）与2D卷积天然兼容，无需设计特殊的稀疏网络。

\textbf{几何与外观的通道分离：}16通道中前8通道编码几何信息（SDF + 顶点变形），后8通道编码高斯外观属性（缩放/不透明度/SH系数）。实验表明分离编码比混合编码收敛更快、渲染结果更清晰。每个三平面初始化为标准差$\sigma=0.01$的随机高斯噪声（而非零初始化），这使得解码后的SDF产生随机碎片化网格面，经几何损失快速筛除后能加速收敛。

\textbf{TriRenderer——从三平面到可渲染图像：}TriRenderer是整个框架的核心枢纽，负责将三平面场解码为标准GS点云并渲染为多视角图像。它包含两条并行分支：

（1）\textbf{几何分支}：解码三平面的前8通道，输出SDF体素场（分辨率$L\times L\times L$）及顶点变形量。通过可微Marching Cube（FlexiCubes）提取三角网格，作为高斯点的``附着表面''。

（2）\textbf{表面采样与旋转固定}：在网格面上随机采样高斯点。每个面的法线$\mathbf{n}=\mathbf{v}_1\times\mathbf{v}_2$直接定义该面所生成高斯的旋转矩阵$\mathbf{R}=[\mathbf{v}_1, \mathbf{n}\times\mathbf{v}_1, \mathbf{n}]$（转换为四元数$\mathbf{q}$），第三个缩放分量$s_3$固定为极小值使高斯``扁平''贴合表面。这意味着旋转参数不需要被学习，由几何隐式确定。

（3）\textbf{GS属性分支}：将采样点坐标查询三平面后8通道，通过三个独立的线性预测头分别解码缩放$s$、不透明度$\alpha$和SH系数。由于不同属性的数值范围差异很大，分离预测头比统一预测效果更好。

TriRenderer为所有物体共享（确保不同三平面的特征分布一致），且\textbf{完全可微分}，因此渲染损失可以直接监督几何编码和外观编码。训练仅需多视角RGBA图像和相机位姿作为监督，不需要任何原始3D数据。

\textbf{训练损失：}总损失包含渲染损失和几何正则化两部分。渲染损失为alpha图$L_1$损失、RGB图像损失（$L_1$ + SSIM的加权和）与感知损失LPIPS的加权组合：

\begin{equation}
L_\text{render} = w_1 L_\text{alpha} + w_2 L_\text{rgb} + w_3 L_\text{pips},
\end{equation}

其中$L_\text{rgb} = \sum_{i=1}^N(1-\beta)\|I_i-\hat{I}_i\|_1 + \beta\,\text{SSIM}(I_i, \hat{I}_i)$。几何正则化$L_\text{geo} = \gamma_1 L_\text{dev} + \gamma_2 L_\text{weighting} + \gamma_3 L_\text{CE} + \gamma_4 L_\text{sign}$约束SDF体素的连通性、权重尺度和符号分布，防止坍缩为空形状。实验确定的关键权重为$w_1=5.0$（alpha图权重最高，强调轮廓准确性）、$w_2=1.0$、$w_3=1.2$、$\beta=0.2$、$\gamma_4=1.0$（空形状惩罚权重最高）。

\textbf{VAE压缩与两阶段潜在扩散：}三平面表示进一步通过UNet结构的VAE压缩至低维潜在空间。VAE对几何通道和外观通道使用\textbf{解耦的编解码器}（实验证明混合编码会导致渲染结果模糊），并将三个平面的batch从$B\times 3\times H\times H\times C$重塑为$3B\times H\times H\times C$以充分利用三平面的结构同质性。VAE训练损失为$L_\text{VAE} = L_\text{Tri} + L_\text{Geo} + L_\text{Render} + \gamma L_\text{KL}$，其中$L_\text{KL}$确保潜在编码不偏离正态分布。

在潜在空间上训练\textbf{两阶段扩散模型}：第一阶段生成几何潜在编码，第二阶段以几何编码为条件生成外观潜在编码。这种逐步生成的设计利用了几何与纹理之间的强相关性——先确定形状再画纹理，符合人类直觉。推理时，扩散模型输出潜在编码$\to$ VAE解码为三平面 $\to$ TriRenderer解码为GS点云，整个流程在1分钟内完成。SH阶数设为1即可在Objaverse数据集上获得满意结果。

\textbf{与同类三平面方法的差异：}与LN3Diff、3DTopia等同样使用三平面的方法相比，DirectTriGS有两个关键优势：（1）使用3DGS而非NeRF作为目标表示，避免了NeRF的伪影、隐式编辑困难和慢速渲染；（2）使用显式网格面进行GS点绑定（而非隐式神经场），几何损失直接约束网格清晰平滑，纹理更精细、模糊更少。此外，与LN3Diff的图像到三平面编码不同，DirectTriGS采用端到端三平面VAE，避免了稀疏输入图像编码不足的问题。

\paragraph{自回归扩散统一框架}

LTM3D（2025.05）~\cite{2505_24245}整合扩散模型和自回归（AR，AutoRegressive）模型的优势，构建统一的条件3D生成框架。

\textbf{动机。}扩散模型擅长建模连续分布但缺乏全局结构推理，自回归模型擅长序列依赖建模但采样慢；LTM3D 通过前缀学习在潜在标记空间统一两种范式，并支持多种 3D 表示和多模态条件输入。

\textbf{与DirectTriGS的差异：}DirectTriGS通过三平面将GS转为2D场再用扩散模型生成，而LTM3D在潜在标记空间用自回归+扩散的混合方式建模。

\textbf{前缀学习（Prefix Learning）：}现有方法直接将图像/文本标记作为条件输入，但条件标记与形状潜在空间之间存在域间隙。LTM3D不改变编码过程，而是在生成过程中桥接这一间隙：初始化一组可学习标记$\mathbf{Q}=\{q_1, \ldots, q_k\}$作为网络参数，通过交叉注意力整合条件语义：$\mathbf{P} = \text{FF}(\text{CrossAttn}(\mathbf{Q}, \mathbf{C}) + \mathbf{Q})$。前缀标记$\mathbf{P}$在自回归生成中充当``前缀上下文''，使生成过程自然对齐条件信息。

\textbf{条件分布建模：}将联合分布$p(\mathbf{X}|\mathbf{C})$因式分解为$\prod_i p(x_i | x_{<i}, \mathbf{C})$，每个标记的条件分布独立建模。采用掩码自编码器（MAE）学习标记间依赖$x_i = f(x_{<i}, \mathbf{C})$，再通过MLP-based DenoiseNet建模单标记分布：$\mathcal{L}(x_i, \hat{x}_i) = \mathbb{E}_{\epsilon,t}[\|\epsilon - \hat{\epsilon}_\phi(x_i|t, \hat{x}_i)\|^2]$。

\textbf{潜在标记重建：}为解决自回归采样早期标记因条件信息不足而产生不确定性，设计重建模块从条件标记直接重建形状标记，与扩散采样标记线性混合：$x_{<i}^F = (1-\gamma_i)\hat{x}_{<i} + \gamma_i x_{<i}^S$。LTM3D支持多种3D表示（包括3DGS）和多模态条件（图像经DinoV2编码，文本经CLIP编码）。

\paragraph{几何引导两阶段生成}

DreamPolisher（2024.03）~\cite{2403_17237}提出基于3DGS的两阶段方法，通过几何先验初始化和视角一致性精炼提升生成质量。

\textbf{动机。}纯 SDS 优化从随机初始化出发，Janus 问题严重且纹理模糊；先用 3D 扩散模型提供几何先验初始化可缓解 Janus 问题，再用视角一致性约束精修纹理细节。

\textbf{与DirectTriGS的差异：}DirectTriGS直接在3DGS上构建生成模型，而DreamPolisher仍使用SDS/ISM优化范式，但通过点云先验和几何约束改进一致性。

\textbf{第一阶段——粗优化：}使用Point-E从文本生成初始3D点云，将各向同性高斯初始化为3DGS表示（而非随机初始化）。这一几何先验有效缓解Janus问题并提升几何一致性。随后用ISM损失优化高斯参数：

\begin{equation}
\nabla_\theta \mathcal{L}_\text{geo} = \mathbb{E}_{t,c}\left[\omega(t)(\hat{\epsilon}(x_t,t,y) - \hat{\epsilon}(x_s,s,\varnothing))\frac{\partial g(\theta,c)}{\partial \theta}\right]
\end{equation}

\textbf{第二阶段——外观精炼：}采样粗3D对象的$N=4$个视角图像，送入ControlNet-tile精炼器增强纹理细节。关键创新在于引入相机编码器——将相机外参矩阵$T_w^c = [R|t]$转换为12维向量后编码为相机感知特征$f_g$，注入ControlNet隐藏状态，使精炼过程条件化于相机方向。

\textbf{视角一致性损失：}通过场景坐标渲染器将像素映射到世界坐标，对两个视角的重叠区域施加一致性约束：$L_\text{VC} = \frac{1}{2WH}\sum_{i,j}(I_{ij} - \hat{I}_{ij})(1-\mathbf{1}[\hat{S}_{ij} \neq S_{ij}])$，确保3D空间中相同位置的像素在不同视角下具有相似颜色。

\paragraph{基于 2DGS 的表面生成}

DirectGaussian（2025.10）~\cite{2510_06967}聚焦于生成3D物体的表面表示，采用2D Gaussian Splatting（Surfels）而非标准3DGS。

\textbf{动机。}标准 3DGS 的高斯基元是体积元素，不适合表达精确表面；2D Gaussian Splatting 的 Surfel 天然贴合表面，缩放限制为二维，在表面生成任务中具有几何优势。

\textbf{与DirectTriGS的差异：}DirectTriGS用三平面中间表示再解码为GS，而DirectGaussian通过多头注意力模型直接生成Surfel参数。

\textbf{Surfel表示：}将3D物体表面参数化为有向盘面，每个Surfel由中心$\mathbf{p}$、旋转矩阵$R = [\mathbf{t}_u, \mathbf{t}_v, \mathbf{n}_p]$、缩放$\mathbf{s} = (s_1, s_2)$、不透明度$\alpha$和颜色$\mathbf{c}$组成。与3DGS不同，Surfel的缩放限制为二维（$s_1 t_u u + s_2 t_v v$），天然贴合表面。

\textbf{粗高斯生成与多视角优化：}由于数据限制无法直接端到端生成高质量Surfel，方法先生成4视角法线图和纹理图（分别由多视角法线扩散模型和法线条件纹理扩散模型生成），再通过多头注意力解码器分别预测几何属性$(\mathbf{p}, \mathbf{s}, R)$和纹理属性$(\mathbf{c}, \alpha)$。

\textbf{360度曲率约束：}在四视角优化中，原始splatting损失不足以引导完整重建。方法引入跨视角法线一致性约束——证明法向曲率一致性等价于法向量一致性：$L_\kappa = \frac{1}{|S_\text{ov}|}\sum_{P \in S_\text{ov}}\|1 - \mathbf{n}_\text{view}(P) \cdot \mathbf{n}_\text{near}(P)\|$，其中$S_\text{ov}$为相邻视角的重叠表面点集。该约束在全局360度范围内强制几何一致性，避免离散Surfel分布导致的视角间几何矛盾。

\subsubsection{场景级文本到3D生成}

本路线介绍从文本生成完整3D场景（而非单个物体）的方法，按中间表示可分为全景图路线、布局图路线和视频轨迹路线三大类。

上述方法主要面向单个物体的生成，而场景级生成面临两个核心挑战：（1）空间范围远大于物体，需要同时处理全局几何一致性和局部细节质量；（2）场景中包含多个物体，其空间关系和物理约束（如重力、遮挡）需要被显式建模。当前主流的"文本 $\rightarrow$ 中间表示 $\rightarrow$ 3DGS"策略根据中间表示的不同可分为三类：全景图路线（如 DreamScene360、HoloDreamer），通过 360 度全景图获取场景全局信息；布局图路线（如 SceneTeller、CausalStruct），通过结构化布局编码物体间的空间关系；以及视频轨迹路线（如 WorldExplorer、PanoWorld），通过自回归视频生成逐步扩展可导航空间。

\paragraph{全景高斯泼溅场景生成}

DreamScene360（2024.04）~\cite{2404_06903}是首个系统的全景 $\to$ 3DGS 场景生成管线。其核心洞察在于：全景图天然包含场景的完整 360 度信息，可避免基于 outpainting 或 SDS 的方法中常见的全局不一致问题。该管线包含三个阶段。

第一阶段为全景生成与自精炼。方法基于 StitchDiffusion（在 MultiDiffusion 框架中引入 LoRA 微调模块）生成 512 $\times$ 1024 的全景图。为确保左右边界连续性，在去噪过程中将左右边界区域拼接后额外执行一轮扩散。随后引入 GPT-4V 进行多轮提示自精炼（prompt self-refinement）：GPT-4V 评估当前全景图的视觉质量并自动修改提示词，迭代优化直至图像与文本语义对齐。这一机制有效缓解了文本到 3D 生成中的提示工程问题。

第二阶段为全景几何场初始化与点云构建。方法利用单目深度估计构建全景几何场（panoramic geometric field），将深度值逐像素投影到 3D 空间生成全局点云，作为 3DGS 的质心初始化。实验中全景图被缩放至 1024 $\times$ 2048 以生成足够密集的初始点云。值得注意的是，DreamScene360 禁用了 3DGS 的致密化过程以避免浮子（floater）伪影。

第三阶段为语义与几何正则化。为处理单视角输入中不可见区域的问题，方法从全景图中提取虚拟视角并施加无监督的语义对齐损失和几何一致性损失，引导高斯优化填补信息缺失区域。

在实验方面，DreamScene360 与 LucidDreamer 进行对比，在 CLIP 嵌入距离、NIQE、BRISQUE 和 Q-Align 等无参考质量指标上均取得显著优势。方法支持室内外多种风格的无约束 360 度场景生成，无需视频或多视角数据。

\paragraph{全景初始化全封闭场景}

HoloDreamer（2024.07）~\cite{2407_15187}针对基于 outpainting 的方法难以保证全局一致性的问题，提出从全景图初始化全封闭场景的策略。方法首先在全景数据库上微调扩散模型，生成符合等距柱状投影几何特征的基础全景图。为增强风格多样性，方法从中提取线稿（lineart）作为几何条件，通过线稿控制的扩散模型生成多种风格化全景，再经超分辨率网络提升分辨率。

全景重建阶段采用增强的两阶段方案：先基于全景深度估计构建初始 3DGS 并进行基础优化，再通过环形融合（circular blending）技术确保全景左右边界无缝衔接，并修复缺失区域以增强场景完整性。与 DreamScene360 相比，HoloDreamer 在风格多样性（通过线稿控制实现）和场景封闭性（通过两阶段修复实现）上有所改进，但未使用 GPT-4V 进行提示自精炼。

\paragraph{15 分钟快速场景生成}

FastScene（2024.05）~\cite{2405_05768}的核心贡献在于将全景修复问题从全景图空间转移到立方体贴图空间，从而规避等距柱状投影的严重畸变。其流程分为三步：首先用 Diffusion360 生成全景图并通过 EGformer 估计深度图；然后提出粗视角合成（Coarse View Synthesis, CVS）——根据球面全景的等距投影公式将像素投影到 3D 球面坐标，再根据新相机位姿重新投影获得带孔洞的退化视角；最后提出渐进式新视角修复（Progressive Novel View Inpainting, PNVI），将全景分解为 6 个立方体贴图面分别用重训练的 AOT-GAN 修复，再重投影回全景空间。多次迭代 PNVI 获得多个干净新视角后，通过多视角投影生成透视图像并用 3DGS 重建。

整个流程在单张 RTX A6000 GPU 上仅需约 15 分钟（全景生成 10 秒 + PNVI 约 2 分钟 + 训练数据准备约 3 分钟 + 3DGS 优化约 10 分钟），是当时最快的全景到 3DGS 场景生成方法。与 DreamScene360 的 GPT-4V 自精炼不同，FastScene 通过立方体贴图分解来处理修复畸变。

\paragraph{文本驱动全景增强}

SceneDreamer360（2024.08）~\cite{2408_13711}聚焦于长文本到高质量全景图的生成，通过三阶段全景增强管线提升分辨率和细节。第一阶段利用 GPT-4 将复杂的场景描述扩展为多样化的训练文本，并在 Habitat Matterport 数据集上对 PanFusion 模型进行 MLP+LoRA 微调以适配长文本输入。第二阶段通过三步超分辨率（超分、细节增强、分辨率提升）将初始全景图逐步提升到高分辨率。第三阶段将高分辨率全景通过等距到透视（e2p）算法分解为多视角图像，利用点云融合初始化方法去除重复点并构建空间一致的初始点云，最终通过 3DGS 训练生成完整的 3D 空间场景。

\paragraph{语言到 3D 房间设计}

SceneTeller（2024.07）~\cite{2407_20727}采用模块化策略，将场景生成分解为布局生成、场景组装和风格化三个独立阶段。在布局生成阶段，方法利用 LLM 从文本描述解析出每个物体的 3D 边界框（包含类别名、中心坐标、尺寸和朝向），并通过 CSS 格式的属性-值对描述将布局信息结构化，以便 LLM 通过上下文学习进行布局推理。在场景组装阶段，通过上下文学习从数据库中检索匹配的 CAD 模型并放入对应的边界框中。在风格化阶段，将组装好的场景初始化为 3DGS 并通过 SDS 优化使其外观与文本提示对齐。与全景图路线不同，SceneTeller 的中间表示是结构化的 3D 布局而非图像，因此天然支持通过额外文本提示修改整个场景或个别物体的外观。

\paragraph{组合式场景生成}

组合式场景生成需要同时处理多个物体的空间关系和物理约束。CompGS（2024.10）~\cite{2410_20723}的核心洞察是：3DGS 的显式表示天然支持局部化渲染，可以单独渲染物体 A 和物体 B 所在的子空间，从而在"分而治之"的意义上实现全局组合生成。方法首先从复杂文本提示生成一张"良好组合"的 2D 图像作为参考，将 2D 组合性先验迁移到 3DGS 初始化中。然后设计动态 SDS 优化策略，自动分解高斯基元为不同实体部分，分别进行实体级优化（关注单个物体的几何细节）和组合级优化（关注物体间的交互与遮挡关系）。

LayoutDreamer（2025.02）~\cite{2502_01949}则利用有向场景图（directed scene graph）从文本中提取物体间的空间依赖关系（如"在...上面"、"靠近..."），通过尺度感知密度调整（scale-aware density adjustment）确保缩放后的 3DGS 保持一致的密度，并通过链式定位方法（chain-based positioning）沿场景图的依赖链依次确定物体位置。最终的物理和布局能量函数确保场景的物理合理性（如物体不会悬浮在空中）。

\paragraph{因果推理驱动场景生成}

CausalStruct（2025.09）~\cite{2509_15249}将因果推理（causal reasoning）嵌入 3D 场景生成，解决多物体放置中的空间歧义问题。与 LayoutDreamer 的场景图不同，CausalStruct 构建的因果图不仅编码物体间的空间关系，还显式建模因果依赖和物理约束（如"桌子因果依赖于地板的支撑"）。方法包含三个阶段。（1）利用 LLM（大语言模型）构建因果图，通过因果顺序（causal order）确定物体的放置顺序，被依赖的物体先放置。（2）通过因果干预（causal intervention）根据物理驱动约束调整空间配置。（3）引入 PID（比例-积分-微分）控制器迭代调整物体尺度和位置，其中 $\tanh$ 函数作为执行器（actuator）平滑控制信号以防止过度修正。3DGS 和 SDS 在最终阶段用于提升形状精度和渲染稳定性。

\paragraph{可导航场景生成}

WorldExplorer（2025.06）~\cite{2506_01799}通过自回归视频轨迹生成构建完全可导航的 3D 场景。方法先创建全景初始化——用 Flux 模型生成四张描述场景不同区域的图像，结合固定相机位姿（均匀分布的偏航角 0/90/180/270 度）和单目深度估计构建初始全景点云。然后利用相机引导的视频扩散模型沿预定义轨迹迭代扩展场景：每次生成一段短视频探索场景的特定部分（"进入并环绕"不同区域），通过场景记忆机制（scene memory）从已生成数据中选择最优条件图像以确保 3D 一致性，碰撞检测机制防止相机穿过几何体。最后将所有视角融合为统一的 3DGS 表示。

PanoWorld（2026.05）~\cite{2605_17916}面向整栋房屋 VR 导览的生成，将全屋合成建模为 360 度全景节点的自回归生成过程。方法的核心设计是楼层平面衍生的 3D 外壳作为全局几何代理（提供墙壁、地板、天花板等低频结构约束），配合动态 3DGS 缓存作为可渲染的空间记忆。每个新节点的生成都依赖此前所有节点的信息：从当前 3DGS 缓存渲染视觉记忆，结合几何代理和相邻全景图合成新的高分辨率全景，再通过前馈全景 LRM（Large Reconstruction Model）将生成的全景提升为局部 3DGS 更新并融入全局缓存。这种"房间感知"的自回归策略确保同一房间内的视角互相增强几何一致性，而被墙壁隔开的视角则不交换外观信息。

\paragraph{单全景图到沉浸场景}

Scene4U（2025.04）~\cite{2504_00387}提出分层 3D 场景重建框架，核心洞察在于：全景图中的前景物体遮挡了大量背景信息，直接从全景图构建 3D 场景会导致大面积信息缺失。方法首先利用开放词汇分割模型（Semantic Segment Anything）和大语言模型（QWen-Plus）将全景图分解为四个语义层级——天空、背景、前景和动态物体层，然后通过 FLUX.1 扩散模型分层修复被遮挡区域（如移除行人后修复被遮挡的建筑物），最后为每层独立初始化 3DGS 并分层优化。在 Scene4U 构建的 WorldVista3D 数据集上，方法在 PSNR、SSIM、LPIPS 等重建指标上均优于 DreamScene360。

ExScene（2025.03）~\cite{2503_23881}设计两阶段管线。第一阶段通过多模态扩散模型生成高保真全景图——使用 LLaVA 从单视角输入图像提取语义描述，通过多模态编码器融合文本特征和视觉特征作为扩散模型的条件输入。第二阶段开发全景深度估计方法：将全景投影到 20 个平面（二十面体均匀覆盖球面），对每个平面独立估计深度，然后通过仿射变换场优化消除不同平面间的尺度-偏移不一致。最后结合几何信息训练初始 3DGS 模型，通过 2D 稳定视频扩散先验进行 GS 精炼。

\paragraph{精确室内场景生成}

GuidedSceneGen（2026.03）~\cite{2603_13910}维持绝对世界坐标系生成度量精确、全局一致的室内场景。方法首先利用 Holodeck 从文本生成 3D 代理环境，提取度量深度和语义分割图（NYUv2-40 标签）并排列为立方体贴图（cubemap）表示——6 张透视图像作为立方体的面，形成 360$^\circ$ 全景图。这一表示相比标准球面全景图减少了像素畸变。然后通过立方体贴图多视角潜在扩散模型（基于 Stable Diffusion 2 并扩展膨胀注意力层），在语义适配器和深度适配器的双重约束下生成全景图。最后从代理深度计算覆盖感知（coverage-aware）的相机轨迹来探索未观测区域，通过相机引导的视频扩散模型合成新视角并融合为 3DGS，所有视角已对齐到绝对世界坐标系。

Rein3D（2026.04）~\cite{2604_10578}采用"恢复-精炼"范式。方法首先从文本条件化全景图及其估计深度构建粗糙 3DGS，然后通过径向探索策略（radial search strategy）从场景中心向外沿轨迹渲染退化全景视频——这些视频揭示了初始 3DGS 中的遮挡区域和几何缺陷。退化视频通过全景视频到视频扩散模型恢复高保真几何和纹理，再经视频超分辨率增强。精炼后的全景视频作为伪真值监督全局 3DGS 的迭代更新。与 GuidedSceneGen 的"布局$\rightarrow$全景$\rightarrow$探索"路线不同，Rein3D 的路线是"全景$\rightarrow$粗GS$\rightarrow$渲染退化视频$\rightarrow$恢复$\rightarrow$精炼GS"。

\paragraph{PSGS与SceneDreamer360的进阶}

PSGS（2026.02）~\cite{2602_00463}提出双层优化架构生成语义一致的全景图。布局推理层利用 T2I-R1 模型的语义级思维链（Chain-of-Thought）框架，通过多种温度设置生成多样化的语义解读并融合，将输入文本解析为结构化的空间布局、物体关系和功能约束。自优化层在 DreamScene360 基础上引入 Qwen MLLM 的迭代反馈机制——扩散模型生成初始全景后，Qwen 分析其质量并生成优化提示词，多轮迭代精炼视觉细节。全景滑动机制（panorama sliding）通过策略性采样重叠的 90 度视角初始化全局一致的 3DGS 点云，克服了传统方法中特征截断和感知盲点的问题。

\paragraph{多模态世界模型}

HY-World 2.0（2026.04）~\cite{2604_14268}是首个统一"生成"与"重建"的开源多模态世界模型框架，支持文本、单图、多视角图像和视频四种输入模态。文本/单图输入触发世界生成模式，通过四阶段管线生成可导航的 3DGS 场景：（1）HY-Pano 2.0 在大规模真实+合成混合全景数据上训练，通过无需显式几何变换的隐式学习实现从透视到全景的空间映射；（2）WorldNav 进行语义感知的轨迹规划，结合 42 视角密集深度采样和 NavMesh 碰撞检测；（3）WorldStereo 2.0 在关键帧潜在空间中运行相机引导视频扩散，通过全局几何记忆（GGM）和空间立体记忆（SSM）确保多轨迹一致性；（4）WorldMirror 2.0 通过前馈 3D 重建将关键帧融合为最终 3DGS。多视角/视频输入触发世界重建模式，WorldMirror 2.0 直接从密集观测恢复精确 3D 结构。

\paragraph{视频扩散直接场景生成}

VideoRFSplat（2025.03）~\cite{2503_15855}提出直接利用视频生成模型生成无约束真实世界场景的 3DGS，无需任何 SDS 优化。方法的核心设计是双流架构——在预训练的 Mochi 视频生成模型（约 100 亿参数）旁附加专用的位姿生成模型（更轻量的 16 块 Transformer），通过通信块（communication block）周期性地交换信息。关键洞察是：将位姿嵌入直接拼接到视频模型中会因分布差异降低文本泛化能力和生成质量，因此位姿模型使用独立的参数空间。异步采样（asynchronous sampling）策略使位姿模态以加速因子 $\alpha=0.2$ 更快地去噪，用已收敛的位姿作为条件化信号来稳定多视角图像生成。高斯泼溅解码器从生成的潜在表示中直接重建 3DGS。在 T3Bench、MVImgNet 和 DL3DV 上的实验表明，VideoRFSplat 在不使用 SDS 精炼的情况下即超越 SplatFlow 和 Director3D 等方法。

\paragraph{统一生成与导航框架}

SplatFlow（2024.11）~\cite{2411_16443}提出统一框架同时支持 3DGS 生成和编辑。方法包含两个核心组件：多视角整流流（Rectified Flow）模型和高斯泼溅解码器（GSDecoder）。多视角 RF 模型基于 Stable Diffusion 3 的冻结编码器，在潜在空间中同时生成 $K$ 个视角的图像、深度图和相机位姿（Pl\"ucker 射线表示），通过膨胀注意力层实现跨视角信息交换。GSDecoder 建立在 SD3 解码器架构上并进行关键修改——调整输入/输出通道以适应图像+深度+射线拼接的输入和像素对齐的 3DGS 参数输出，修改注意力层为跨视角注意力。由于 RF 模型与 SD3 共享潜在空间，SplatFlow 可利用无需训练的 Repaint 和 SDEdit 技术实现物体替换、新视角合成和相机位姿估计。

\paragraph{SIC3D: 风格化3DGS生成}

除了通用场景生成，3DGS+文本还渗透到特定领域。SIC3D（2026.04）~\cite{2604_08760}引入风格图像条件控制，通过变分风格化分数蒸馏（VSSD）损失从参考图像迁移风格到3DGS。

\paragraph{3DBonsai与PlantDreamer: 自然物体重建}

在自然物体重建方面，3DBonsai（2025.04）~\cite{2504_01619}设计了可训练的3D空间殖民算法生成盆景结构，作为3DGS的几何先验；PlantDreamer（2025.05）~\cite{2505_15528}利用深度ControlNet和可调高斯剔除算法生成高保真合成植物模型。

\paragraph{SplatFont3D: 3D艺术字体生成}

在字体方面，SplatFont3D（2025.12）~\cite{2512_00413}专注于3D艺术字体生成，通过Glyph2Cloud模块和动态组件分配策略实现部件级风格控制。

\paragraph{ClotheDreamer与GarmentDreamer: 可穿戴3D服装生成}

在服装方面，ClotheDreamer（2024.06）~\cite{2406_16815}和GarmentDreamer（2024.05）~\cite{2405_12420}分别面向可穿戴3D服装生成，前者提出解耦服装高斯（DCGS）分别优化服装和身体，后者利用3DGS引导和神经纹理场生成多样化的服装几何与纹理。

\paragraph{ART3D: 艺术场景生成}

在艺术创作方面，ART3D（2024.05）~\cite{2405_10508}将扩散模型与3DGS结合用于艺术场景生成，通过图像语义迁移算法和深度一致性模块桥接艺术与真实图像的域差异。

\paragraph{Interactive3D: 场景级可控生成}

Interactive3D（2024.04）~\cite{2404_16510}是场景级可控生成的代表性方法，赋予用户在生成过程中精确控制的能力。方法的核心设计是在两个级联阶段使用互补的 3D 表示：第一阶段使用 3DGS，利用高斯基元的独立性和显式性实现灵活的交互操作——添加/移除组件（通过级联高斯集合或按部件索引删除）、可变形/刚性拖拽（如拖动恐龙头部朝右看，类似 DragGAN 的 3D 版本）、几何变换和语义编辑（通过文本替换改变物体风格）。第二阶段将高斯通过快速 NeRF 蒸馏转换为 InstantNGP，通过交互式哈希精炼（Interactive Hash Refinement）模块添加细节并提取几何——该模块在哈希表中添加额外的可训练特征层，仅精炼用户感兴趣的区域。整个流程在 A100 GPU 上 20k 步（第一阶段 10k + 第二阶段 10k）完成，在 CLIP R-Precision 上超越 ProlificDreamer。

\paragraph{ESCT3D: 提示自优化可控生成}

Interactive3D 提供了精细的几何控制能力，但从不同角度提升可控性的方法还有多种。ESCT3D（2025.04）~\cite{2504_10316}从提示工程角度切入，利用 GPT-4V 进行提示自优化，并支持多种条件输入（风格、边缘、涂鸦、姿态等组合），通过多视角信息（深度、掩码、特征和图像）解决 Janus 问题。

\paragraph{TV-3DG: 分类器分数匹配可控生成}

TV-3DG（2024.10）~\cite{2410_21299}从蒸馏策略角度改进可控性，提出分类器分数匹配（CSM），去除 SDS 中的差异项并使用确定性噪声添加过程，结合视觉提示和注意力融合实现定制化 3D 生成。

\paragraph{MultiDreamer3D: 多概念定制生成}

MultiDreamer3D（2025.01）~\cite{2501_13449}面向多概念定制生成的组合问题：先用 LLM 布局控制器生成 3D 边界框，再创建粗点云放入边界框并初始化带概念标签的 3DGS，最后通过概念感知区间分数匹配精炼。

\paragraph{VRsketch2Gaussian: VR草图交互生成}

VRsketch2Gaussian（2025.03）~\cite{2503_12383}则引入全新交互模态——VR 草图，通过 Sketch-CLIP 特征对齐和解耦的条件控制（VR 草图控制几何、文本控制外观）实现高效高保真生成。

\paragraph{垂直领域场景生成}

DreamSat（2024.10）~\cite{2410_05097}将Zero123 XL微调到高质量航天器数据集上，集成到DreamGaussian框架中实现航天器3D重建。DreamSea（2025.03）~\cite{2503_06784}利用水下机器人采集的真实海底图像训练扩散模型，生成包含 RGBD 四通道的海底图像（同时预测颜色和几何），再通过深度融合构建 3DGS 模型，实现水下地形的光照真实新视角渲染——这是 3DGS 生成技术在海洋科学领域的首次应用探索。

\paragraph{城市场景生成}

MeSS（2025.08）~\cite{2508_15169}利用城市网格模型作为几何先验，是城市场景生成路线的代表性方法。与全景图路线不同，MeSS 的输入是具有语义和实例标签但无纹理的 3D 城市网格模型。方法采用稀疏到密的视角生成策略：第一阶段通过级联 ControlNet（ControlNet-S 用于几何条件，ControlNet-n 用于外推生成）沿预定义相机路径自回归生成关键视角图像，深度以色图编码保留远近场深度差异；第二阶段从 3DGS 中间渲染缺失视角，通过外观引导修复（Appearance Guided Inpainting, AGInpaint）填补伪影；第三阶段通过全局一致性对齐消除不同视角学习的高斯泼溅面（Gaussian surfels）之间的视觉不一致。在 Unreal Engine 5 的 City Sample 数据集上，MeSS 在 FID（Frechet Inception Distance，Frechet起始距离）和 KID（Kernel Inception Distance）指标上优于 WonderWorld 和 VistaDream。

ScenDi（2026.01）~\cite{2601_15221}整合 3D 和 2D 扩散模型。3D 生成器是一个 3D 潜在扩散模型，包含体素到 3DGS 的 VQ-VAE——直接从 2D 监督和伪 3D 真值（来自度量单目深度估计）学习，避免逐场景优化的高昂代价。VQ-VAE 将彩色体素网格编码为潜在码再解码为 3D 高斯基元。2D 视频扩散模型以 3D 生成器渲染的低分辨率图像为条件增强外观细节并合成体积外的背景内容。方法支持通过 3D 边界框、路网图或文本提示进行可控合成。

\paragraph{前向场景与组合优化}

RealmDreamer（2024.04）~\cite{2404_07199}专注于前向场景（forward-facing scenes）的文本到 3D 生成。方法的核心洞察是利用 2D inpainting 扩散模型为未知区域提供低方差监督，同时配合深度扩散模型的几何蒸馏赋予高保真几何。初始化策略是关键：方法首先生成参考图像并通过单目深度模型提升为点云，然后通过 outpainting 向左右两侧扩展生成辅助视角的点云，并预计算遮挡体积（occlusion volume）标记未确定的 3D 区域。在优化阶段，以初始场景估计为条件的 inpainting 模型为被遮挡区域提供约束，深度蒸馏损失提升几何质量，最终通过锐度过滤器微调获得高质量 3DGS。RealmDreamer 无需视频或多视角数据即可合成多种风格的高质量 3D 场景。

DreamScape（2024.04）~\cite{2404_09227}提出 3D 高斯引导（3D Gaussian Guide, $3DG^2$）机制，利用 LLM 从文本编码语义原语、空间变换和关系，实现局部到全局的优化。$3DG^2$ 的参数包括每个物体的类别、初始化细节、空间变换（位置、尺度、旋转）和文本描述，这些参数由 LLM 生成并在优化中精炼。局部优化阶段采用渐进式尺度控制（progressive scale control）针对局部物体生成中的训练不稳定问题；全局优化阶段对整个 $3DG^2$ 进行全局一致性训练。碰撞关系建模通过碰撞训练损失确保物理正确性。方法还设计了针对雨雪等弥散物体的稀疏初始化和致密化策略。

\paragraph{基于词汇丰富度的GS初始化: 智能初始化策略}

上述方法聚焦于生成管线的主流程，还有一些工作从特定环节切入改进效率和质量。在初始化策略方面，基于词汇丰富度的GS初始化框架（2024.08）~\cite{2408_01269}指出简单球体初始化难以处理复杂文本描述的形状，提出将3D高斯聚合到空间均匀体素中，设计全局信息感知（GIP）块和高斯-文本融合（GTF）块，使每个高斯能吸收来自其他区域的空间信息和来自文本的语义信息。

\paragraph{PanoDreamer: 全景初始化场景生成}

PanoDreamer（2025.04）~\cite{2504_05152}利用大语言模型和warp-refine管线生成360度全景，然后提升为3D点云，再用多种方法生成与初始点云一致的额外图像来扩展和精炼，最终通过3DGS创建完整3D场景。

\paragraph{BloomScene: 轻量级结构化场景生成}

BloomScene（2025.01）~\cite{2501_10462}提出轻量级结构化3DGS框架用于跨模态场景生成，包含渐进式点云重建和分层深度先验正则化机制，以及结构化哈希网格压缩冗余。

\paragraph{Hash3D: 无训练通用加速}

在加速与优化方面，Hash3D（2024.04）~\cite{2404_06091}提出无需训练的通用3D生成加速方法——观察到相邻时间步和相机角度渲染的特征图存在冗余，通过自适应网格哈希复用特征图，将文本到3D优化加速1.3--4倍。

\paragraph{Dream3DVG: 双分支渐进细节控制}

Dream3DVG（2025.05）~\cite{2505_21377}提出双分支优化框架——辅助3DGS分支桥接文本提示与矢量图形的域差异，通过调度分类器自由引导实现从粗到细的渐进细节控制。

\paragraph{GaussVideoDreamer: 渐进式视频修复场景生成}

GaussVideoDreamer（2025.04）~\cite{2504_10001}通过渐进式视频修复策略和不一致感知高斯泼溅（Inconsistency-Aware GS）实现单图3D场景生成。

\paragraph{EarthCrafter: 地理尺度3DGS生成}

\begin{takeaway}
文本到3D生成将3DGS与扩散模型结合，核心范式是 SDS/VSD 损失引导高斯优化。技术演进从单物体生成到场景级生成、从静态到动态、从文本驱动到多模态驱动。关键挑战包括多视角一致性（Janus问题）、生成质量和优化效率的平衡。GaussVideoDreamer 的不一致感知策略和 EarthCrafter 的地理尺度生成代表了当前的前沿。
\end{takeaway}

在规模化扩展方面，EarthCrafter（2025.07）~\cite{2507_16535}将3DGS生成扩展到地理尺度——在50000个600m$\times$600m航拍场景上训练双稀疏潜在扩散模型，支持从语义引导的城市布局生成到无条件地形合成的多种应用。

%=====================================================================


\section{图像到3D生成}
\label{sec:08_2}

单张图像到三维重建（single-image 3D reconstruction）是计算机视觉中最经典且最具挑战性的问题之一。其核心困难在于深度歧义（depth ambiguity）：单张二维图像丢失了完整的几何信息，同一张图像可能对应无数种不同的三维结构。近年来，基于 3DGS 的图像到3D生成方法取得了突破性进展。基本思路是利用预训练的2D扩散模型（diffusion model）作为先验知识，从单张输入图像推理出不可见区域的外观与几何，再通过 3DGS 表示构建完整的三维模型。

本节按照技术路线将这些方法组织为四个类别：（1）前馈重建路线，即通过单次网络推理直接输出3D模型的方法；（2）优化式重建路线，即通过多视角扩散模型生成多视角图像并迭代优化的方法；（3）大规模重建基础模型路线，即从大规模3D数据训练通用重建模型的方法；（4）高质量/高保真重建路线，即通过细化机制或特殊约束提升生成质量的方法。此外还有一些从特殊角度出发的方法，在末尾单独讨论。

以上文本到3D生成方法主要关注从零开始创建3D内容。然而在实际应用中，用户往往手头已有一张照片，希望直接从中恢复3D结构。下一节转向这一更具应用价值的问题——图像到3D生成。

%---------------------------------------------------------------------
\subsection{单图到3D重建}
\label{sec:08_2_1}

本节讨论5条技术路线：（1）前馈重建路线；（2）优化式重建路线；（3）大规模重建基础模型路线；（4）高质量与高保真重建路线；（5）其他单图到3D方法。

%---------------------------------------------------------------------
\subsubsection{前馈重建路线}
\label{sec:08_2_1_1}

本路线介绍通过单次网络推理直接从输入图像输出完整3DGS表示的方法，核心优势在于推理速度快（秒级甚至亚秒级），无需逐样本迭代优化。

前馈重建方法的核心思想是：训练一个编码器-解码器网络，使其能够从单张图像一次性推理出完整的3DGS表示。这类方法的优势在于推理速度快（通常在秒级甚至亚秒级），且不需要针对每个输入进行耗时的迭代优化。

\paragraph{锚点方法：前馈重建}

\textit{从优化到前馈的范式转换。}\;上一节中介绍的 SDS 优化路线同样可以迁移到图像到3D任务：给定单张输入图像，以预训练的多视角扩散模型（如 Zero123）为引导，通过数万步 SDS 梯度迭代优化 3DGS 参数。这条路线虽然生成质量较高，但优化一个物体需要数分钟，且 SDS 固有的过平滑和几何失真问题依然存在。前馈重建路线正是为了解决这一效率瓶颈而生：训练一个编码器-解码器网络，使其在单次前向推理中直接从输入图像输出完整的 3DGS 表示，推理速度从分钟级压缩到秒级甚至亚秒级。

\textit{像素-高斯对应关系的瓶颈。}\;前馈路线的代表性先驱是 Splatter Image，它将每个输入像素通过 U-Net 映射为一个 3D 高斯基元——像素数即为高斯数，像素的空间位置即决定高斯的初始位置。这种像素-高斯固定对应（pixel-Gaussian correspondence）的设计简洁高效，但存在根本局限：高斯数量被锁定为像素数（如 $64 \times 64 = 4096$），无法根据物体复杂度自适应增减；更严重的是，大量高斯被迫分配到图像背景区域，而非集中在物体表面的几何与纹理关键位置上。换言之，像素对应关系是一种"空间均匀铺开"的策略，它浪费了大量表示能力去拟合对应关系本身，而非表达目标物体的真实结构。

\textit{改了什么：从逐像素映射到可学习查询。}\;LeanGaussian（2024.04）~\cite{2404_16323} 的核心贡献是彻底打破像素-高斯固定对应关系，用 Deformable Transformer 中的 $N$ 个可学习查询（learnable query, $\mathbf{q} \in \mathbb{R}^{N \times C}$）替代逐像素映射。每个查询代表一个独立的 3D 高斯椭球，其数量 $N$ 和空间位置完全由网络在训练中自由学习决定，不再与输入图像的像素布局有任何绑定。直觉上，如果说 Splatter Image 是"把像素格子直接倒进 3D 空间"，那么 LeanGaussian 就是"让网络自己决定在哪里放多少个高斯椭球"。

\textit{实现原理。}\;LeanGaussian 的完整管线包含三个阶段。

第一阶段是\textbf{深度感知特征提取}。输入图像分别通过 U-Net 和 DINOv2 提取像素对齐特征和深度感知特征。U-Net 利用逐像素对应关系保证输入视角的像素位置一致性；DINOv2 在深度预测任务上预训练，提供物体结构和显著内容的位置先验。两类特征通过特征金字塔网络（Feature Pyramid Network, FPN）风格的融合块聚合：DINOv2 特征上采样后与 U-Net 特征相加，经卷积层处理后下采样（步长和卷积核大小均为 14），得到融合特征 $\mathbf{F}_{\mathrm{img}}$ 作为后续 Transformer 解码器的 key 和 value。

第二阶段是\textbf{3D 高斯初始化}。$N$ 个查询 $\mathbf{q}_{\mathrm{init}} \in \mathbb{R}^{N \times C}$ 随机初始化，通过一个 splatting 头 $\mathcal{S}$ 预测初始高斯参数 $G_{\mathrm{init}} = \mathcal{S}(\mathbf{q}_{\mathrm{init}})$。splatting 头使用 Xavier 均匀初始化权重，偏置通过预设超参数控制高斯参数的初始范围，避免随机初始化导致的收敛困难。

第三阶段是\textbf{可变形 Transformer 解码器的逐层精化}。解码器由 $L$ 层组成，每层包含可变形交叉注意力（GaussianDFA）、自注意力和前馈网络三个子模块。核心创新在于\textbf{3D 参考点投影机制}：每个 3D 高斯的中心被定义为 3D 参考点，通过针孔相机模型投影到图像特征平面上：

\begin{equation}
\label{eq:leangaussian_proj}
P^l = \frac{C_{\mathrm{intrin}} \cdot \mu^l}{d^l}
\end{equation}

其中 $C_{\mathrm{intrin}}$ 是相机内参矩阵，$\mu^l$ 和 $d^l$ 分别是第 $l$ 层高斯的中心坐标和深度值。投影后的 2D 参考点 $P^l$ 作为可变形注意力的锚点位置。网络进一步学习采样偏移量 $\Delta S^l = \mathrm{MLP}^S(\mathbf{q}^l)$，得到采样点 $S^l = P^l + \Delta S^l$，在图像特征上通过双线性插值提取特征值。注意力权重由查询经 $\mathrm{MLP}^A$ 计算：

\begin{equation}
\label{eq:leangaussian_attn}
\mathbf{q}_n^{l\prime} = \sum_{p=1}^{N_P} A_p \cdot \mathrm{Bilinear}(\mathbf{F}_{\mathrm{img}}, P^l + \Delta S_p^l), \quad A_p = \mathrm{MLP}^A(\mathbf{q}_n^l)
\end{equation}

其中 $N_P$ 是采样点数量。更新后的查询依次经过自注意力（在高斯椭球之间交换信息）和前馈网络：

\begin{equation}
\label{eq:leangaussian_decoder}
\mathbf{q}^{l+1} = \mathrm{FFN}\!\left(\mathrm{SelfAttn}\!\left(\mathrm{GaussianDFA}(\mathbf{q}^l, P^l, \mathbf{F}_{\mathrm{img}})\right)\right)
\end{equation}

\textit{白话解释}：公式~(\ref{eq:leangaussian_attn})描述了第 $n$ 个高斯的查询如何从图像特征中"采样"信息——注意力权重 $A_p$ 决定了每个采样位置的重要性，而采样位置由 3D 参考点投影到图像后加上学习的偏移量决定。公式~(\ref{eq:leangaussian_decoder})则说明每层解码器先用可变形注意力从图像提取信息，再用自注意力让高斯之间交换信息，最后通过前馈网络输出更新后的查询。

\textit{多层精化与参数增量更新。}\;LeanGaussian 采用逐层增量更新的策略：每层 splatting 头 $\mathcal{S}^l$ 预测参数增量 $\Delta G^l = \mathcal{S}^l(\mathbf{q}^l)$，然后通过 $G^{l+1} = G^l \oplus \Delta G^l$ 更新高斯参数，其中 $\oplus$ 表示对旋转参数用乘法更新、对其余参数（位置、尺度、不透明度、球谐系数）用加法更新。每层渲染并计算监督损失，实现由粗到细的迭代精化。高斯中心坐标不直接回归 3D 坐标，而是参数化为深度 $d$ 和三个偏移量 $(\delta_x, \delta_y, \delta_z)$：$\mu = (u_1 d + \delta_x,\; u_2 d + \delta_y,\; d + \delta_z)$，其中 $\mathbf{u}_1, \mathbf{u}_2$ 为相机射线的方向向量。这一参数化借鉴自 Splatter Image，有助于收敛。

\textit{训练目标。}\;总损失函数为特征提取器输出和每层解码器输出的渲染图像与真实图像之间的 MSE 损失加 LPIPS 感知损失的加权组合：

\begin{equation}
\label{eq:leangaussian_loss}
\mathcal{L} = \frac{1}{|D|}\sum_{I_{\mathrm{gt}} \in D}\left(\lambda_E \mathcal{L}_{\mathrm{MSE}}(\hat{I}^E, I_{\mathrm{gt}}) + \sum_{l=1}^{L}\left(\lambda_D \mathcal{L}_{\mathrm{MSE}}(\hat{I}^D_l, I_{\mathrm{gt}}) + \lambda_{\mathrm{LPIPS}} \mathcal{L}_{\mathrm{LPIPS}}(\hat{I}^D_l, I_{\mathrm{gt}})\right)\right)
\end{equation}

其中 $\hat{I}^E$ 是特征提取器直接渲染的图像，$\hat{I}^D_l$ 是第 $l$ 层解码器渲染的图像，$\lambda_E$、$\lambda_D$、$\lambda_{\mathrm{LPIPS}}$ 为各项损失的权重系数。

\textit{设计动机总结。}\;LeanGaussian 的核心动机来自一个简单观察：像素-高斯固定对应关系本质上是一种"以输入空间布局为中心"的设计，它假设输入图像的像素结构与 3D 物体的几何结构之间存在直接映射。然而，实际物体表面的几何复杂度与输入像素的空间分布并不匹配——物体某些区域需要更多高斯来精细建模（如边缘、纹理丰富的区域），而背景像素对应的那些高斯则完全浪费了。通过用可学习查询替代像素映射，LeanGaussian 将高斯分布的控制权从输入布局交还给网络学习，使高斯数量和位置可以自由适配目标物体的真实几何需求。配合 3D 参考点投影的可变形注意力机制，每个高斯能够精确地从图像中提取最相关的特征信息。

\textit{注解}：LeanGaussian 在 ShapeNet SRN 数据集上实现 25.44 dB PSNR，在 GSO 数据集上实现 22.36 dB PSNR，3D 重建速度为 7.2 FPS，渲染速度达 500 FPS。每个物体使用 $N$ 个可学习查询（数量可配置），相比 Splatter Image 的固定像素对应策略，LeanGaussian 在保持前馈推理速度的同时显著提升了重建质量，标志着图像到 3D 前馈重建路线从"像素绑定"到"自由分布"的重要演进。

\paragraph{Hierarchical Splatter Image: 层级像素对应重建}

与 LeanGaussian 的无对应关系设计形成鲜明对比，针对 Splatter Image 像素→单一高斯一对一映射无法重建被遮挡背面区域的局限，Hierarchical Splatter Image（2024.05）~\cite{2405_20310} 保留了像素对应关系，但将其从一对一推广为一对多的层级结构：每个像素对应一个父高斯和若干子高斯（child Gaussian）。子高斯通过轻量级条件 MLP 生成，其输入为父高斯的图像特征 $G^{\mathrm{Parent}}(I, P)$ 和目标相机视角的嵌入。层级结构的设计动机在于：原始的 Splatter Image 虽然能精确预测可见区域的高斯参数，但对输入视角中被遮挡的背面细节重建效果较差。子高斯专门负责补充这些不可见区域，通过显式编码父高斯与目标相机位姿之间的空间关系来推断遮挡部分。该模型采用两阶段训练：前 10K 次迭代仅训练 U-Net 以确保父高斯的合理预测，之后加入条件 MLP 进行端到端联合训练。在 ShapeNet-SRN 的 Chair 数据集上，Hierarchical Splatter Image 在 PSNR、SSIM 和 LPIPS 三项指标上均达到最优，尤其在遮挡区域的重建精度显著优于 Splatter Image。

\paragraph{GSD: 视角引导扩散采样重建}

针对前馈重建方法直接回归高斯参数无法建模分布先验的局限，GSD（2024.07）~\cite{2407_04237} 则从生成建模的角度出发，利用 Diffusion Transformer（DiT）学习3D高斯基元的无条件分布先验。与直接回归高斯参数的前馈方法不同，GSD 的核心创新在于视角引导采样（view-guided sampling）：在推理时，将输入图像通过 GS 的可微渲染（splatting）函数反向传播像素级梯度到3D空间，用这些梯度条件化 GS DiT 的去噪过程。这种细粒度条件化机制避免了传统图像编码器的信息压缩损失，且天然支持任意相机参数。此外，GSD 还联合训练一个辅助2D扩散模型，通过3D到2D的"抛光-复用"循环（polishing and re-using）进一步提升纹理细节。在 CO3Dv2 的10个类别上，GSD 在 PSNR、SSIM、LPIPS 以及3D几何指标 F-score 和 Chamfer Distance 上均优于 NerFormer、ViewFormer 和 SparseFusion。

\paragraph{DiffusionGS: 单阶段3D扩散重建}

DiffusionGS（2024.11）~\cite{2411_14384} 提出了一种单阶段3D扩散模型，其关键设计是在每个去噪时间步直接输出像素对齐的3D高斯点云（而非2D图像），从而在扩散过程中强制执行视图一致性。每个像素对应一个高斯基元 $G_t^{(k)} = (\mu_t^{(k)}, \Sigma_t^{(k)}, \alpha_t^{(k)}, c_t^{(k)})$，其中中心位置 $\mu_t^{(k)}$ 由像素射线参数化，确保高斯的空间定位。为实现通用3D先验，DiffusionGS 提出场景-物体混合训练策略：在 Objaverse（730K 物体）和 MVImgNet（220K 物体）上训练物体先验，在 RealEstate10K 和 DL3DV10K 上训练场景先验，并设计相对位姿条件化（RPPC）统一物体和场景的相机参数化方式。推理仅需约6秒，在物体级和场景级重建上分别实现了 2.20 dB 和 1.34 dB 的 PSNR 提升，同时比已有方法快5倍以上。

\paragraph{DSplats: 隐空间高斯重建}

DSplats（2024.12）~\cite{2412_09648} 重新审视了高斯重建模型的骨干架构：传统方法使用像素级 U-Net，与预训练的 Latent Diffusion Model（如 Stable Diffusion）不兼容。DSplats 将骨干替换为预训练的 LDM，并引入 VAE 编码器/解码器在像素空间与隐空间之间映射。然而，标准 LDM 仅针对图像生成优化，因此 DSplats 对其进行多项适应性改造：添加噪声编码器处理多视角输入、在去噪网络中注入3D感知的相机位姿条件，以及通过可微高斯光栅化将隐空间的中间特征解码为显式3D高斯。整个框架以单次前向推理的方式联合完成多视角去噪和3D重建。在 GSO 数据集上达到了 20.38 dB PSNR 和 0.842 SSIM 的重建质量。

\paragraph{DreamLifting: 即插即用多视角提升}

DreamLifting（2025.09）~\cite{2509_07435} 作为即插即用模块，将多视角扩散模型提升为 3D 资产生成器。

%---------------------------------------------------------------------
\subsubsection{优化式重建路线}
\label{sec:08_2_1_2}

本路线介绍先利用多视角扩散模型生成多视角图像作为监督信号，再通过可微分渲染迭代优化3DGS参数的方法，生成质量通常高于前馈路线但需要更多优化时间。

优化式重建方法的核心思路是：先利用多视角扩散模型（multi-view diffusion model）从输入图像生成多个视角的图像，然后以这些多视角图像作为监督信号，通过可微分渲染优化3DGS参数。这类方法通常能够生成质量较高的结果，但需要一定的优化时间。

\paragraph{锚点方法：图像提升重建}

前文介绍的前馈重建路线（如 LeanGaussian、DiffusionGS）通过单次网络推理直接输出3DGS参数，速度快但受限于训练数据的几何先验——当输入视角稀缺时，不可见区域的几何只能靠网络"猜"，质量天花板明显。优化式路线走了另一条路：不直接预测3D，而是先用扩散模型生成多视角图像作为"伪监督"，再通过迭代优化将这些2D观测提升为3D。LiftImage3D（2024.12）~\cite{2412_09597} 正是这条路线的代表性工作，它将视频扩散先验与3DGS重建统一为端到端管线，系统性地解决了视频扩散先验用于单图3D重建时的三大挑战：（1）大相机运动下视频生成质量退化、（2）生成帧缺乏精确的相机位姿、（3）扩散去噪过程引入的几何扭曲破坏3D一致性。

\textbf{核心改动}：传统的"视频到3D"管线直接用运动可控的视频扩散模型（LVDM，Latent Video Diffusion Model，潜在视频扩散模型）生成完整视频序列，再用 COLMAP 估计位姿。但大相机运动会导致 LVDM 生成模糊帧，而 COLMAP 对扩散模型产生的几何不一致极度脆弱，往往直接失败。LiftImage3D 的核心思路是\textbf{化大为小、分而治之}。将大相机运动分解为多个小运动片段逐段生成（关节式轨迹策略），用鲁棒的神经匹配模型 MASt3R 替代脆弱的 COLMAP，并引入扭曲感知3DGS显式建模并吸收各帧的独立几何扭曲。

\textbf{三步管线详解}如下。

\textit{第一步：关节式轨迹策略（articulated trajectory strategy）}。直觉上，视频扩散模型擅长生成小幅度、平滑运动的视频，但在大范围相机运动时帧质量急剧下降。关节式轨迹策略借鉴机械臂的关节运动思想：不要求模型一步到位生成大范围环视视频，而是将相机轨迹沿 $D$ 个方向（实践中取 $D=4$，即上下左右）分解，每次仅沿一个方向生成 $l$ 帧（默认 $l=8$），并将前一段的最后一帧作为下一段生成的锚点图像。形式化地，设运动可控 LVDM 为 $f_v$，输入图像为 $x$，相机位姿为 $(R_i, T_i)$，则第一阶段的生成为 $x_i^* = f_v(x, R_i T_i)$（$i \in [0, 1, \ldots, l \times D]$），第二阶段以各方向终点帧为锚点继续延伸 $x_{i}^* = f_v(x_{d_i}^*, R_i T_i)$。以两段累积为例，总共生成 $L = l \times D + (l-1) \times (D-1) = 56$ 帧，实现大范围视角覆盖。

\textit{第二步：鲁棒神经匹配与点云注册}。扩散模型的去噪过程引入时间噪声，导致生成帧之间存在固有几何不一致，传统 SfM（如 COLMAP）在这种数据上极易失败。LiftImage3D 采用基于回归的神经匹配模型 MASt3R 替代传统特征匹配：MASt3R 直接从两张图像回归点云图（point map），通过 Procrustes 对齐估计相对位姿。为降低 $L$ 帧的 $O(L^2)$ 匹配复杂度，该方法利用视频帧的时序先验——每帧仅与其前一帧配对，将匹配对数从 $L(L-1)$ 降至 $L-1$。所有帧共享同一相机内参，通过 Weiszfeld 算法联合优化焦距 $f^*$：

\begin{equation}
\label{eq:liftimage3d_focal}
f^* = \arg\min_f \sum_{i=0}^{W} \sum_{j=0}^{H} O^{i,j} \left\| (i^\prime, j^\prime) - f \frac{(P^{i,j,0},\, P^{i,j,1})}{P^{i,j,2}} \right\|
\end{equation}

其中 $(i^\prime, j^\prime) = (i - \frac{W}{2}, j - \frac{H}{2})$ 为以图像中心为原点的像素坐标，$O^{i,j}$ 为有效点掩码，$P^{i,j,\cdot}$ 为 MASt3R 回归的3D点坐标。优化后的焦距和各帧相对位姿共同完成点云注册，为后续3DGS提供初始化。

\textit{第三步：扭曲感知3DGS优化（distortion-aware 3DGS）}。这是该方法最关键的创新。直觉上，扩散模型生成的每一帧都带有独立的几何扭曲（如局部膨胀、边缘偏移），如果直接用这些帧优化标准3DGS，扭曲会被"烘焙"进3D表示中，导致渲染伪影。LiftImage3D 的解法是为每一帧引入独立的\textbf{扭曲场}（distortion field），将扭曲从规范3DGS中分离出来：

\begin{equation}
\label{eq:liftimage3d_distortion}
F: (G + \Delta G, \, t) \to G^*
\end{equation}

其中 $G$ 是规范高斯（canonical 3DGS），$\Delta G = (\Delta X, \Delta r, \Delta s)$ 分别表示位置、旋转和缩放的扭曲增量，$t = (t_i, t_j)$ 是沿两个方向的二维时间戳（用于区分各帧），$G^*$ 是扭曲后的高斯。扭曲场网络 $F$ 的结构借鉴 4D 高斯中的形变场设计：采用 K-Planes 将5D神经体素（3D 空间 $xyz$ + 2D 时间 $(t_i, t_j)$）分解为 9 个多分辨率平面——6 个空间平面 $\{(x,y), (x,z), (y,z)\}$ 配对 2 个时间轴 $\{t_i, t_j\}$ 各 3 个，加上 3 个纯空间平面 $(x,y)$、$(x,z)$、$(y,z)$——再由小型 MLP 融合为扭曲特征。训练完成后，扭曲增量 $\Delta G$ 被丢弃，仅保留无扭曲的规范高斯 $G$ 作为最终3D表示。

训练损失函数包含五项约束：

\begin{equation}
\label{eq:liftimage3d_loss}
\mathcal{L} = \mathcal{L}_{\mathrm{RGB}}^{L1} + \mathcal{L}_{\mathrm{LPIPS}} + \mathcal{L}_{\mathrm{depth}}^{L1} + \mathcal{L}_{\mathrm{TV}} + \mathcal{L}_{\mathrm{distort}}
\end{equation}

其中 $\mathcal{L}_{\mathrm{RGB}}^{L1}$ 和 $\mathcal{L}_{\mathrm{LPIPS}}$ 监督渲染图像与生成帧的纹理相似性，$\mathcal{L}_{\mathrm{depth}}^{L1}$ 基于校准后的深度先验约束几何（将 MASt3R 的粗但绝对深度与 Depth Anything V2 的精细但相对深度通过中值标定融合），$\mathcal{L}_{\mathrm{TV}}$ 为网格化总变分正则化防止扭曲无序增长。最关键的 $\mathcal{L}_{\mathrm{distort}}$ 约束规范空间（即输入图像对应的 $t_i = t_j = 0$）的扭曲趋于零：

\begin{equation}
\label{eq:liftimage3d_distort_reg}
\mathcal{L}_{\mathrm{distort}} = |\Delta X - 0|_{t_i=t_j=0} + |\Delta r - 0|_{t_i=t_j=0} + |\Delta s - 0|_{t_i=t_j=0}
\end{equation}

\textit{公式注解}：公式~(\ref{eq:liftimage3d_distort_reg}) 确保规范空间（对应输入图像的精确无扭曲视角）的扭曲增量为零，即规范高斯 $G$ 直接对齐输入图像。而其他帧的扭曲增量则自由学习以吸收扩散模型的几何偏差。这种设计的精妙之处在于：将"噪声"吸收到各帧独立的扭曲中，使最终的规范高斯成为干净无扭曲的3D表示。

\textit{注解}：LiftImage3D 的设计动机可以概括为\textbf{承认并建模不完美}——与其要求视频扩散模型生成完美的3D一致帧（目前做不到），不如显式建模其不完美，在优化过程中分离和吸收。这种思路与前馈路线形成鲜明对比：前馈方法假设网络已学会足够的3D先验，一次推理即出结果；优化式方法则允许通过迭代修正来弥补先验不足，代价是需要更多计算时间。在 LLFF、DL3DV 和 Tanks and Temples 数据集上，LiftImage3D 均取得了当时最优性能。

\paragraph{GSV3D: 视频扩散高斯光栅化解码}

与 LiftImage3D 利用视频扩散先验的策略不同，GSV3D（2025.03）~\cite{2503_06136} 直接利用 Stable Video Diffusion（SV3D）生成多视角图像，并设计高斯光栅化解码器（Gaussian Splatting Decoder）将 SV3D 的多视角隐变量转化为显式3D高斯表示。该解码器的核心是将视频扩散模型每帧的隐变量解码为3D高斯参数，通过可微渲染的几何蒸馏（geometric distillation）为去噪网络提供3D约束。这种设计的关键在于：多视角隐变量不再仅仅用于生成2D图像，而是被转化为可渲染的3D结构，使扩散模型在训练过程中获得3D感知能力。

\paragraph{Hi3D: 视频扩散作为3D感知序列生成器}

Hi3D（2024.09）~\cite{2409_07452} 开创性地将视频扩散模型重新定义为3D感知的序列图像生成器。该方法的核心动机是：传统多视角扩散模型（如 Zero123、SyncDreamer）仅在2D空间建立跨视图注意力，缺乏内在的3D一致性；而视频扩散模型在预训练过程中已隐式学习到时序一致性先验，天然适合建模多视角之间的几何关系。Hi3D 采用两阶段管线：第一阶段为预训练的 Stable Video Diffusion 注入相机位姿条件（Pl\"{u}cker 嵌入），生成 $512 \times 512$ 分辨率的低分辨率多视角轨道视频；第二阶段训练3D感知的视频到视频细化器（refiner），以深度图为条件将分辨率提升至 $1024 \times 1024$，这是当时图像到3D领域的最高分辨率。最后通过3DGS插值增强更多新视角，并用 SDF 重建高保真网格。

\paragraph{Vista3D: 粗精两阶段暗面探索重建}

Vista3D（2024.09）~\cite{2409_12193} 采用了粗-精两阶段策略，并提出了"3D暗面探索"（3D darkside exploration）的视角：将不可见区域的生成视为从单张图像发掘隐藏维度的问题，而非简单的重建任务。粗阶段利用3DGS快速（约30秒）生成初始几何，采用 Top-K 致密化和透射率/尺度正则化策略避免过度膨胀；精阶段直接从3DGS中提取 SDF，使用 FlexiCubes 可微分等值面表示进一步优化几何和纹理。Vista3D 的另一个核心贡献是扩散先验组合（prior composition）：通过梯度幅度约束同时使用两个独立的扩散先验——Zero-1-to-3 负责可见面的保真度，Stable Diffusion 负责不可见面的多样性——实现了生成一致性与多样性的平衡。

\paragraph{GaussVideoDreamer: 不一致性感知场景级视频重建}

在场景级视频生成方向上，GaussVideoDreamer（2025.04）~\cite{2504_10001} 提出了不一致性感知的 3DGS（Inconsistency-Aware Gaussian Splatting）方法来解决多视角扩散模型固有的几何-纹理不一致问题。该方法首先通过 warp-and-inpaint 管线用遮挡体积（occlusion volume）技术构建粗场景，然后引入不一致性感知机制：对参考视角的已验证点与新生成点的渲染结果分别建立3D一致性掩码，以此指导后续去噪过程。生成的新视角被分为参考点和新增点两类，通过视角特定的细化掩码（refinement mask）确保可靠区域保持不变、不确定区域被扩散模型修复。该机制实现了32\%更高的 LLaVA-IQA 分数和至少2倍的加速。

\paragraph{ExScene: 360度沉浸式场景重建}

ExScene（2025.03）~\cite{2503_23881} 设计了两阶段管线以实现360度沉浸式场景重建：第一阶段通过多模态扩散模型（融合 LLaVA 语义描述和视觉特征）生成高保真全景图，并提出消除投影畸变的正则化全景深度估计方法初始化3DGS；第二阶段引入高斯补丁（Gaussian Patching）模块，将相机轨迹一致性和颜色-几何先验注入 Stable Video Diffusion 去噪过程，修复初始化3D模型中的视角遮挡伪影。

\paragraph{ERGO: 超额风险引导多目标优化}

ERGO（2026.02）~\cite{2602_10278} 从统计学习角度提出超额风险引导（Excess-Risk-Guided Optimization），将3DGS优化构建为多目标加权优化问题，通过将每个视角的损失分解为衡量当前参数与最优参数差距的超额风险和不可约噪声的贝叶斯误差，动态调整各视角-目标的损失权重，以 Pareto 最优的方式平衡几何感知和纹理感知目标。

%---------------------------------------------------------------------
\subsubsection{大规模重建基础模型路线}
\label{sec:08_2_1_3}

本路线介绍在大规模3D数据上训练通用前馈重建网络（Large Reconstruction Model, LRM）的方法，使其能够从单张图像直接预测3D表示。

随着训练数据和模型规模的增长，研究者开始探索大规模重建基础模型（Large Reconstruction Model, LRM），即在大规模3D数据上训练通用的前馈重建网络，使其能够从单张图像直接预测3D表示。

\paragraph{锚点方法：状态空间序列建模}

MVGamba（2024.06）~\cite{2406_06367} 是一种基于状态空间模型（State Space Model, SSM）的轻量级高斯重建模型。已有的基于 Transformer 的高斯重建模型往往面临多视角不一致和纹理模糊的问题——MVGamba 将此归因于 Transformer 架构在追求强大表达力时牺牲了多视角信息传播的效率。

MVGamba 的核心架构采用了类 RNN 的状态空间模型替代 Transformer，其高斯重建器在生成高斯序列的同时，通过因果上下文传播多视角信息：

\begin{equation}
\label{eq:mvgamba_ssm}
\mathbf{h}_t = \bar{\mathbf{A}} \, \mathbf{h}_{t-1} + \bar{\mathbf{B}} \, \mathbf{x}_t, \quad \mathbf{y}_t = \bar{\mathbf{C}} \, \mathbf{h}_t
\end{equation}

\textit{白话解释}：公式~(\ref{eq:mvgamba_ssm})描述了 SSM 的核心递推关系。隐状态$\mathbf{h}_t$在时间步$t$由前一步隐状态$\mathbf{h}_{t-1}$经矩阵$\bar{\mathbf{A}}$变换后加上当前输入$\mathbf{x}_t$经$\bar{\mathbf{B}}$变换的结果。输出$\mathbf{y}_t$（即高斯属性）由隐状态经$\bar{\mathbf{C}}$投影得到。$\bar{\mathbf{A}}$的因果性保证了多视角信息的有序传播，实现跨视图自精化（cross-view self-refinement）。

\textit{注解}：SSM 的线性复杂度使得 MVGamba 能够用仅约0.1倍的模型参数量超越基于 Transformer 的 SOTA 基线。MVGamba 统一了单图、稀疏图和文本驱动的3D生成任务。

\paragraph{3D-Adapter: 即插即用3D几何感知注入}

与 MVGamba 的 SSM 架构不同，3D-Adapter（2024.10）~\cite{2410_18974} 采用了"即插即用"的设计思路：作为一个并行分支模块将3D几何感知注入预训练的图像扩散模型。其核心机制是3D反馈增强（3D Feedback Augmentation）——区别于之前的 I/O 同步方法（将3D重建/渲染插入去噪循环），3D-Adapter 保持原始网络拓扑不变，而是在基模型旁附加一个3D感知分支：将中间多视角特征解码为3D表示，再将渲染的 RGBD 视角重新编码来增强基模型的特征。这一设计避免了对原始去噪流程的干扰和误差累积，同时避免了多视角分数平均导致的模式坍塌。该方法提供了基于高斯泼溅的快速前馈版本和基于神经场的无需训练版本。

\paragraph{Direct 3D: 对极注意力多视角生成}

Direct 3D（2024.11）~\cite{2411_10947} 则直接利用改造后的 Stable Diffusion 模型同时生成多视角深度图、RGB 图像和高斯特征图。其关键创新是在隐空间到像素空间的 VAE 解码器中引入对极注意力（epipolar attention）：对于目标视角中每个像素，仅在对应极线上采样源视角的特征，利用已知的相对相机位姿建立几何约束。这一机制确保了深度图的像素级多视角一致性——简单的跨视图注意力（cross-view attention）足以保证 RGB 图像的语义一致性，但深度图需要更严格的几何约束。来自不同视角的 RGB、深度和高斯特征被融合为3D点云，既可提取为纹理网格也可通过3DGS渲染。

\paragraph{NavCrafter: 视频扩散驱动的场景探索}

NavCrafter（2026.04）~\cite{2604_02828} 从视频扩散模型中学习3D先验，旨在解决单图场景探索中的两大挑战：可控性和一致性。其方法包含三个模块：可控新视角合成模块采用三阶段架构（双分支相机注入、LoRA 调制、注意力引导），通过 Pl\"{u}cker 射线嵌入将相机轨迹注入 DiT 去噪器，实现精确的相机控制；碰撞感知轨迹规划器迭代地探索新视角，通过可微变换将物体级坐标系与场景级坐标系对齐；增强版3DGS管线引入深度对齐监督（确保几何精度）和结构正则化（防止过拟合），并结合图像扩散细化模块提升视觉保真度。

\paragraph{One2Scene: 全景图多视角立体匹配重建}

One2Scene（2026.02）~\cite{2602_19766} 将全景图投影为多个稀疏锚点视角，将重建任务重新表述为多视角立体匹配。其方法分为三步：首先用 Hunyuan-Pano-DiT 从单图生成360度全景图以获取全局覆盖；然后提出前馈3D高斯支架（3D Geometric Scaffold）模型，为全景图的每个像素预测3D高斯参数，其中关键设计是视角投影模块——利用双向特征融合（bidirectional feature fusion）在锚点视角之间传递信息，保证跨视图一致性；最后以该几何支架为先验，驱动合成网络在任意相机位姿下生成照片级真实的新视角。这种"先建支架、后画细节"的策略有效解决了大视角变化下全局语义漂移和几何不一致的问题。

\paragraph{UnCommon Objects in 3D: 罕见物体类别3D生成}

UnCommon Objects in 3D（2025.01）~\cite{2501_07574} 专注于罕见物体类别的 3D 生成，通过分布外泛化扩展了重建模型的适用范围。

%---------------------------------------------------------------------
\subsubsection{高质量与高保真重建路线}
\label{sec:08_2_1_4}

本路线介绍在已有框架基础上引入额外约束机制或细化策略的方法，旨在提升生成结果的几何精度、纹理细节和物理合理性。

随着前馈和优化式方法的基本框架趋于成熟，研究重点逐渐转向如何提升生成结果的几何精度、纹理细节和物理合理性。这类方法通常在已有框架基础上引入额外的约束机制或细化策略。

\paragraph{锚点方法：RGB-N 分支重建}

GS-RGBN（2025.04）~\cite{2504_01512} 的核心洞察是：当前单视图3D生成方法同时受到两个问题的困扰——2D图像的几何歧义（geometric ambiguity）和3DGS缺乏结构化表示。该论文提出了一种混合体素-高斯表示（hybrid Voxel-Gaussian representation）来同时缓解这两个问题。

GS-RGBN 的关键设计在于用3D体素表示显式编码几何信息，并将高斯基元锚定在体素结构上。其融合模块将 RGB 特征和表面法线（surface normal）特征对齐：

\begin{equation}
\label{eq:gsrgbn_fusion}
\mathbf{V} = \mathrm{Conv3D}\!\left(\mathrm{Concat}\!\left[\mathbf{F}_{\mathrm{RGB}}, \, \mathbf{F}_{\mathrm{Normal}}\right]\right), \quad \mathbf{G}_k = f_{\theta}\!\left(\mathbf{V}(\boldsymbol{\mu}_k)\right)
\end{equation}

\textit{白话解释}：公式~(\ref{eq:gsrgbn_fusion})中，$\mathbf{F}_{\mathrm{RGB}}$和$\mathbf{F}_{\mathrm{Normal}}$分别是从2D图像估计的RGB特征和法线特征，在3D空间中拼接后通过3D卷积融合为体素特征$\mathbf{V}$。然后，每个高斯基元$\mathbf{G}_k$的属性由其中心$\boldsymbol{\mu}_k$位置的体素特征经网络$f_\theta$预测得到。

\textit{注解}：体素表示提供了显式的3D几何先验，消除了2D图像的几何歧义；同时体素结构约束了高斯的学习过程，使优化更容易找到好的局部最优解。

\paragraph{PhyCAGE: 物理仿真增强组合式资产生成}

与 GS-RGBN 关注几何结构不同，PhyCAGE（2024.11）~\cite{2411_18548} 聚焦于组合式3D资产生成（compositional 3D asset generation）——当输入图像包含两个交互物体（如"一只穿着毛衣的青蛙"）时，需要分别生成各组件的3D表示并确保其物理关系合理。该方法引入了物理仿真增强的分数蒸馏采样（PSE-SDS），将 SDS 损失的梯度作为物理仿真的初始速度，让可微仿真器充当物理引导的优化器，逐步修正高斯基元的位置至物理兼容状态（如确保支撑关系、避免穿透）。具体流程为：先用 Grounded-SAM 分割前景/背景并分别用 SyncDreamer 生成多视角图像，再分别重建两个3DGS模型，最后通过 PSE-SDS 优化物理合理性。这是首个从单图生成物理合理的组合式3D资产的方法。

\paragraph{RIGI: 不确定性感知伪影消除}

RIGI（2024.11）~\cite{2411_18866} 从不确定性感知学习的角度解决多视角扩散模型输出不一致导致的3DGS伪影问题。该方法包含两步：不确定性估计——同时优化两个独立的高斯模型，通过计算它们渲染结果之间的绝对差异生成逐像素不确定性图，差异越大意味着该区域的多视角一致性越差；不确定性正则化——基于不确定性图动态调整逐像素的优化强度，对高不确定性区域（通常是物体边缘、遮挡边界等歧义区域）降低损失权重，从而避免冲突的梯度信号导致伪影。此外，RIGI 还引入了深度正则化和法线平滑等额外约束进一步提升几何质量。

\paragraph{ScalingGaussian: 3D/2D扩散模型协同细节提升}

ScalingGaussian（2024.07）~\cite{2407_19035} 结合3D和2D扩散模型的优势来提升3D内容生成的细节质量。其管线分为三步：首先用3D扩散模型（基于 PC$^2$ 的点云扩散）生成初始稀疏点云，通过局部密度加权选择进行致密化以解决点云稀疏导致的效率问题；然后用2D扩散模型的 SDS 损失引导高斯的克隆和分裂，逐步丰富3D表示的几何和纹理；最后将3DGS转换为纹理网格，并引入 GPP（Gaussian-Projected Polygon）损失优化表面纹理——该损失将网格的纹理质量约束直接与扩散模型的判断对齐，显著提升了纹理边缘的清晰度。

\paragraph{ContrastiveGaussian: 对比学习纹理细化}

ContrastiveGaussian（2025.04）~\cite{2504_08100} 将对比学习引入3DGS生成过程以解决 SDS 优化导致的纹理模糊和几何不一致问题。其核心思路是：利用2D扩散模型为新视角生成参考图像，通过超分辨率增强输入质量后，将渲染图像与参考图像构成正样本对、与无关图像构成负样本对，用感知损失（perceptual loss）引导高斯参数优化。该方法还提出了 Quantity-Aware Triplet Loss（QA-Triplet Loss），根据正负样本在特征空间中的距离动态调整 triplet 损失的边际，解决了训练中正负样本数量不均衡的问题。最后通过扩散模型的多步 MSE 损失细化粗纹理。

\paragraph{DreamSat: 航天器领域适配重建}

DreamSat（2024.10）~\cite{2410_05097} 针对航天器这一垂直领域，通过领域适配实现高质量单图到3D重建。该方法首先构建了包含190个高质量航天器3D模型的专用数据集，在此基础上微调 Zero123 XL 的多视角生成能力；然后将其集成到 DreamGaussian 的 SDS 优化框架中，利用3DGS作为3D表示。DreamSat 还针对航天器的特殊几何特征（如太阳能板、天线等细长结构）设计了专项优化策略。这是将通用图像到3D管线适配到特定工程领域的代表性工作。

\paragraph{ERGO: 超额风险引导几何-纹理双目标优化}

ERGO（2026.02）~\cite{2602_10278} 在高质量路线中将超额风险分解与几何-纹理双目标优化相结合。具体而言，ERGO 定义了两组损失函数：几何感知目标（通过边缘检测聚焦于几何一致性关键区域）和纹理感知目标（通过高通滤波聚焦于纹理丰富的细节区域）。对每个辅助视角，ERGO 估计其相对于参考视图的超额风险——超额风险越高，说明该视角的伪标签噪声越大——并据此降低该视角在优化中的权重。这种自适应权重调整机制使优化过程能自动聚焦于高质量视角，在 GSO 和 OmniObject3D 数据集上显著优于均匀加权的基线方法。

\paragraph{GaussianZoom: 渐进式放大生成3D重建}

GaussianZoom（2026.05）~\cite{2605_18252} 提出渐进式放大生成3D重建框架，其核心洞察是：极限放大本质上是渐进式生成过程而非单次上采样。在每个放大步骤中，GaussianZoom 使用基于深度的特征变形（depth-based feature warping）替代传统的光流变形——利用3DGS重建的可靠深度图，通过精确的几何投影建立跨视图像素对应关系，有效解决了大视差和遮挡问题。同时引入 VLM（视觉-语言模型）驱动的语义条件化，为超分辨率模块提供高层语义引导，使其能合成超出输入分辨率的新细节而非简单锐化。该方法还设计了可扩展的连续细节层次（LOD）结构，动态调制不同尺度高斯基元的可见性，实现了无锯齿的平滑缩放渲染。

%---------------------------------------------------------------------
\subsubsection{其他单图到3D方法}
\label{sec:08_2_1_5}

本路线介绍从特殊角度或特殊场景出发的单图到3D方法，包括视频扩散先验、像素级对齐、生成式补全、全景场景重建、可编辑控制等差异化路线。

除上述技术路线外，还有一些方法从特殊角度或特殊场景出发，提出了差异化的解决方案。

\paragraph{视频扩散先验路线} FlashDreamer（2025.03）~\cite{2503_00726} 利用预训练的视觉-语言模型（LLaMA-3.1-8B）为输入图像生成描述性文本提示，引导 Stable Diffusion 从不同视角生成图像，然后通过3DGS合并与对齐（merge and alignment）过程将各视角的3D表示融合为完整场景。核心流程为：先用 Flash3D 从单图构建初始3D场景，渲染新视角后用扩散模型修复空白区域，然后将修复后的图像重建为新的3DGS并与已有场景对齐合并。为保持多视角一致性，该方法在所有视角统一使用同一文本描述作为扩散提示。该流程迭代执行，无需额外训练即可从单图逐步扩展为综合3D场景。

\paragraph{像素级对齐路线} Image2GS（2025.07）~\cite{2507_00916} 将图像到3D问题解耦为"提升"（lifting，将可见区域转换为3D）和"补全"（completion，生成不可见区域）两个子任务，专注于确定性的提升任务。其核心观察是：判别性模型被强行用于生成性任务时，会预测多种可能结果的"均值"，导致输出模糊。通过将任务限制在可见区域，判别性模型可以输出确定性预测。具体实现中，Image2GS 基于 DepthAnythingV2（ViT 骨干）构建，通过预优化每个场景的3DGS推导逐视图可见性掩码（visibility mask），在训练中排除源视角中不可见区域的损失，显著提升了可见区域的重建质量。

\paragraph{生成式补全路线} Complete Gaussian Splats（2025.08）~\cite{2508_21542} 提出用生成式而非回归式的方法处理被遮挡和超出视锥的区域。其核心创新是解决了"无真值3D数据的扩散模型训练"问题：由于不存在大规模的3D Splatter Image 数据集，该方法设计了变分自重建器（Variational AutoReconstructor），仅从2D图像数据集（如 CO3D、ScanNet）通过可微渲染（前向模型）自监督学习 Splatter Image 的隐空间分布。随后在此隐空间上训练扩散模型，实现多样化的3D补全。与回归式方法只能输出模糊的"均值"预测不同，该生成式方法能为同一输入图像产出多种合理的3D补全结果。

\paragraph{全景场景重建路线} Scene4U（2025.04）~\cite{2504_00387} 从全景图出发，结合开放词汇分割模型和大语言模型将全景图分解为多层（如前景物体、背景墙壁、地板等），通过分层修复模块恢复被遮挡区域，再用3DGS进行分层优化，最终生成支持自由探索的沉浸式3D场景。该方法的关键创新在于分层策略：将复杂的全景3D重建分解为多个较简单的子任务，每层独立处理遮挡关系和几何约束，最后合成为完整场景。在 LPIPS 上提升了24.24\%，BRISQUE 上提升了24.40\%。

\paragraph{可编辑和可控制路线} Dragen3D（2025.02）~\cite{2502_16475} 引入锚点-高斯变分自编码器（Anchor-GS VAE）和种子点驱动策略，使用户能直接在3D空间编辑生成结果的几何形状。具体流程为：首先训练 Anchor-GS VAE 将复杂3D信息编码为紧凑隐空间表示；然后从输入图像生成一组稀疏3D种子点作为粗几何骨架；接着通过整流流模型（Rectified Flow）将种子点信息映射到 Anchor-GS VAE 的隐空间，解码为3DGS。用户通过拖拽种子点即可改变最终生成的3D形状——种子点的变形会更新对应的隐编码，从而产生几何一致的变形输出。这种"稀疏控制点驱动"的交互方式极大地降低了3D内容编辑的门槛。3DisGS（2025.04）~\cite{2504_04190} 首次实现了无监督可解释的3DGS重建。该方法提出双分支重建架构：几何分支使用风格化折叠模块（style-based folding module）从2D网格生成3D点云，外观分支基于三平面（triplane）表示生成纹理特征，两者结合实现粗粒度的几何-外观解耦。在细粒度层面，3DisGS 通过解耦表示学习（DRL）的编码器-适配器模块，自动发现3D场景中潜在语义因素（如"家具""装饰""人物"等）的独立表示。损失函数中引入互信息最小化，确保各语义因素之间相互独立。

\paragraph{SMAL-pets: 参数化模型驱动的动物化身生成}

SMAL-pets（2026.03）~\cite{2603_17131} 将3DGS与 SMAL（Skinned Multi-Animal Linear）参数化模型结合，从单张图像生成高质量、可动画化的动物（狗）化身。该方法采用混合表示：SMAL 网格提供结构先验（形状、姿态、可变形顶点偏移），3DGS 提供高保真渲染。重建分为两个阶段——网格约束的初始优化阶段用 SMAL 顶点位置初始化高斯中心，无约束的体积细化阶段释放高斯以捕捉超出网格分辨率的细节。蒙皮过程将高斯基元绑定到 SMAL 骨骼上，通过操纵 SMAL 参数即可驱动化身执行新姿态。

\paragraph{GenMOJO: 多物体遮挡物体分解生成}

GenMOJO（2025.06）~\cite{2506_12716} 则针对多物体遮挡的单目视频，通过物体分解策略将场景拆分为独立的前景物体和背景，每个对象用独立的可变形3DGS表示。优化时同时执行帧中心的渲染误差优化（联合 splatting 所有对象的高斯并计算与输入帧的差异）和对象中心的 SDS 优化（从不可见视角渲染单个对象并用扩散先验引导），通过可微变换桥接两个坐标系。

\paragraph{GEN3D: 域无关单图3D场景生成}

GEN3D（2025.11）~\cite{2511_14291} 提出域无关的单图 3D 场景生成方法，从任意领域图像生成 3DGS 场景。

\begin{takeaway}
图像到3D生成从单张或多张图像重建3DGS场景，核心挑战是从有限视角推断完整3D结构。大规模前馈模型（LRM、LGM）通过Transformer架构实现秒级生成，多物体遮挡场景需要物体分解策略。GenMOJO 的双优化范式（帧中心渲染误差+对象中心 SDS）是解决遮挡问题的有效方案。
\end{takeaway}

%=====================================================================

\subsection{近期方法脉络：生成先验与稀疏重建、Diffusion-aligned 3D Gaussian generation、Pixel-space diffusion for Gaussian splat generation等}

这一组近期工作延续 3DGS、3DGS generation，主要把改动集中在增量改进、扩展、生成路线、应用等环节。下面只保留相对既有路线的变化，便于和前文方法直接对照。

\paragraph{FlowObject}
FlowObject（2026.06）~\cite{2606_19019} 延续 3DGS 在“生成先验与稀疏重建”上的基本设定，主要改变是：FlowObject：用flow-matching生成先验补全未观测几何，并结合3DGS优化保持可见表面保真。3DGS 的显式高斯表示和可微渲染能力保持为基础。
其设计动机是：针对生成模型synthetic bias和3DGS不可见面缺失的互补局限，提出生成先验受观测流引导的对象重建框架。
具体而言，方法以稀疏视角观测约束生成式对象先验，再利用3DGS优化保留可见表面细节，从而补全不可见几何。
适用边界是：生成先验可能仍覆盖观测证据；对象类别和纹理复杂度影响补全可靠性。

\paragraph{FLUX3D}
FLUX3D（2026.06）~\cite{2606_24874} 延续 3DGS generation 在“Diffusion-aligned 3D Gaussian generation”上的基本设定，主要改变是：扩散对齐的稀疏体素表示与 2D-3D token 对齐机制。3DGS 的显式高斯表示和可微渲染能力保持为基础。
其设计动机是：识别语义 2D 特征和稀疏 3D latent 之间的表示瓶颈，并设计可扩展机制对齐图像 token 与稀疏体素 latent。
具体而言，使用更偏重重建线索的 2D 表示构建稀疏体素 latent，并在扩散 transformer 中加入密集 2D token 到稀疏 3D voxel latent 的对齐模块，最终生成 3D Gaussian 表示。
适用边界是：依赖大规模扩散先验和稀疏体素设计，可能对训练成本、输入域泛化和复杂拓扑存在限制。

\paragraph{PixGS}
PixGS（2026.07）~\cite{2607_01803} 可作为“Pixel-space diffusion for Gaussian splat generation”中的新锚点候选。相较于 3DGS generation，主要改变是：PixGS：避免复杂 cascade 和 latent-space 瓶颈的 pixel-space diffusion 3DGS 生成。3DGS 的显式高斯表示和可微渲染能力保持为基础。
其设计动机是：针对高质量 3D 数据稀缺、2D 生成器视角不一致和复杂级联难扩展的问题，提出更直接的 3DGS 生成路径。
具体而言，不同于把大规模 text-to-image latent diffusion 复杂适配到 3DGS 的方法，PixGS 通过 pixel-space diffusion 直接生成 Gaussian splat assets，以降低压缩 latent 和多组件串联造成的 artifacts。
这一锚点判断来自当前分类中的机制关系：候选锚点：生成路线，且相关性 0.95；后续方法可围绕 直接生成 3D Gaussian Splats 而非通过复杂级联间接构造。 继续比较和扩展。
适用边界是：仍受训练数据规模、扩散采样成本和 3D 监督质量限制；复杂拓扑和高频细节需进一步验证。

\paragraph{OrbitForge}
OrbitForge（2026.06）~\cite{2606_24799} 延续 3DGS 在“Video-prior anchored Gaussian scene generation”上的基本设定，主要改变是：冻结视频先验生成轨道视频，再用 Gaussian Splatting 重建作为几何锚点。3DGS 的显式高斯表示和可微渲染能力保持为基础。
其设计动机是：提出 reconstruction-anchored video synthesis，把初始 3DGS 重建作为锚点来改善视频多视角一致性并生成可探索场景。
具体而言，先由文本到视频模型生成闭环或近闭环视频，再通过 Gaussian Splatting 优化得到初步 3D 重建，并将重建反馈给后续视频合成和一致性约束。
适用边界是：质量受视频生成先验、相机轨道覆盖和初始 3DGS 重建稳定性限制；复杂开放场景可能仍有漂移。
\section{4D动态生成}
\label{sec:08_3}

4D动态生成是3D内容创建领域的前沿方向，其目标是同时建模三维空间的几何外观和随时间变化的动态运动。相比静态3D生成，4D生成需要额外处理时间维度的一致性：生成的3D模型不仅要在空间上多视角一致，还要在时间上运动连贯、无闪烁。基于3DGS的4D生成方法利用高斯基元的显式表示优势，通过形变场（deformation field）、时空建模（spatio-temporal modeling）等机制实现动态内容的生成。前面两节讨论了从文本或图像到静态3D内容的生成。然而，现实世界是动态的。要让生成结果"动起来"，需要在3DGS的基础上引入时间维度。

本节将这些方法按技术路线组织为三个类别：（1）视频/文本驱动的4D生成；（2）动态3DGS生成与形变建模；（3）运动迁移与动画生成。

%---------------------------------------------------------------------
\subsection{4D动态生成}
\label{sec:08_3_1}

本节讨论3条技术路线：（1）视频/文本驱动的4D生成；（2）动态3DGS生成与形变建模；（3）运动迁移与动画生成。

%---------------------------------------------------------------------
\subsubsection{视频/文本驱动的4D生成}
\label{sec:08_3_1_1}

本路线介绍利用视频扩散模型生成多视角动态视频序列或通过文本提示经分数蒸馏引导4D模型优化的方法，核心挑战在于同时保证空间一致性和时间一致性。

视频和文本驱动的4D生成是当前最主流的技术路线。其基本思路是：利用视频扩散模型生成多视角的动态视频序列，或者利用文本提示通过分数蒸馏引导4D模型的优化。核心挑战在于如何同时保证空间一致性和时间一致性。

\paragraph{锚点方法：4D 扩散蒸馏}

DreamGaussian4D（2023.12）~\cite{2312_17142} 是首个将3DGS引入4D生成的高效框架，能够从单张图像在约10分钟内生成14帧的4D动画（含纹理细化约10分钟，不含则约6.5分钟）。在此之前的 MAV3D、4D-fy 等方法依赖隐式 NeRF 表示和视频 SDS 优化，生成单个4D场景需要数小时，且难以有效控制生成的运动。DreamGaussian4D 的核心洞察是：利用显式4D表示（3DGS + 形变场）替代隐式 NeRF，用驱动视频的显式监督替代视频 SDS 的隐式蒸馏，从而同时实现高效率和运动可控性。整个框架分为两大阶段。

（1）图像到4D GS生成阶段——又细分为静态生成和动态生成两步。

静态生成：先用 DreamGaussianHD 从输入图像生成静态3DGS。DreamGaussianHD 是对 DreamGaussian 的工程改进，包含两个关键实践：一是多视角优化，每步采样16个视角而非1个随机视角进行 SDS，显著改善未观测区域的几何质量（代价是显存和耗时增加）；二是固定白色背景，因为大多数3D感知图像扩散模型的训练对象使用白色背景，黑色背景会引入额外噪声导致模糊。静态优化约2分钟（500次迭代）。

动态生成：利用 HexPlane 形变场为静态3DGS 引入时间维度的运动。具体流程如下：

首先，用图像到视频扩散模型（如 Stable Video Diffusion）从输入图像生成一段驱动视频 $I_{Ref,1:T} = f_\theta(\epsilon; I_{Input})$。驱动视频的概念是该方法的核心创新之一。与 MAV3D、4D-fy 等直接使用视频扩散模型 SDS 损失的方法不同，DreamGaussian4D 将视频作为显式监督信号，用户可以通过选择不同随机种子或不同驱动视频来灵活控制4D运动的类型和风格，实现了运动的可控性和多样性。

然后，训练 HexPlane 形变场来预测每个高斯基元在位置 $(x,y,z)$ 和时间戳 $t$ 下的位移、旋转和缩放变化。HexPlane 将4D时空场分解为六个特征平面（跨每对坐标轴），将4D表示的复杂度从 $O(N^4)$ 降至 $O(N^2)$：

\begin{equation}
\label{eq:hexplane}
\Delta \boldsymbol{\mu}(x,y,z,t) = \sum_{(i,j) \in \mathcal{P}} \mathbf{M}_{ij}\!\left(\mathrm{proj}_{ij}(x,y,z,t)\right)
\end{equation}

\textit{白话解释}：公式~(\ref{eq:hexplane})描述了HexPlane如何计算高斯基元在时空位置$(x,y,z,t)$的位移$\Delta \boldsymbol{\mu}$。$\mathcal{P}$是平面对的集合（如 XY-T 表示将$(x,y)$投影到一个平面、$t$投影到另一个平面），$\mathbf{M}_{ij}$是每对平面上的特征矩阵，通过双线性插值查询投影位置的特征并求和得到最终位移。

\textit{注解}：HexPlane 的分解策略不仅降低了参数量，更重要的是将4D场表示为一组可学习的4D基函数的加权和，这天然地对4D特征施加了正则化，保证了运动的局部空间刚性和时间平滑性。通过控制 HexPlane 的空间和时间分辨率（默认 $H=32, T=32$），可以调节运动的刚性与平滑程度。论文的消融实验表明，在纯 MLP、逐帧独立3DGS、HexPlane 三种变形表示中，HexPlane 始终取得最优结果——MLP 容易产生局部抖动，逐帧拟合则缺乏时间一致性且耗时极长。

形变场优化采用双重视角策略：在参考视角下用驱动视频的 MSE 损失提供运动监督，在其他视角用 Zero-1-to-3-XL 的 SDS 损失传播运动到未观测区域。一个关键训练技巧是零初始化（zero-initialization）：变形网络的最后几层线性层初始化为零输出（通过残差连接确保梯度非零），使得训练初始时动态模型与静态模型一致，避免因初始化偏移导致的颜色分布异常（如静态阶段熊猫背面为黑白相间，不良初始化会使动态优化后变为全黑）。此外，得益于静态模型的质量，SDS 可以从较低的噪声水平 $T_{max}=0.5$ 开始（而非通常的 0.98），加速收敛。动态优化约4.5分钟（200次迭代）。

（2）视频到视频纹理细化阶段——生成的4D GS 可以导出为动画网格序列，但 SDS 的固有模糊性会导致网格纹理不够清晰。纹理细化模块在 UV 空间操作：先沿一条相机轨迹渲染视频，添加噪声水平 $T=0.7$ 的噪声，再用图像到视频扩散模型去噪，通过 MSE 损失反向传播优化所有时间步的纹理贴图。与 DreamGaussian 的逐帧独立细化不同，视频到视频的 pipeline 天然保持时间一致性，避免了相邻帧闪烁。此阶段约3.5分钟（50次迭代）。

\textit{整体性能}：在单张80GB A100 GPU上，完整流程约10分钟。用户研究（38名参与者）显示，DreamGaussian4D 在图像对齐、3D外观和运动质量三个维度上均显著优于 MAV3D 和 Animate124。CLIP 图像相似度指标也显示出明显优势。在视频到4D的定量评测基准上，两种变体（Ours-Fast: 16帧/500迭代，Ours: 32帧/500迭代$\times$4视角）均以远短于基线的时间取得更优结果。

\paragraph{Diffusion4D与EG4D: 无蒸馏4D生成}

与 DreamGaussian4D 的 SDS 优化策略不同，Diffusion4D（2024.05）~\cite{2405_16645} 和 EG4D（2024.05）~\cite{2405_18132} 提出无需分数蒸馏（score distillation）的4D生成方法。Diffusion4D 利用视频扩散模型的时序一致性，将其迁移为4D生成所需的时空一致性，在精心策划的动态3D数据集上训练4D感知视频扩散模型，生成动态3D资产的轨道视角图像，然后通过粗到细的显式4D构建过程在几分钟内生成高保真4D资产。EG4D 同样不依赖 SDS，而是通过显式生成多视角视频，利用注意力注入策略保证时序一致，再基于3DGS的鲁棒动态重建方法生成高质量4D资产。

\paragraph{PLA4D: 显式像素级对齐4D生成}

PLA4D（2024.05）~\cite{2405_19957} 从另一个角度解决多先验冲突问题。该方法观察到多个SDS损失的梯度方向可能相反——这在数学上等价于在潜空间中强制折中两个相互冲突的目标，导致4D模型陷入两难。核心思想是：用显式参照系将冲突"外化"到像素空间，从而绕过了潜空间中等价方向无法兼容的根本困境。然而，锚点视频的质量成为新的瓶颈——若锚点存在运动模糊或几何畸变，对齐误差会直接传播至4D目标。为此，PLA4D 提出从隐式潜空间对齐转向显式像素级对齐：以文本生成的视频作为锚点参考，在像素空间中统一不同扩散模型引导的渲染过程。静态层面，通过焦距对齐估计每帧对应的焦距长度，配合高斯-网格对比学习从未知视角获取几何监督；动态层面，通过运动对齐将锚点视频的运动传递到4D目标，配合T-MV细化方法修复动态多视角渲染中的表面裂缝。PLA4D 完全基于开源模型构建，每个样本仅需15分钟（0.6K次迭代）即可生成。

\paragraph{STP4D: 扩散模型直接生成4DGS}

STP4D（2025.04）~\cite{2504_18318} 则代表了一条全新的技术路线——首次用扩散模型直接生成4DGS，而非通过SDS优化或视频重建。该方法的核心设计是时空-提示（spatio-temporal-prompt）三重一致性建模，集成了三个关键模块：时变提示嵌入（TPE）将文本特征映射为逐帧差异化提示信息，通过交叉注意力嵌入每帧高斯中，确保生成内容与文本高度对齐；几何信息增强（GIE）利用 K-Planes 将4D高斯分解为三组低秩矩阵，通过 GroupFormer 提取组间和组内的时空特征，增强几何表示；时间扩展形变（TED）以锚帧为基础，通过可学习权重池扩展到完整帧序列，保证运动连贯。该方法基于 DDIM 采样直接从噪声生成4DGS，每个资产生成仅需约4.6秒（24帧），比 Diffusion4D 快约100倍。

%---------------------------------------------------------------------
\subsubsection{动态3DGS生成与形变建模}
\label{sec:08_3_1_2}

本路线介绍通过形变场、时空建模等机制实现动态3D高斯生成的方法，通常先构建静态3D模型再通过形变场驱动其运动。

动态3DGS生成方法的核心在于如何有效地表示和建模3D高斯基元随时间的形变。与直接生成4D内容不同，这类方法通常先构建静态3D模型，再通过形变场驱动其运动。

\paragraph{锚点方法：4D 时空重建}

4DSTR（2025.11）~\cite{2511_07241} 针对基于可变形4D高斯的生成管线中两个根本性缺陷，提出了时空校正（spatial-temporal rectification）机制。

\textbf{Baseline 做法及其局限。}前述可变形4D Gaussian Splatting（2024）~\cite{wu20244d} 将3D高斯扩展到时间域，其形变场 $F(S, t) = [X_t, s_t, r_t, o_t, c_t]$ 为每个时间步独立预测位置、缩放、旋转等属性。然而，这种逐帧独立处理存在两个问题：其一，缺少显式的时序关联——相邻帧的高斯属性各自回归，容易出现不连贯的闪烁和运动抖动；其二，高斯点数量在所有帧保持恒定——这意味纹理突变区域（如角色突然张嘴）无法获得足够的高斯点来表示细节，而静态区域则存在冗余。

\textbf{改了什么：Mamba 时序缓冲区 + 逐帧自适应致密化。}4DSTR 的核心改动有两方面。第一，引入基于 Mamba 的时序缓冲区（temporal buffer），将多帧高斯属性作为时序 token 序列进行相关性建模，并回归缩放和旋转的校正残差。第二，提出逐帧自适应致密化与修剪（per-frame adaptive densification and pruning），根据每帧的梯度分布动态调整该帧的高斯点数量。

\textit{直觉。}想象拍摄一部逐帧动画：如果每帧的道具都独立摆放，它们之间没有任何关联，播放时就会出现不自然的跳变。4DSTR 的时序缓冲区就像一位"时序校对员"，它不仅看当前帧，还保留最近若干帧的高斯属性作为历史记录，通过 Mamba 的选择性状态空间模型（selective state-space model）从历史中提取时序模式，对当前帧的属性进行校正。同时，自适应致密化相当于给每帧分配不同数量的画笔——运动剧烈的帧需要更多画笔来捕捉细节。

\textit{时序缓冲区与 Mamba 相关性建模。}具体而言，时序缓冲区 $M^G \in \mathbb{R}^{T \times d_G}$ 维护最近 $T$ 帧的高斯属性向量（$d_G=11$，包含位置、缩放、旋转和不透明度）。当新的帧 $t$ 的特征 $F_t$ 到来时，缓冲区以滑动窗口方式更新：丢弃最远的 $F_{t-T-1}$，将 $F_t$ 与最近 $T-1$ 帧拼接为 $M^G_{t-1} = \{F_{t-T}, \ldots, F_{t-1}, F_t\}$，送入 Mamba 模块提取时序相关性特征：

\begin{equation}
\label{eq:4dstr_mamba}
\hat{F}_t = \mathrm{Mamba}(M^G_{t-1})
\end{equation}

其中 Mamba 指标准选择性状态空间模型（2024）~\cite{gu2023mamba}。$\hat{F}_t$ 编码了当前帧与历史帧之间的时序关联信息，随后用于回归校正残差。

\textit{属性校正。}时序相关性特征 $\hat{F}_t$ 与当前帧和历史帧的属性通过动态加权融合，回归缩放和旋转的校正残差。以缩放校正为例（旋转采用相同流程）：

\begin{equation}
\label{eq:4dstr_rectify_scale}
\Delta s_t = W(\hat{F}_t \odot s_t \odot s_{t-1}), \quad \tilde{s}_t = s_t + \Delta s_t
\end{equation}

其中 $W$ 为动态加权方法（2023）~\cite{aydemir2023adapt}，$s_{t-1}$ 是最近历史帧的缩放，$\Delta s_t$ 为回归的缩放残差，$\tilde{s}_t$ 为校正后的缩放。直观理解：校正量由"时序上下文"（$\hat{F}_t$）、"当前值"（$s_t$）和"前一帧值"（$s_{t-1}$）三者共同决定，确保缩放的变化与历史趋势一致。

\textit{逐帧自适应致密化与修剪。}传统4D-GS 在所有帧使用相同的高斯点数量，4DSTR 观察到训练过程中各高斯点的累积梯度 $G(p)$ 呈对数正态分布（log-normal distribution）。为此，4DSTR 为每一帧独立计算致密化阈值：

\begin{equation}
\label{eq:4dstr_densify}
\tau_t = Q_{(1-\beta)}\!\left(\{G(p) \mid p \in \mathcal{F}_t\}\right)
\end{equation}

其中 $Q_{(1-\beta)}(\cdot)$ 表示对所有高斯点累积梯度的 $(1-\beta)$ 分位数，$\mathcal{F}_t$ 为帧 $t$ 的高斯点集合。只有 $G(p) > \tau_t$ 的点才会被致密化。修剪则基于不透明度阈值 $\tau_o$ 和世界空间缩放约束 $[s_{\min}, s_{\max}]$：

\begin{equation}
\label{eq:4dstr_prune}
\mathcal{F}'_t = \{p \in \mathcal{F}_t \mid (\alpha(p) > \tau_o) \wedge (s_{\min} < s_p^t < s_{\max})\}
\end{equation}

这确保纹理突变区域自动获得更多高斯点，而静态区域去除冗余。为防止致密化/修剪破坏帧间高斯点的对应关系（temporal correspondence），4DSTR 还设计了逐帧索引机制，显式关联每个高斯点与其所属帧。

\textit{训练目标。}4DSTR 使用多视角 SDS 损失结合集合平均损失（collective average loss, CAL）进行优化。CAL 聚合一个子片段内 $T_s$ 帧的总损失取平均：$L_{\mathrm{CAL}} = \frac{1}{T_s} \sum_{t=1}^{T_s} L_t$，弥补逐帧损失无法约束跨帧一致性的不足。

\paragraph{MVG4D: 图像矩阵监督4D生成}

与 4DSTR 的时空校正策略不同，MVG4D（2025.07）~\cite{2507_18371} 走了一条"先生成监督信号、再优化4D表示"的路线。其图像矩阵模块首先利用视频扩散模型从输入图像生成动态视频序列，再通过微调的多视角扩散模型为每帧合成多视角图像，形成兼具时序一致性和空间多样性的图像矩阵。这些图像作为密集监督信号优化3D高斯点云，最后通过轻量级形变网络（基于4D-GS的 HexPlane 形变场）扩展到时间域。为解决不同视角参数导致的视角不对齐问题，MVG4D 在训练时引入了精心对齐的监督策略，最小化生成的新视角图像与对应输入帧之间的感知差异。该方法在 CLIP-I、PSNR、FVD 等指标上均优于基线方法。

\paragraph{Gaussian Variation Field: 高斯变分场4D建模}

Gaussian Variation Field（2025.07）~\cite{2507_23785} 提出了一种高效的4D生成建模框架，其核心创新在于绕过了逐实例拟合的瓶颈。传统方法需要为每个4D动画序列单独拟合动态3DGS，耗时且难以扩展。该方法提出的 Direct 4DMesh-to-GS Variation Field VAE 直接从3D动画网格数据编码规范高斯（canonical GS）和时序变化量（Gaussian Variation Fields），将高维动画序列压缩到紧凑的512维隐空间。为此，VAE 引入了网格引导的插值机制，通过自适应半径的 K 近邻插值将网格位移场传递到高斯点，增强了运动感知能力。在此基础上训练的 Gaussian Variation Field 扩散模型（基于 DiT 架构）在合成数据上训练后，展现了对野外视频输入的出色泛化能力。

\paragraph{PhysGM: 前馈物理属性预测4D生成}

PhysGM（2025.08）~\cite{2508_13911} 是一种前馈生成框架，从单张图像同时预测3D高斯表示和物理属性（如弹性、质量等），实现即时的物理模拟和高保真4D渲染。与以往需要逐场景优化物理属性的方法（如 PhysGaussian）不同，PhysGM 通过摊销推理（amortized inference）将物理属性预测转化为单次前馈过程。该方法提出了 PhysAssets 数据集（50K+ 3D资产），并使用 Direct Preference Optimization（DPO）将模拟结果与物理合理的参考视频对齐，确保生成内容的物理合理性。

%---------------------------------------------------------------------
\subsubsection{运动迁移与动画生成}
\label{sec:08_3_1_3}

本路线介绍将特定运动模式应用到生成的3D模型上的方法，使得静态的3D资产能够按照用户意图进行动画展示，涵盖姿态驱动、交互式控制、骨骼绑定等技术。

运动迁移与动画生成方法关注如何将特定的运动模式应用到生成的3D模型上，使得静态的3D资产能够按照用户意图进行动画展示。

\paragraph{锚点方法：角色化身重建}

CharacterShot（2025.08）~\cite{2508_07409} 的核心贡献在于解决了4D角色动画生成中一个长期被忽视的问题——\textbf{可控性}：从用户提供的单张角色图像和2D姿态序列出发，生成连贯、可控的动态3D角色动画。

\textbf{Baseline 做法及其局限。}前述4D生成方法（如 Stag4D、SV4D）的典型管线是：先用扩散模型生成单视角角色视频，再基于该视频优化4D高斯表示。然而，这条管线存在两个根本缺陷。第一，从自定义角色图像生成特定运动的视频本身就极其复杂且代价高昂，普通用户难以获得满意的驱动视频。第二，即使获得了多视角视频，这些视频之间往往存在细微的跨视角不对齐（slight misalignment），导致4D高斯优化过程中出现离群点和噪声——具体表现为角色手部突然消失或可见的几何伪影，这正是逐帧独立优化缺乏几何约束的后果。

\textbf{CharacterShot 的三阶段管线。}为解决上述问题，CharacterShot 设计了从姿态控制到4D输出的端到端框架，分为三个阶段：

\textit{阶段一：姿态驱动的2D角色动画。}以 DiT（Diffusion Transformer）架构的图生视频模型 CogVideoX 为基座，将姿态信息编码为与视频隐变量同维度的条件隐变量 $z_p \in \mathbb{R}^{f \times h \times w \times 16}$，并与视频隐变量 $z_i$ 和参考图像隐变量 $z_r$ 在通道维度拼接：

\begin{equation}
\label{eq:charactershot_stage1}
z_0 = \big([z_r,\; z_i],\; [z_p',\; z_p]\big), \quad z_0 \in \mathbb{R}^{(f+1) \times h \times w \times 32}
\end{equation}

其中 $z_p'$ 是参考图像对应的姿态隐变量，$f$ 为帧数，$h \times w$ 为空间分辨率。训练时随机替换参考帧及其姿态，增强模型对姿态输入不对齐的鲁棒性。

\textit{阶段二：多视角视频生成。}将阶段一的2D输出扩展到 $V$ 个视角，引入相机编码器将 Pl\"ucker 嵌入（一种将相机内外参编码为 $6 \times 8h \times 8w$ 特征图的方法）编码为相机 token，注入每个视角的输入序列。

\textbf{双注意力模块（dual-attention module）}是此阶段的关键创新：先前方法（如 SV4D）分别计算空间注意力、时间注意力和视角注意力，这种分离式设计不利于跨维度信息的隐式传递。CharacterShot 改用两个并行的3D全注意力块——分别以形状 $V \times (f+1) \times N$ 和 $(f+1) \times V \times N$ 重排 token——同时建模时空相关性和空间-视角相关性，实现连贯的多视角视频生成。

\textit{阶段三：邻居约束的4D高斯优化。}获得多视角视频后，采用由粗到精的优化框架建模可变形3D高斯。粗阶段用 $L_1$ 损失快速构建初始4D空间；精阶段逐步精化可变形高斯。核心创新是\textbf{邻居约束}（neighbor constraint），用于解决多视角视频不对齐导致的优化伪影。直觉上，该约束要求每个3D点与相邻点之间的相对空间配置在时间上保持一致——即使点本身移动，其与邻居的距离关系应平稳变化。

具体而言，对每个3D点 $i$，计算其在时刻 $t$ 和 $t-1$ 相邻点群中心的距离：

\begin{equation}
\label{eq:charactershot_neighbor_dist}
\mathbf{L}_i^t = \mathbf{u}_i^t - \frac{1}{|N(i)|} \sum_{j \in N(i)} \mathbf{u}_j^t
\end{equation}

其中 $N(i)$ 是点 $i$ 的邻居集合，$\mathbf{u}_i^t$ 是点 $i$ 在时刻 $t$ 的3D坐标。邻居损失通过一个二值门控机制仅对出现异常位移的点对激活：

\begin{equation}
\label{eq:charactershot_neighbor_loss}
m_i = \mathbb{1}\!\big[\|\mathbf{u}_i^t - \mathbf{u}_i^{t-1}\| > \delta\big], \quad m_{ij} = m_i \cdot m_j, \quad \mathcal{L}_{\mathrm{neighbor}} = \sum_{(i,j) \in E} \big\|\mathbf{L}_i^t - \mathbf{L}_i^{t-1}\big\|^2 \cdot w_{ij} \cdot m_{ij}
\end{equation}

其中 $\delta$ 是预设的位移阈值（用于识别离群点），$w_{ij} = \|\mathbf{u}_i^{t-1} - \mathbf{u}_j^{t-1}\|$ 是空间边权重。精阶段的完整损失为：

\begin{equation}
\label{eq:charactershot_fine_loss}
\mathcal{L}_{\mathrm{fine}} = \lambda_1 \mathcal{L}_1 + \lambda_2 \mathcal{L}_{\mathrm{LPIPS}} + \lambda_3 \mathcal{L}_{\mathrm{neighbor}} + \lambda_4 \mathcal{L}_{\mathrm{TV}}
\end{equation}

\textit{白话解释}：公式~(\ref{eq:charactershot_neighbor_loss})的设计非常巧妙。$m_i$ 是一个指示函数——只有当点 $i$ 的位移超过阈值 $\delta$ 时才激活。这意味着邻居约束只惩罚那些\textbf{已经出现异常位移的离群点}，而对正常运动的点不施加额外约束，避免了过度约束导致运动僵硬。$w_{ij}$ 以相邻点间的空间距离为权重，使得距离越近的邻居对约束贡献越大，这符合局部刚性的物理直觉——一个点应该和它最近的邻居一起运动。

\textit{注解}：为支撑训练和评测，CharacterShot 构建了大规模数据集 Character4D，从 VRoidHub 平台筛选 13,115 个高质量角色模型，绑定 Mixamo 的 40 种骨骼动画（跳舞、唱歌、跳跃等），在 Blender 中渲染 21 个视角的多视角视频。渲染配置为：相机半径 2.5、视场角 $40^\circ$、A-pose 初始姿态。同时发布了评测基准 CharacterBench。

\paragraph{Drag4D: 交互式4D环境生成}

与 CharacterShot 的姿态驱动动画不同，Drag4D（2025.09）~\cite{2509_21888} 引入交互式4D环境生成框架，将用户定义的3D轨迹控制与文本驱动3D背景生成统一。该方法分三阶段工作：首先通过全景图生成和深度估计构建3D背景场景（使用2D高斯泼溅表示，并引入基于深度-法线相似度的像素级自适应加权机制减少噪声区域影响）；然后通过3D复制粘贴将参考图像中的前景物体提取为360度3D对象，并通过碰撞和重力损失实现物理感知的空间定位；最后通过部件增强的运动条件视频扩散模型，使用户可以定义3D轨迹控制物体运动方向，并与高质量3D背景协调融合。

\paragraph{DynaPose4D: 姿态关键点监督4D生成}

DynaPose4D（2025.10）~\cite{2510_22473} 将4DGS与类别无关姿态估计（CAPE）技术结合。其核心思路是用姿态关键点作为额外的监督信号：先用 Zero-1-to-3 从单张静态图像构建3D模型，再通过 Stable Video Diffusion 生成驱动视频，最后用4DGS扩展到动态4D内容。关键的姿态对齐损失通过 PoseAnything 预测的逐帧关键点与生成结果的关键点进行匹配，确保运动过渡的时序连贯性。实验表明，去除姿态监督后生成内容出现时间抖动和空间不一致，验证了姿态关键点引导的有效性。

\paragraph{TiP4GEN: 全景-透视双分支4D生成}

TiP4GEN（2025.08）~\cite{2508_12415} 面向全景动态场景生成，提出了双分支生成模型架构。全景分支以全局文本提示引导，确保场景整体连贯性；透视分支以多个局部文本提示引导，增强内容多样性和细粒度控制。两个分支通过双向交叉注意力机制实现信息交换。为克服全景视频数据的稀缺性，TiP4D 利用预训练的透视视频扩散模型作为基础，通过混合训练策略迁移其生成能力。在4D重建阶段，几何对齐重建模型对场景点云进行时空对齐，并估计相机姿态以初始化场景，生成360度沉浸式动态场景。

\paragraph{GS-DiT: 伪4D高斯场视频可控生成}

GS-DiT（2025.01）~\cite{2501_02690} 的目标是为视频生成赋予4D可控性（如多机位拍摄、推拉变焦等高级摄影效果）。其核心思路是构建伪4D高斯场作为桥梁：首先提出高效密集3D点追踪方法（D3D-PT），从输入视频中估计所有像素的3D轨迹，直接构建伪4D高斯场而无需优化；然后用该高斯场按原始相机参数渲染视频，作为条件输入微调预训练 DiT，使其学会从渲染视频恢复原始视频。推理时，只需更换相机参数即可在保持相同动态内容的同时生成新视角视频。该方法的关键创新在于松耦合的深度设计——D3D-PT 的深度估计不与2D追踪紧耦合，使其能灵活适应真实场景中不同的深度分布。

\paragraph{TextMesh4D: 文本到4D网格直接生成}

TextMesh4D（2025.06）~\cite{2506_24121} 是首个实现文本到4D网格直接生成的方法。与基于 NeRF 或3DGS的4D生成不同，网格表示在现代图形管线中具有优先地位，但其拓扑约束限制了形变的灵活性。TextMesh4D 的核心创新是 Jacobian 形变场（JDF），将形变单位从顶点提升到面片：每个三角面使用独立的 Jacobian 矩阵建模变换，打破了顶点邻居间的约束，使变形不受拓扑限制。形变后的表面通过 Poisson 积分重建，即使在高方差梯度下也能保持光滑。配合局部-全局语义正则化器（LGSR）——局部利用网格拓扑连接性稳定刚性和抑制随机噪声，全局保持帧间语义身份一致——确保帧间语义连贯。

\paragraph{DreamMesh4D: 网格-高斯混合4D生成}

DreamMesh4D（2024.10）~\cite{2410_06756} 采用网格-高斯混合表示（基于 SuGaR），将高斯泼溅绑定到网格三角面上实现联合纹理和几何优化。该方法在静态阶段生成高质量3D模型后，在动态阶段使用测地线距离构建变形图，通过稀疏控制节点驱动运动。其核心创新是自适应混合蒙皮算法，结合了线性混合蒙皮（LBS）和双四元数蒙皮（DQS）的优势：LBS 公式简洁、动画自然但存在体积丢失伪影，DQS 消除该伪影但无法处理非刚性变形。DreamMesh4D 在两者之间自适应切换，实现了更鲁棒的变形效果。该方法的输出兼容现代图形管线，在3D游戏和影视行业具有应用潜力。

\paragraph{HoloTime: 360度4D场景生成}

HoloTime（2025.04）~\cite{2504_21650} 提出了360度4D场景生成框架。该方法首先构建了 360World 数据集——首个固定机位全景视频数据集，然后训练全景动画器（Panoramic Animator），采用两阶段运动引导生成策略：先生成低分辨率粗视频提供全局运动引导，再生成高分辨率视频细化细节。为增强模型泛化能力，训练时混合使用全景视频和风景延时视频。在4D重建阶段，通过全景光流估计和窄视场深度估计进行时空深度估计，确保跨时间和空间的深度一致性，再将全景视频转化为动态点云，优化4D-GS表示，重建时空一致的4D场景。

\paragraph{In-2-4D: 4D中间帧插值}

In-2-4D（2025.04）~\cite{2504_08366} 提出了全新的4D中间帧插值任务——仅从两张单视角图像（分别捕获物体的两种运动状态）生成4D动态内容。该方法采用分治策略处理大范围、非线性运动：利用基础视频插值模型（如 DynamiCrafter）生成初始中间帧，然后通过 DIFT 特征相似度自适应识别关键帧，将复杂运动分解为时间局部、运动简单的片段。对每个片段，先构建关键帧的3DGS表示，再通过形变场将时序帧内的运动注入3DGS。最后，通过插值边界形变场并对齐引导视频，将独立生成的3D运动段合并为全局平滑的4D内容。

\paragraph{Gen4D: 全自动4D人体动画生成}

Gen4D（2025.06）~\cite{2506_05397} 提供了全自动4D人体动画生成管线，无需手动3D建模或场景设计。该方法首先从公开互联网视频中利用 SMPLerX 提取专家3D运动数据，然后通过文本提示引导的扩散模型优化规范空间的高斯头像（在 SMPLX 网格上初始化高斯点，通过双分支 SDS 损失同时优化 RGB 和深度），生成多样化的人体外观。最后通过人体感知的背景合成器（根据姿态上下文生成与人体匹配的场景背景）完成整体动画生成。该方法还发布了 SportPAL 运动场景合成数据集（50万+帧，覆盖棒球、冰球、足球三个运动项目，包含2D/3D姿态标注、SMPLX参数、分割掩码等多模态标注）。

\paragraph{OneTo3D: 文本指令驱动骨骼动画}

OneTo3D（2024.05）~\cite{2405_06547} 结合了隐式生成与显式控制的优势。该方法先用基于3DGS的模型（依赖 Zero-1-to-3）从单图生成初始3D模型，然后在 Blender 中设计自适应骨骼绑定机制——分析3D模型的几何参数信息（宽度、高度、骨骼旋转角度等），自动生成并微调适配人体骨架。核心创新在于文本指令解析算法：分析用户命令的时态和语法，提取各骨骼的位置和方向数据，支持多子命令分割和连续编辑。骨骼驱动的关键帧插入机制配合 Blender 接口实现可重复的复杂动作序列，从而实现精确的3D模型动作控制和无限时长的语义连续3D视频生成。

\paragraph{WorldPrompter: 全景视频中间表示场景生成}

WorldPrompter（2025.04）~\cite{2504_02045} 验证了360度全景视频作为3D场景生成中间表示的有效性。该方法收集了约1700个全景视频和20万+全景图像数据集，训练基于 DiT 的全景视频生成器，可生成长达10.6秒（128帧）的360度全景视频，模拟人在虚拟环境中行走的效果。关键设计包括混合训练策略（混合室内全景视频和户外全景图像增强泛化能力）和掩码扩散损失（确保生成无伪影）。生成的全景视频经过 COLMAP 估计相机姿态，再通过 Long-LRM 前馈3D重建器转化为高斯泼溅。该方法实现了COLMAP匹配率94.6\%的高质量全景重建，远超360DVD的41.8\%。

\paragraph{Track4DGen: 点追踪先验注入4D生成}

Track4DGen（2025.12）~\cite{2512_06158} 提出两阶段框架，核心创新在于将基础点追踪器（CoTracker3）的运动先验显式注入4D生成的两个阶段。第一阶段在多视角视频扩散模型中，通过分析 U-Net 各块特征的时序对应能力，发现第二时空上采样块中时间运动模块的特征最适合追踪。据此，在训练时引入密集特征级点追踪监督，强制扩散特征编码帧间对应关系，有效抑制外观漂移。第二阶段在4D-GS重建中，将携带追踪先验的扩散特征与标准 HexPlane 特征拼接，形成混合运动编码，并通过4D SH建模增强动态表示，生成时序稳定且可文本编辑的4D资产。

\paragraph{Text-to-3D with Physics-Grounded Motion: 物理驱动运动文本生成}

Text-to-3D with Physics-Grounded Motion（2024.12）~\cite{2412_05560} 从文本同时生成 3D 模型和物理驱动的运动。

\begin{takeaway}
4D 动态生成将3DGS生成扩展到时域，核心挑战是时序一致性和运动自然性。Track4DGen 将点追踪先验注入扩散模型和 4D-GS 重建两阶段，有效抑制外观漂移。技术路线从简单的 SDS+时间变形走向结构化的运动先验注入，物理驱动方法则为运动生成提供力学约束。
\end{takeaway}


\section{3D生成模型与GAN}
\label{sec:08_4}

前两节讨论的方法以单张图像或文本提示为输入，通过优化或前馈网络生成3D内容。本节转向另一类思路：直接学习3D资产的概率分布，从中采样生成全新的3D高斯表示。这类方法不依赖逐样本优化，而是训练一个生成模型（GAN、扩散模型、VAE 或自回归模型），在推理时一步到位地输出3DGS。

核心挑战在于3D高斯的离散性与无结构特性。与图像的规则网格不同，一组3D高斯是无序点集，其数量、排列顺序和属性值均不固定，传统生成模型难以直接建模。

前面的4D生成方法关注的是让3D内容"动起来"，而本节回到静态3D内容本身，但换一个角度：不是从文本/图像"推导"出3D，而是让模型"学会"3D高斯的概率分布，直接从中采样生成全新的3DGS。这种范式有望实现真正的"一步到位"生成。本节按技术路线组织如下：首先介绍3D GAN方法，将经典EG3D范式迁移到3DGS；然后讨论扩散模型如何直接在3DGS空间做生成；接着介绍自回归和Transformer-based方法；最后分析高斯原生生成方法，包括VAE和潜在空间扩散。

\subsection{生成模型范式}
\label{sec:08_4_1}

本节讨论6条技术路线：（1）3D GAN：从EG3D到高斯；（2）扩散模型直接生成3DGS；（3）潜在空间与VAE方法；（4）自回归与Transformer-based方法；（5）蒸馏与适配器方法；（6）其他3D生成模型。

\subsubsection{3D GAN：从EG3D到高斯}

本路线介绍将成熟的2D/3D对抗生成框架（如StyleGAN、EG3D）迁移到3DGS空间的方法，用高斯解码器替换NeRF渲染分支，实现高效且可编辑的3D生成。

3D GAN路线的思路是将成熟的2D/3D对抗生成框架（如StyleGAN、EG3D）迁移到3DGS空间，用高斯解码器替换NeRF渲染分支，实现高效且可编辑的3D生成。

3D感知GAN的代表工作EG3D通过三平面（tri-plane）+ NeRF体渲染实现了高质量3D生成。然而，NeRF渲染的巨大计算开销限制了在低功耗设备上的部署，且隐式表示难以直接集成到显式3D场景中。3DGS的显式性和高效渲染能力为3D GAN提供了理想的替代渲染器。

\paragraph{从 NeRF 到高斯的 3D 感知 GAN 解码器}

3D感知GAN的代表工作EG3D通过三平面+ NeRF体渲染实现了高质量3D生成，但NeRF渲染的巨大计算开销使其无法在低功耗设备（如移动端和VR/AR头显）上部署，且基于神经网络的隐式表示难以直接集成到显式3D场景（如VR环境或游戏引擎）中。3DGS的显式性和高效渲染能力恰好弥补了这两方面的不足。

Gaussian Splatting Decoder（2024.04）~\cite{2404_10625}的核心思路是：保留EG3D成熟的生成管线（StyleGAN2 backbone + 三平面特征），但用一个可训练的高斯解码器替换NeRF渲染分支，将隐式三平面映射为显式3D高斯。为降低复杂度，方法省略了球谐系数的估计，仅预测位置、颜色、旋转、缩放和不透明度五个基本属性。值得强调的是，解码器不使用任何超分辨率模块，直接以GAN最终输出的相同分辨率渲染高斯场景。这使得解码后的场景可直接导出至3D建模环境，并支持任意分辨率和宽高比的实时渲染。

位置初始化是关键难点。3D高斯具有多个属性（位置、颜色、旋转、缩放、不透明度），自由度极高，直接从梯度学习位置非常困难。当高斯位置错误时，优化器倾向于降低不透明度使其消失或调整颜色以适应当前位置，而非将其移动到正确位置。标准3DGS通过迭代式的分裂/克隆/剪枝操作来克服这一问题，但生成器必须单次前馈完成，无法使用这种迭代策略。

方法的解决方案是利用预训练3D GAN三平面中已编码的几何信息：在 $128\times128\times128$ 的均匀采样点上，用预训练NeRF MLP解码不透明度，再通过Marching Cubes估计表面，从中随机采样50万个点并略微向中心内插，形成"厚表面"作为高斯位置的初始化。消融实验表明，虽然网格采样在定量指标上略优，但在视觉上会产生水平和竖直条纹伪影。Marching Cubes采样在指标与视觉质量之间取得了最佳平衡。值得注意的是，从不透明度中采样表面的计算开销较大，但每个ID/隐向量只需执行一次，之后高斯场景即可高效渲染。

在解码器架构上，方法采用顺序解码（sequential decoding）策略。先预测颜色，再将颜色信息与三平面特征一起送入下一个解码器预测不透明度，依次迭代直到解码全部属性（颜色$\to$不透明度$\to$旋转$\to$缩放$\to$位置偏移），而非一次性预测所有属性。设计直觉在于建立属性间的依赖关系：例如缩放解码器可以从颜色和旋转信息中推断该高斯应有多大。同时，每个解码器只负责预测一部分属性，有效降低了单个解码器需要处理的自由度，使其更容易特化。每个子解码器内部使用三个隐藏层（各128个神经元，GELU激活），输出层无激活函数，唯独缩放解码器使用反Softplus激活以防止高斯过大导致光栅化时GPU显存暴涨。消融实验进一步验证了这一设计：并行解码器（同时预测所有属性）与顺序解码器性能接近，但反转顺序的解码器显著下降，说明解码顺序对性能有重要影响。

除了优化解码器权重外，方法还对预训练StyleGAN backbone进行微调。动机在于：NeRF渲染仅需查询特定位置的颜色和密度，而高斯还额外拥有缩放和旋转属性，会影响周围区域，因此三平面特征需要适配这种本质不同的表示形式。消融实验证实，关闭backbone微调后ID相似度从0.883降至0.858，说明微调确实帮助三平面特征为高斯解码提供了更好的基础。训练时损失函数由五项组成：L1损失、SSIM结构相似性损失、LPIPS感知损失、基于ArcFace的ID相似性损失（仅对正面视角计算），以及基于Sobel边缘检测的梯度损失（引导解码器关注需要精细结构的区域）。损失权重设为 $(\lambda_1, \lambda_2, \lambda_3, \lambda_4, \lambda_5) = (0.2, 0.5, 1.0, 1.0, 0.2)$。消融实验揭示了一个有趣发现：LPIPS损失对渲染清晰度至关重要。禁用后渲染结果显著模糊，而禁用L1或SSIM主要在边缘产生伪影，说明L1/SSIM负责构建粗粒度几何，LPIPS则提供精细结构的梯度信号。训练使用Adam优化器（学习率0.0009），共500k次迭代。

这条路线的优势在于：生成质量继承自成熟的3D GAN框架，而输出是可直接渲染、可编辑的3D高斯，渲染速度约为原始GAN的4倍。将渲染分辨率提高4倍后，帧率仍为原始GAN在低分辨率下的3倍以上。方法还首次将GAN反演（GAN inversion）和GAN编辑应用到显式3D高斯场景。用户在潜空间中修改隐编码，对应的3D高斯场景随之变化，支持潜在插值和风格编辑。实验在三个预训练GAN上验证：EG3D-FFHQ、EG3D-LPFF和PanoHead-FFHQH。在 $3D_{\text{FFHQ}}$ 模型上ID相似度最高达到0.968。

方法也观察到一定局限：由于省略了球谐系数，解码器无法学习视角依赖的外观（如瞳孔方向），眼睛渲染可能出现轻微模糊或始终朝前注视的现象，未来可通过引入球谐系数来缓解。此外，高斯表示还减少了EG3D/PanoHead中常见的锯齿和纹理粘连伪影。在原始NeRF渲染中旋转相机时，头发或皮肤纹理会出现位移而非随3D头部移动，而固定的高斯3D场景消除了这一问题。

\paragraph{分层多尺度高斯对抗生成}

前述方法通过替换EG3D管线中的NeRF解码器实现高斯生成，但仍保留了三平面+NeRF的中间表示。GSGAN则更进一步，直接在对抗训练框架中生成3D高斯，完全绕过NeRF中间表示。

GSGAN（2024.06）~\cite{2406_02968}直接在对抗训练框架中生成3D高斯，无需NeRF中间表示。但方法发现，朴素生成器（直接合成一组无约束的高斯）在对抗训练中面临两个问题：训练早期所有高斯从渲染图像中消失（训练不稳定），以及收敛后出现长条形高斯导致的视觉伪影。根本原因在于3DGS的非可微场景依赖特性（如SfM初始化和启发式致密化）无法直接迁移到GAN框架中。

为解决这些问题，GSGAN引入分层多尺度高斯表示。层级从粗到细编号为 $l \in \{0, \ldots, L-1\}$，不同层级的精细高斯由其对应的粗层级高斯参数化——精细高斯的位置被约束在粗高斯的局部坐标系内，且尺度随层级加深单调递减：

\begin{equation}
\boldsymbol{\mu}^{(l+1)} = \boldsymbol{\mu}^{(l)} + R^{(l)} S^{(l)} \Delta\boldsymbol{\mu}^{(l)}, \quad s^{(l+1)} < s^{(l)}
\end{equation}

其中 $\boldsymbol{\mu}^{(l)}$、$s^{(l)}$ 和 $R^{(l)}, S^{(l)}$ 分别为第 $l$ 层高斯的位置、尺度和旋转/缩放矩阵，$\Delta\boldsymbol{\mu}^{(l)}$ 是由 tanh 激活约束的局部位置偏移。尺度差使用 softplus 取负值确保单调递减。直觉上，粗层级建模场景的大尺度结构，精细层级逐步补充细节。层级间通过注意力层（底层用自注意力，高层用KNN（K-Nearest Neighbors，K近邻）局部注意力，$k=16$）建模高斯交互，KNN正则化损失鼓励邻居高斯具有一致的属性，进一步稳定训练。实验在FFHQ和AFHQ-Cat数据集（$256\times256$ 和 $512\times512$）上验证：GSGAN在FID上与Mimic3D可比，渲染速度提升约100倍，训练时间仅为Mimic3D的约44\%（28 RTX A6000天 vs. 64 A100天）。

\paragraph{单目数据集前馈 3D 感知生成}

F3D-Gaus（2025.01）~\cite{2501_06714}面向一个更困难的问题：仅从单目数据集（如ImageNet，含420K训练图像）学习3D感知生成，无需多视角或动态视频数据。多视角数据集获取成本高、需要专用设备，而单目数据集可从网络大量获取，具有更好的可扩展性。

方法提出基于像素对齐高斯（pixel-aligned Gaussian Splatting）的前馈管线：U-Net backbone 以 RGB-D 输入，预测像素对齐的高斯属性图（位置偏移、不透明度、颜色、缩放、旋转），深度图由现成的单目深度估计模型提供。核心挑战在于单目数据集仅提供正面监督，侧面视角的渲染质量无法直接验证。为此，方法引入循环聚合（cycle-aggregative）自监督策略：先从正面输入预测 $GS_0$，渲染侧面视角得到 $(I_1, D_1)$，再用该侧面输入预测 $GS_1$，将两个GS在高斯级别做互补聚合——在重叠区域取属性平均，在非重叠区域保留各自的高斯基元。重建损失 $\mathcal{L}_{recon} = \|I_0 - I_0^{GT}\|_1 + \|D_0 - D_0^{GT}\|_1$ 保证正面视角质量，聚合损失确保跨视角一致性。此外，方法利用视频扩散模型先验（CogVideoX）进行几何感知精炼，在宽视角场景中提升3D纹理质量。在ImageNet全量评估上，F3D-Gaus在FID和非平坦度分数（NFS）上均优于G3DR基线。

\subsubsection{扩散模型直接生成3DGS}

本路线介绍将扩散模型的去噪范式直接搬到3D高斯空间的方法，核心难点在于处理无序点集的排列不变性。

与GAN路线并行，另一条思路是将扩散模型的去噪范式直接搬到3D高斯空间。核心难点是处理无序点集的排列不变性——3D高斯不像图像像素那样有固定的排列顺序。

\paragraph{函数化高斯扩散}

在3DGS空间做生成面临的核心障碍是高斯的离散性与无序性：不像图像像素那样排列在规则网格上，3D高斯是一个无序点集，其数量不固定、排列顺序无意义。这意味着成熟的2D扩散模型（如 Stable Diffusion）无法直接搬用到3DGS——扩散模型需要输入具有固定维度和结构，而3D高斯点集恰恰不具备这些性质。

DiffGS（2024.10）~\cite{2410_19657}的核心洞察是绕过离散障碍：不直接对离散高斯点集做扩散，而是将3DGS分解为三个连续函数，在函数的连续表示上训练潜在扩散模型。这一思路类似于用 SDF（符号距离函数）表示离散网格。连续函数不依赖点的数量和排列顺序，天然兼容扩散模型的生成范式。具体而言，方法定义三个解耦的高斯泼溅函数：
\begin{itemize}
\item 概率函数 $\mathrm{GauPF}(\mathbf{x})$：建模空间中任意位置 $\mathbf{x}$ 存在高斯基元的概率密度。直觉上，离所有高斯中心越远的点，被高斯占据的概率越低。真值由查询点到最近高斯中心的最小距离定义，经截断函数过滤极端值后，再用连续映射将距离转换为 $[0,1]$ 区间的概率值
\item 颜色函数 $\mathrm{GauCF}(\mathbf{x})$：给定已知的高斯基元位置 $\mathbf{x}$，预测该位置的颜色属性 $c$
\item 变换函数 $\mathrm{GauTF}(\mathbf{x})$：给定已知的高斯基元位置 $\mathbf{x}$，预测旋转 $r$、缩放 $s$ 和不透明度 $o$ 等几何变换属性
\end{itemize}
三函数的解耦设计反映了不同的建模逻辑：GauPF 从3D空间中探索几何结构（``这个位置是否存在高斯基元？''），而 GauCF 和 GauTF 在已知几何的基础上预测外观和变换属性（``这个位置的高斯长什么样？''）。这种分离使得几何生成和属性预测可以独立优化。

\textbf{高斯VAE与潜在扩散。}三个连续函数虽然化解了离散性障碍，但它们仍然是高维的神经场表示，无法直接作为扩散模型的输入。DiffGS 训练一个高斯变分自编码器（VAE）来完成压缩。GS编码器将整个3DGS形状编码为一个潜向量 $z$，三平面解码器将 $z$ 解码为特征三平面 $t \in \mathbb{R}^{H \times W \times C \times 3}$（三个正交特征平面 $t_{XY}$、$t_{XZ}$、$t_{YZ}$）。对于3D空间中任意查询点 $\mathbf{x}$，将其投影到三个平面并三线性插值得到特征 $f$，再分别送入三个神经预测器 $\phi_{\mathrm{pf}}$、$\phi_{\mathrm{cf}}$、$\phi_{\mathrm{tf}}$ 预测概率、颜色和变换。VAE 的训练目标包含属性预测的 $L_1$ 损失和潜空间的 KL 散度正则化。

在训练好的 VAE 基础上，DiffGS 在潜空间上训练潜在扩散模型（LDM），采用 DALL-E 2 架构作为扩散 backbone。为支持可控生成，通过交叉注意力机制注入条件信号（文本、图像或部分3DGS），条件编码器将输入条件投影为嵌入向量后与扩散生成器融合。

\textbf{八叉树引导的高斯提取。}生成过程的最后一步是从生成的连续函数中提取离散的3D高斯，类似于 Marching Cubes 从 SDF 中提取网格。关键挑战是从 GauPF 表示的概率场中精确定位高斯中心。朴素方法（密集采样+保留高概率点）计算量大且难以准确捕捉概率峰值。DiffGS 设计了两步提取算法。

第一步是八叉树引导的渐进采样。将3D空间递归细分为八叉树，每层仅对概率值高于阈值的区域做进一步细分（``只在有高斯的地方细化搜索''），经 $L$ 层后在最高概率区域均匀采样 $N$ 个代理点作为粗略的高斯中心。第二步是将代理点位置设为可学习参数，以最大化 GauPF 概率为目标做梯度优化，使代理点精确收敛到概率峰值位置。这一设计的核心优势是 $N$ 可以设为任意数值，意味着 DiffGS 能以任意数量和密度生成高斯基元，不受训练数据中高斯数量的限制。提取出中心位置后，颜色和变换属性由 GauCF 和 GauTF 从生成的三平面中直接预测。

DiffGS 支持无条件生成和条件生成（文本/图像/部分3DGS），还可完成高斯补全（Gaussian Splatting Completion）和点云到高斯生成（Point-to-Gaussian）等任务。在 ShapeNet 数据集上，DiffGS 在 FID 和 KID 上均优于 GET3D 和 DiffTF 等基线。350K 高斯的提取仅需约 2.5 秒，兼具生成质量和效率。

\paragraph{大规模图像扩散迁移到 3DGS}

DiffSplat（2025.01）~\cite{2501_16764}的策略是直接利用大规模文本到图像扩散模型的丰富2D先验，其核心发现在于：图像扩散模型已隐式学习了3D几何估计能力（如深度、法线预测），因此可以将3D高斯属性视为"特殊风格的图像"来生成。关键设计是将3D高斯组织为多视角泼溅网格（multi-view splat grids），每个像素对齐一个高斯基元的全部属性，网格结构与2D图像兼容。为适配预训练图像扩散模型的潜在空间，方法引入"泼溅潜变量"（splat latents）：将高斯属性的多个通道通过VAE压缩到与图像潜变量相同维度的空间，仅需复制预训练VAE卷积权重的前4个通道即可完成初始化。训练流程分两步：先训练一个仅42M参数的轻量级Transformer重建模型（不到LGM/GRM的10\%），在0.1秒内从多视角图像回归高斯网格，用于在G-Objaverse的265K物体上高效构建3D数据集；然后在泼溅潜变量上微调各种预训练扩散模型（SD1.5、SDXL、PixArt-$\Sigma$、SD3），同时使用扩散损失和3D渲染损失确保跨视角一致性。整个训练仅需8 A100 GPU约3天。由于与图像扩散模型兼容，DiffSplat能够无缝迁移ControlNet、IP-Adapter等2D技术到3D生成。

\paragraph{单阶段 3D 扩散}

DiffusionGS（2024.11）~\cite{2411_14384}指出先前方法的两阶段范式（先生成多视角图像，再重建3D）存在两个根本缺陷：多视角扩散模型缺乏3D模型，视角不一致时容易崩溃；且提示视角方向受限。方法提出单阶段3D扩散模型，在每个去噪时间步直接输出像素对齐的3D高斯点云 $G_\theta = \{g_t^{(k)}(\mu_t^{(k)}, \Sigma_t^{(k)}, o_t^{(k)}, c_t^{(k)})\}$，而非中间的2D图像，天然保证了视角一致性。高斯中心位置通过 $\mu_t^{(k)} = o^{(k)} + u_t^{(k)} d^{(k)}$ 参数化，其中距离 $u_t^{(k)} = w_t^{(k)} u_{near} + (1-w_t^{(k)}) u_{far}$ 由网络预测权重控制。为同时学习物体和场景级别的3D先验，方法开发场景-物体混合训练策略，并设计参考点Pl\"ucker坐标（RPPC）编码每个射线上距世界坐标系原点最近的点，帮助模型感知深度和3D几何。在Objaverse+MVImgNet（物体）和RealEstate10K+DL3DV10K（场景）上分别训练后，DiffusionGS在ABO和GSO上取得2.20 dB的PSNR提升，场景重建上取得1.34 dB的PSNR提升，推理速度约6秒，比DMV3D快5倍以上。

\paragraph{多视角去噪与高斯重建融合}

DSplats（2024.12）~\cite{2412_09648}将多视角扩散去噪与3D高斯重建统一在单一框架中，核心架构创新是用潜在扩散模型替换传统高斯重建器的U-Net骨干。传统高斯重建器（如LGM）使用像素级U-Net，与潜在扩散模型不兼容。DSplats将骨干替换为潜在扩散模型，引入VAE编解码器在潜空间和图像空间之间映射，并在U-Net的最后上采样块中嵌入极线Transformer（epipolar transformer）预测3D高斯参数。去噪过程中，中间潜变量先映射到3D高斯表示，再渲染回潜变量，形成自监督循环——3D表示充当自然的归纳偏置，确保跨视角几何一致性。训练采用DDPM余弦噪声调度1000步，推理用DDIM加速至50步。在Google Scanned Objects数据集上，DSplats取得PSNR 20.38、SSIM 0.842的成绩，在PSNR和LPIPS上均优于所有对比方法。

此外，Sampling 3D Gaussian Scenes（2024.06）~\cite{2406_13099} 利用潜在扩散模型在秒级采样 3D 高斯场景。

\paragraph{场景级高斯扩散生成}

GGS（Generative Gaussian Splatting）（2025.03）~\cite{2503_13272}将3DGS特征场与预训练潜在视频扩散模型集成，解决了一个关键难题：视频扩散模型的压缩潜空间本身不具备3D一致性，且标准扩散模型预测噪声而非去噪结果，但3D表示应该反映去噪后的3D场景。方法的解决方案是在特征空间（而非潜空间）学习3D表示——在U-Net最后上采样块后引入极线Transformer预测3D高斯参数，高斯渲染出特征图后经2D解码器映射到潜变量。这种"特征空间3D表示"设计使模型兼容 $v$-prediction训练范式。此外，方法设计3D解码器直接从生成的特征图预测解码后的3D场景，进一步提升一致性。在RealEstate10K和ScanNet+基准上，GGS在生成3D场景的FID上比无3D表示的基线提升约20\%，且可直接生成可编辑的3D高斯场景。

\subsubsection{潜在空间与VAE方法}

直接在3D高斯空间做扩散面临数据维度高、训练成本大的问题。潜在空间方法先通过VAE或自编码器将3DGS压缩到低维潜在空间，再在潜在空间上训练生成模型，兼顾效率和质量。这类方法的关键创新在于如何将离散无序的3D高斯转化为结构化表示（如UV图、标记序列），使其能与现有2D生成模型兼容。

\paragraph{球面映射结构化高斯}

UVGS（2025.02）~\cite{2502_01846}的核心洞察是：3D高斯集合本质上是排列不变的（permutation invariant）无序点集，不同高斯基元之间没有任何结构对应关系。这使得标准神经网络难以从中提取有意义的特征，也阻碍了与2D生成模型的直接对接。如果能将这个无序集合转化为结构化的2D表示，就能直接借用成熟的图像生成架构。

为此，UVGS引入球面映射（spherical mapping）将3DGS转换为2D UV图。具体流程是：首先将3DGS物体包围在与其同心的球体内（球体半径取物体质心到最远点的欧氏距离），然后对每个高斯基元 $g_i = \{\mu_i, r_i, s_i, o_i, c_i\} \in \mathbb{R}^{14}$ 计算球面坐标 $(r_i, \theta_i, \phi_i)$，其中径向距离 $r_i = \sqrt{x_i^2 + y_i^2 + z_i^2}$，方位角 $\theta_i = \text{atan2}(y_i, x_i)$，极角 $\phi_i = \arccos(z_i / r_i)$。将方位角和极角归一化后映射为UV图像坐标，即得到 $M \times N \times 14$ 通道的UVGS图 $\mathbf{U} \in \mathbb{R}^{M \times N \times 14}$。这种变换的关键优势在于：无论输入高斯的随机排列顺序如何，始终映射到相同的UVGS表示，天然解决了排列不变性问题；同时3D空间中相邻的高斯点会映射到UV图上的邻近像素，保留了局部空间对应关系。

球面映射不可避免地遇到多对一冲突——多个高斯基元可能映射到同一UV像素。UVGS采用动态选择（dynamic selection）策略：在每个像素位置只保留沿该射线方向不透明度最高的高斯。消融实验表明，单层UVGS（即仅保留最高不透明度的一个高斯）的重建PSNR已超过30 dB，对于一般物体已足够。对于复杂物体或真实场景，UVGS支持多层UVGS堆叠（将物体置于多层同心球体内，每层取第 $K$ 高不透明度的高斯），最多可达12层，4层UVGS即可接近拟合3DGS的重建质量。

然而，14通道的异构属性直接与预训练2D模型不兼容，更关键的是，不同属性的值分布差异巨大。邻近高斯在位置和颜色上通常平滑变化，但在旋转、缩放和不透明度上变化剧烈，直接混合处理会导致梯度异常和收敛缓慢。为此，方法设计了Super UVGS映射网络。它采用三分支CNN架构分别处理位置（径向距离 $\mu$）、变换（旋转 $r$ 和缩放 $s$）和外观（不透明度 $o$ 和颜色 $c$）三类特征，各分支隐式学习处理其独特的属性分布特性。三个特征图拼接后送入中央分支（由多层卷积 + BatchNorm + ReLU 组成，最后一层使用 $\tanh$ 激活防止梯度爆炸），融合并压缩为3通道Super UVGS图像 $\mathbf{S} \in \mathbb{R}^{M \times N \times 3}$。反向映射网络按相反结构将3通道Super UVGS还原为完整的14通道UVGS。训练目标采用MSE损失和LPIPS感知损失的加权组合 $\mathcal{L}_{\text{uvgs}} = \mathcal{L}_{\text{mse}} + \lambda \cdot \mathcal{L}_{\text{UV-lpips}}$，其中LPIPS分别在位置、颜色、缩放和旋转四个属性上计算后求和，$\lambda$ 从0线性增长至10。

令人惊讶的是，Super UVGS作为结构化的3通道图像，可以被标准图像自编码器、KL-VAE和VQ-VAE无需任何微调即直接重建，3DGS存储压缩率超过99\%（Super UVGS本身已压缩89.7\%，再加VAE压缩超过99\%）。方法还构建了迄今最大规模的3DGS数据集——将Objaverse的约40万个物体通过88视角渲染并拟合为3DGS表示，再经球面映射转化为UVGS图，固定为 $512 \times 512$ 分辨率（单层可存储最多约26.2万个独立高斯）。在此基础上，直接使用Stable Diffusion的VAE将Super UVGS映射到潜空间，训练潜在扩散模型（LDM）：无条件LDM使用DDPM余弦噪声调度（前向1000步，DDIM采样20步），文本条件LDM则通过CLIP文本编码器和交叉注意力引导生成。UVGS还首次支持高斯修复（inpainting）——在UV空间中擦除部分区域后由扩散模型补全缺失的高斯基元。在ShapeNet-cars上，UVGS的无条件生成FID优于GaussianCube和EG3D，且生成结果天然具备多视角一致性（无需额外的渲染损失或SDS损失）。

\paragraph{场景级 3D 高斯标记化}

Can3Tok（2025.08）~\cite{2508_01464}是首个面向场景级3DGS的VAE。场景级生成面临独特挑战：不同场景尺度不一致（COLMAP的SfM初始化使得全局场景尺度和每个高斯基元的缩放值均不具有度量意义），高斯数量差异巨大（单个场景可达数十万基元），且3DGS数据本质上是高度非结构化的——异构特征（几何、外观、光照混杂）与点云式的不规则性并存。实验表明，直接将现有3D VAE（如PointNet自编码器、PointTransformer、L3DG等）应用于场景级3DGS，即使参数量和网络层数超过Can3Tok，在几百个场景上都无法收敛，更不具备对未见场景的泛化能力。Can3Tok从架构设计和数据处理两个层面系统性地解决了这些难题。

\textbf{架构设计。}
Can3Tok采用基于Transformer的VAE编码器-解码器结构。编码器的输入是经过傅里叶位置编码（Fourier positional encoding，最大频带 $L_B=51$）的高斯基元参数 $G \in \mathbb{R}^{N \times (2L_B+C)}$，同时为每个高斯附加以其最近体素中心坐标的傅里叶编码作为结构化锚点（体素分辨率 $V=40$），引入规整的几何先验以降低编码器学习负担。由于单个场景的高斯数量 $N$ 超过万级，直接自注意力计算代价过高，编码器采用PerceiverIO式的交叉注意力：输入经线性层映射为$key$和$value$，与一个可学习规范查询（learnable canonical query）进行交叉注意力运算。该规范查询初始化为规则体素网格坐标及其描述子 $\text{query} \in \mathbb{R}^{M \times (P+Q)}$，其中 $M$ 为规范体素数量，$P$ 为位置维度，$Q$ 为描述子维度（$Q=256$，隐层维度768），并在训练过程中持续优化。这种设计不仅大幅降低计算开销（将 $O(N^2)$ 自注意力降为 $O(NM)$ 交叉注意力），还在非结构化3DGS数据上引入了结构化的几何先验。交叉注意力输出后，经8个自注意力块进一步挖掘潜在标记之间的亲和关系。瓶颈层将编码结果投影为均值 $\mu$ 和对数方差 $\sigma^2$，通过重参数化技巧采样潜在嵌入 $z = \mu + \epsilon \cdot \exp(0.5 \cdot \sigma^2)$，潜在空间维度为 $64 \times 64 \times 4$，恰好与Stable Diffusion的潜在空间大小一致。解码器接收潜在采样 $z$，经线性层和16个自注意力块处理后，由多层感知机映射回连续3D空间的3DGS参数。由于输入和输出均在有界空间内，MLP虽参数有限但足以完成潜在到3D的映射。训练目标为L2重建损失与KL散度正则化的加权组合：$\mathcal{L} = \text{Dist}(GS_\text{rec}, GS_\text{gt}) + \beta \cdot \mathcal{L}_\text{KL}(z, \mathcal{N}(0, I))$，其中 $\beta = 10^{-6}$。

\textbf{规范3D标记化与数据处理。}
场景级3DGS表示学习的核心瓶颈在于尺度不一致。由于COLMAP估计的相机位姿和三角化点没有度量约束，不同场景的全局尺度各异，且每个高斯基元的三维缩放值也缺乏统一参照。Can3Tok提出全面的3DGS数据处理管线解决此问题。首先，\emph{归一化（Normalization）}：将高斯中心均值平移至世界坐标原点，然后将所有高斯约束到半径 $r$ 的球体内（$r = \|x + \text{translate}\|_2 \times 1.1$），同时等比例缩放高斯的缩放因子 $s$ 和关联的训练视角相机中心 $T_i$，保持相机朝向不变。这一变换确保归一化前后渲染结果像素一致，同时将所有场景统一到有界尺度范围。其次，\emph{语义感知过滤（Semantic-aware Filtering）}：场景级3DGS重建常因观测不足而包含浮点噪声，这些噪声会严重损害潜在表示的高频细节。方法使用LangSam（文本引导的SAM变体）在场景中间帧上以"最显著区域"为提示提取语义掩码，从掩码内选取一个种子高斯，再通过K-NN在3D空间中逐步扩展，直到达到预设数量 $N=40K$，从而保留语义最重要、质量最高的高斯基元子集。此外，方法还采用随机SO(3)旋转作为数据增广策略，增加训练多样性。

\textbf{实验验证。}
在DL3DV-10K数据集（6:1训练/测试划分）上，Can3Tok是唯一能在千级场景上成功收敛并泛化到未见场景的方法。所有对比方法均出现100\%失败率（失败定义为重建L2误差超过1000.0，输出完全不可辨认）：PointNet和L3DG在超过500个场景时即无法收敛；PointTransformer虽在L2误差上略优，但重建呈现拉伸和扭曲模式。Can3Tok的单次前向编解码仅需约0.06秒（8张A100训练5天），与扩散模块兼容，可实现前馈式3D生成。潜在空间分析揭示其编码同时捕获空间和语义结构：同一场景不同SO(3)旋转的潜在嵌入在t-SNE空间中自动聚拢（无需显式约束），而不同场景的嵌入彼此分离；同一场景经不同过滤/裁剪后覆盖相同语义内容时，其潜在嵌入也保持接近。消融实验表明：去除可学习查询则因内存溢出无法训练（凸显交叉注意力的计算效率优势）；去除归一化则无法泛化到测试数据（揭示模型设计本身无法弥补尺度不一致问题）；去除语义过滤则高频细节显著退化；去除体素坐标附加则局部细节层次降低。方法还展示了下游生成应用：基于Can3Tok潜在空间训练文本条件扩散UNet实现text-to-3DGS生成（FID达28.32，PointTransformer为153.76），以及基于图像回归编码器实现image-to-3DGS生成。

\paragraph{可学习密度控制高斯扩散}

DeG（Density-Sampled Gaussians）（2026.05）~\cite{2605_16355}从密度控制角度重新思考高斯生成。不同于将高斯中心固定在体素网格或数组上的做法（如GaussianCube），DeG将几何分布与基元属性解耦。

\textbf{DeG-VAE架构。}
方法的编码器对3D资产 $\mathcal{X}$ 渲染 $K$ 个视角，分别用DINOv3提取语义特征、FLUX.2 VAE提取高频纹理特征。同时从资产表面采样稠密点云，将每个点投影到多视角特征图上并取平均，得到两个互补的特征增强点集。随后通过最远点采样（FPS）选取 $M$ 个代表性中心，经Transformer集合编码器压缩为固定大小的潜标记集 $\mathcal{Z}=\{z_i \in \mathbb{R}^C\}_{i=1}^{M}$。解码器则分为密度采样和属性预测两个阶段。

\textbf{八叉树密度函数与可微密度优化。}
DeG的核心创新是将高斯中心预测建模为从可学习空间概率密度 $p(x \mid \mathcal{Z})$ 中采样的随机过程。该密度定义在 $L$ 级八叉树上，以 $(2^L)^3$ 的有效分辨率实现稀疏计算。八叉树上的联合概率被因子化为逐层的8-way条件分布：

\begin{equation}
p(x \mid \mathcal{Z}) = \prod_{l=1}^{L} p(x_{0:l} \mid x_{0:l-1}, \mathcal{Z})
\end{equation}

其中每一项由一个共享Transformer通过交叉注意力查询潜标记 $\mathcal{Z}$ 并输出8个logit。推理时通过高效祖代采样：维护活跃单元前沿，逐层评估logit、路由样本、剪枝空分支，最终从叶单元内均匀采样得到连续锚点位置 $X = \{x_i\}_{i=1}^{P}$。参数 $P$ 不由架构固定，可在推理时动态调节以权衡渲染速度与质量。

密度优化的关键挑战在于：锚点是采样的离散操作，渲染损失 $\mathcal{L}_\text{render}$ 无法直接对密度参数 $\theta$ 反向传播。DeG利用策略梯度和差异奖励（difference reward）推导出密度梯度信号——
\begin{equation}
\nabla_\theta \mathbb{E}[\mathcal{L}_\text{render}] = \mathbb{E}\left[\nabla_\theta \log p(x_j \mid \theta) \cdot \Delta_j\right]
\end{equation}

其中 $\Delta_j = \mathcal{L}(\{x_i\}_{i=1}^{P}) - \mathcal{L}(\{x_i\}_{i \neq j}^{P})$ 衡量锚点 $x_j$ 对渲染损失的边际贡献。该信号使密度在基元显著降低重建误差的位置被增强。为实现高效计算，DeG在3DGS光栅化的CUDA反向传播过程中，利用已有的透射率和累积背景色，以几乎零额外开销同时计算所有基元级的贡献，再按锚点聚合——实现了可微的"致密化与剪枝"。

VAE训练采用三阶段课程：(1) 仅用结构监督 $\mathcal{L}_\text{struct}$ 建立粗糙几何外壳（约占总训练时间6\%）；(2) 固定密度，用渲染损失训练属性解码器（可选加速阶段）；(3) 全参数端到端训练，密度梯度信号 $\mathcal{L}_\text{rdg}$ 引导密度重分配。

\textbf{局部扩展机制。}
每个锚点 $x_i$ 通过属性解码器生成 $K$ 个高斯基元，带有学习到的局部偏移：$\{g_i^k\}_{k=1}^{K} = f_\text{attr}(x_i, \mathcal{Z})$，最终产生 $N = P \times K$ 个可渲染高斯。这种层级结构（全局密度采样+局部扩展）使模型在均匀区域使用少量锚点，在复杂细节处密集填充。

\textbf{VecSeq机制。}
为解决无序集合上扩散训练的排列歧义问题（$M!$ 种可能的噪声-数据配对导致收敛缓慢、生成模糊），DeG提出VecSeq规范序列化方法。核心思想是将无序潜标记对齐到一个固定的、确定性的空间结构：3D Sobol序列 $S = \{s_j\}_{j=1}^{M}$——一种低差异准随机序列，对单位立方体 $[0,1]^3$ 的覆盖比标准随机采样更均匀。训练时，通过最优传输计算资产特定的FPS点到Sobol锚点的最优匹配 $\pi^*$，据此重排潜标记，使第 $j$ 个标记始终关联到第 $j$ 个Sobol锚点 $s_j$ 的空间邻域。该匹配作为离线预处理一次完成，训练和推理时零额外开销。重排后的标记序列附加Sobol位置的正弦位置编码后送入Flow Matching扩散Transformer（S3-DiT骨干），将困难的集合生成问题转化为稳定的序列生成任务。

\textbf{实验结果。}
在Toys4K数据集上，DeG-VAE在相同高斯预算下于PSNR、SSIM和LPIPS三项指标上均显著优于TRELLIS和UniLat3D。值得注意的是，DeG达到与TRELLIS相同的视觉质量（LPIPS）仅需不到一半的高斯数量。在单图条件3D生成任务上，DeG取得最高的图像条件对齐分数（CLIP-I）和多数分布度量（FID$_\text{incep}$、FID$_\text{dinov2}$等），在用户研究中也以显著优势获偏好。消融实验证实：密度梯度 $\mathcal{L}_\text{rdg}$ 在低预算区间带来最大增益（自适应分配在容量受限时最有价值），VecSeq重排序相比VecSet基线显著提升条件对齐和分布质量。模型在Objaverse/XL上训练（VAE约10天+扩散约11天，32×A800），推理时可通过调节采样锚点数 $P$ 在33K至262K高斯之间灵活权衡质量与效率。

\paragraph{可拖拽几何一致性生成}

Dragen3D（2025.02）~\cite{2502_16475}引入了可控3D高斯生成的能力。方法的Anchor-GS VAE将点云和单张图像编码为锚点潜变量，再解码为3D高斯。关键创新是种子点驱动策略（Seed-Point-Driven strategy）：先生成稀疏种子点作为粗几何表示，再通过整流流模型（Rectified Flow）将种子点映射到锚点潜变量——流匹配目标 $\mathcal{L}(\theta) = \mathbb{E}_{t, x_0, x_1, y}\|v_\theta(x_t, t, y) - (x_1 - x_0)\|^2$ 学习从种子点潜空间到锚点潜空间的映射。由于种子点稀疏且易于学习，几何一致性得到保证。用户可以直观地拖拽种子点来变形最终3D高斯的几何形状，变化通过锚点潜变量传播到生成结果。这是首个无需2D扩散先验即可实现几何可控3D高斯生成和编辑的方法。

\paragraph{教师引导 2D 监督 3D 扩散}

A Lesson in Splats（2024.12）~\cite{2412_00623}解决了纯2D监督训练3D扩散模型的难题。标准训练范式要求去噪过程和监督信号在同一模态中，但大规模3D数据集稀缺。方法的关键洞察是解耦带噪样本与监督信号的来源：将带噪样本保留在3D高斯空间，而将监督信号解耦到2D图像空间。具体地，一个预训练的确定性"嘈杂教师"模型从单张图像重建3D高斯，其输出虽然不够精确（模糊、缺乏多样性），但经过足够噪声步 $t^*$ 后，教师模型产生的带噪样本 $s_t = \alpha_t s_0^\theta + \sqrt{1-\alpha_t}\epsilon$ 可以与从真实3D数据出发的带噪样本分布重叠。因此，可以用教师模型的输出替代不存在的3D真值来启动扩散训练，而监督仍使用2D多视角图像。训练分两阶段：先引导阶段（bootstrapping）用教师输出训练去噪器，再用多步去噪的渲染损失精调。在物体和场景级数据集上的验证表明，该方法始终优于确定性教师模型。

Predicting 3D Structure（2026.05）~\cite{2605_10830} 通过潜在后验采样预测 3D 结构。

\subsubsection{自回归与Transformer-based方法}

本路线介绍利用自回归模型和Transformer架构直接生成3D高斯表示的方法，需要解决高斯的排列顺序和层级结构问题。

自回归模型通过逐步预测下一个标记来生成内容。将这种范式扩展到3D，需要解决高斯的排列顺序和层级结构问题。

\paragraph{多尺度自回归高斯生成}

自回归模型在2D图像上已取得突破性进展（如下一尺度预测范式），将其扩展到3D高斯生成面临两个核心难题：(1) 高斯是无序的3D基元，不存在自然的标记序列；(2) 不同物体的高斯数量差异悬殊，难以统一建模。ARGS（Auto-Regressive Gaussian Splatting）（2026.04）~\cite{2604_00494}的解决思路是：先将高斯集简化为一棵层次化二叉树，再反转该过程作为生成路径，从而将无序高斯转化为有序的粗到细生成序列。

\textbf{基于矩守恒的简化原则。}方法的起点是高斯的矩性质——零阶矩 $m_0$ 对应总质量（不透明度与协方差行列式的函数），一阶矩 $m_1$ 指向高斯中心。ARGS的简化原则是：在每一步中，选取协方差矩阵行列式最小的高斯（即空间最紧凑、对渲染贡献最小的那个），将其与加权距离最近的邻居高斯融合。之所以选择行列式最小的而非最远的，是因为行列式直接度量高斯的体积大小，体积越小则投影面积越小、视觉影响越低，从而保证简化对渲染质量的损失最小。融合操作本身并不简单叠加属性——为正确处理两个高斯的空间重叠，ARGS引入了一组辅助"交叉高斯"：其不透明度 $\alpha_c = \alpha_1 \alpha_2$，协方差 $\Sigma_c = (\Sigma_1^{-1} + \Sigma_2^{-1})^{-1}$，均值 $\mathbf{u}_c = \Sigma_c(\Sigma_1^{-1}\mathbf{u}_1 + \Sigma_2^{-1}\mathbf{u}_2)$。融合后的零阶和一阶矩分别从两个源高斯减去交叉高斯的贡献（$m_{0,3} = m_{0,1} + m_{0,2} - m_{0,c}$，$m_{1,3} = m_{1,1} + m_{1,2} - m_{1,c}$），新均值由矩比值恢复 $\mathbf{u}_3 = m_{1,3}/m_{0,3}$，协方差和球谐特征则以零阶矩为权重做加权平均。整个简化过程用堆优先队列驱动，复杂度接近线性。

\textbf{从线性序列到层次化树结构。}直接拟合简化序列意味着自回归需要 $n-1$ 个串行步骤（每步生成一个高斯），对于动辄数万高斯的物体完全不可行。ARGS的关键改进是将线性简化序列重组为层次化二叉树：以最终融合后的单一高斯为根节点，每次融合操作的逆过程即为一个父节点分裂为两个子节点，递归展开直到所有叶节点均为原始高斯。由此，$n$ 个高斯被组织为 $\mathcal{O}(\log n)$ 层的树，每层包含该尺度下的全部高斯节点。生成时只需 $\mathcal{O}(\log n)$ 步即可完成：第 $t$ 步预测第 $t$ 层全部节点的分裂结果，而非逐点生成。

\textbf{树基Transformer与层次注意力。}ARGS采用解码器-only Transformer作为自回归骨干（24层、1536维、24头），沿树的层级由粗到细逐步预测。所有连续高斯属性（位置、不透明度、缩放、旋转、球谐系数）经对数量化后离散化为256类，每类通过可学习查找表嵌入，拼接后投影为统一标记。位置编码采用三维旋转位置编码（3D RoPE），将3D坐标分解为 $(x, y, z)$ 三组独立的旋转编码，保持空间等变性。注意力机制的设计经历了三级演进：(1) 标准因果注意力——严格逐标记生成，$O(n)$ 步，最慢；(2) 层次空间注意力——同层内并行自注意力，层间因果依赖，降至 $O(\log n)$ 步；(3) 树注意力——在层次空间注意力的基础上，将前序层的内部（非叶）节点也纳入当前层的注意力计算，使叶节点能直接访问祖先信息，在保持 $O(\log n)$ 步的同时增强长程依赖建模。模型以标准交叉熵损失训练，同时预测每个节点是否可分裂（布尔值）以及分裂后两个子节点的全部量化属性。

\textbf{噪声增强训练与局部坐标表示。}教师强制（teacher-forcing）训练下，模型只见过完美真值前驱，一旦推理时某步预测出错，误差会沿层级逐级放大（粗层错误污染所有细层）。为此，ARGS引入噪声增强训练：以比例 $\rho \in (0,1)$ 随机扰动一批高斯的属性，但监督信号保持真值不变，迫使模型学会检测和修正输入中的错误，从而在推理时具备自纠错能力。此外，ARGS还将高斯属性转换到父节点定义的局部坐标系中，使子节点的属性分布更紧凑、更容易学习。

\textbf{实验与应用。}在ShapeSplat和ModelSplat数据集上，ARGS的简化方法在多级压缩下保持高渲染保真度（PSNR下降可控）。单类（飞机类别）生成实验中，模型的可分裂检测准确率达99.60\%，高斯属性预测准确率91.45\%。层级化结构还带来多项应用价值：场景压缩与渐进传输（类可伸缩图像编解码器）、类LLM的灵活生成长度控制（任意截断点均可解码为完整场景）、以及多粒度可编辑生成（粗层控制全局形状姿态，细层调整纹理和局部几何）。当前局限包括：仅支持无条件类别级生成（尚无条件生成）、单高斯单标记的表示效率较低、以及简化过程的 $O(n^2)$ 复杂度在大规模场景下仍需优化。

\paragraph{状态空间序列建模统一 3D 生成}

MVGamba（2024.06）~\cite{2406_06367}发现现有高斯重建模型的Transformer架构虽强但计算昂贵——其二次复杂度在建模长高斯序列时成为瓶颈，常需牺牲多视角信息传播来控制开销。方法提出基于状态空间模型（SSM）的多视角高斯重建器，核心创新是以RNN风格传播因果上下文实现跨视角自精炼，同时以线性复杂度生成长高斯序列建模细节。图像经Pl\"ucker射线嵌入和卷积标记化后，由SSM处理器沿展平的多视角标记序列处理，再经轻量级高斯解码器预测每个高斯的属性。配合现成多视角扩散模型（MVDream、ImageDream），MVGamba统一了单图、稀疏图和文本驱动的3D生成任务，模型大小仅为对比方法的约0.1倍。

\paragraph{潜在标记空间建模}

LTM3D（2025.05）~\cite{2505_24245}将扩散模型和自回归模型的优势整合到一个框架中，用于条件3D形状生成。方法的核心设计包括三个组件：(1) 前缀学习（Prefix Learning）将条件标记与形状潜在标记对齐——不直接将图像/文本标记作为条件，而是学习一组可训练的前缀标记，通过交叉注意力融合条件信息，桥接模态鸿沟；(2) 条件分布建模骨干网络利用掩码自编码器和扩散模型增强标记依赖学习，使模型能灵活选择自回归或扩散生成策略；(3) 潜在标记重建模块结合重建引导采样降低不确定性，从随机初始化的标记中收敛到更准确的形状。LTM3D在标记空间操作，通过替换编解码器即可支持SDF、点云、网格和3DGS等多种3D表示，在图像和文本条件生成任务上均优于现有方法。

UniGS（2025.02）~\cite{2502_17860} 提出统一的语言-图像-3D 预训练框架，利用 3DGS 作为 3D 表示桥接多模态学习。核心思想是将每个高斯基元的全部属性（位置、颜色、缩放、旋转、不透明度）拼接为高维特征 $g = (\mu, c, \alpha, s, R)$，通过跨模态对比学习在语言、图像和3D三个模态间建立共享嵌入空间。

\textbf{高斯感知引导（Gaussian-Aware Guidance）。}UniGS的关键创新是利用3DGS渲染过程中的固有信息指导表征学习：在渲染管线中，高斯基元经投影和$\alpha$-混合产生像素级颜色，这一可微过程天然将3D几何与2D图像关联。UniGS通过高斯感知引导模块将渲染梯度反向传播到3D表征，使模型不仅学习到高斯属性的统计分布，还能理解属性之间的几何语义关系（如缩放和旋转如何共同影响渲染结果）。

\textbf{数据规模与泛化性。}UniGS通过集成现有3D数据集（ShapeNet、Objaverse等）构建了大规模3DGS训练语料，支持从10000个3D点的密集采样。实验表明，随着训练数据规模和采样点数量增加，UniGS在零样本分类和检索任务上持续提升，且预训练表征可迁移至2DGS等其他3DGS衍生方法，展现出跨表示的鲁棒性。

\subsubsection{蒸馏与适配器方法}

本路线介绍利用已训练好的2D模型（尤其是多视角扩散模型），通过蒸馏或适配器将其知识迁移到3D高斯生成的高效路线。

与其从零训练3D生成模型，另一条高效路线是利用已训练好的2D模型（尤其是多视角扩散模型），通过蒸馏或适配器将其知识迁移到3D高斯生成。

\paragraph{多视角扩散蒸馏为 3D 生成器}

DD3G（2025.04）~\cite{2504_00457}的核心思路是将预训练多视角扩散模型（MV-DM）的视觉知识和空间几何信息压缩到3D高斯生成器PEPD中，实现单张图像到3D高斯的快速转换。与SDS（Score Distillation Sampling）利用随机采样梯度不同，DD3G选择模拟多视角扩散模型的确定性ODE轨迹进行蒸馏——利用DDIM的确定性反向采样特性，离线收集大量$(I, C, II, OI)$四元组（初始图像、相机姿态、噪声图像、ODE输出图像）作为训练数据，从而避免了SDS中噪声时间步调优的敏感性。

PEPD生成器采用两阶段设计。第一阶段为\textbf{模式提取}（Pattern Extraction, PE）：鉴于MV-DM的反向轨迹受相机参数影响，该方法首先用Plücker射线嵌入将相机参数密集编码并与噪声拼接为随机变量$NC$（包含$n$个视角对应的子随机变量）。在图像到3D的提升过程中，不可见区域必须从可见区域推断——不同类型物体的推断模式各异（如对称物体可直接复制可见区域，不规则物体则不能过度依赖），因此PE阶段在输入图像引导下对$NC$进行初始映射：通过交叉注意力（Cross Attention）层注入图像信息，通过3D自注意力（3D Self-Attention）层实现$n$个视角子变量之间的交互，中间堆叠ResBlock和SpatialTransformer3D层提取提升模式，最终输出模式标记。

第二阶段为\textbf{渐进解码}（Progressive Decoding, PD）：3D高斯的各属性（位置$\mu_i$、缩放$s_i$、旋转$q_i$、颜色$c_i$、不透明度$\alpha_i$）相互耦合，用2D图像监督时多种梯度方向可能相互干扰。此前工作采用双分支结构将位置预测独立为点云生成任务，但降低了计算效率且阻断了位置与其他属性的交互。DD3G设计了单分支网络，以PointTransformer块为骨干，通过渐进策略从不同网络层级解码不同属性——模仿人类构建3D物体的方式，先解码$\mu_i$获得基本形状，再解码$s_i$和$q_i$捕获轮廓，最后解码$c_i$和$\alpha_i$获取细节信息。这种设计既为不同属性提供了相对独立的优化路径，又允许充分交互。此外，每个输出token解码多个3D高斯以提升生成细节，且不对3D高斯施加结构约束以保持其表征能力。

在训练目标上，DD3G设计了\textbf{显式监督}与\textbf{隐式验证}的联合优化框架。

\textbf{显式监督}通过模拟MV-DM的ODE轨迹实现知识蒸馏：将PEPD生成的3D高斯可微渲染为多视角图像，与ODE输出图像计算MSE和LPIPS损失。然而仅靠显式监督会导致几何不一致性——在扩散加速领域已有研究表明，用距离度量（MSE、LPIPS等）迁移分布会造成知识损失，而对抗目标可有效缓解。但在3D蒸馏场景中，数据有限且非结构化表征复杂，训练鲁棒的3D判别器极为困难。

\textbf{隐式验证}巧妙地利用MV-DM本身作为验证器：当噪声强度控制在较低水平时，SDS能在特征空间中隐式验证样本是否符合真实分布（概念上类似判别器损失），且不施加特定形式约束。具体实现为多视角SDS损失$L_{IV}$，噪声时间步$t$限制在较小范围内以确保不破坏生成样本的内部一致性。

由于训练早期生成质量较差、隐式验证难以提供有效梯度，DD3G引入\textbf{课程学习}策略：显式监督保持不变，隐式验证的权重$\beta_i$随训练进度从0逐渐增大至预设上限（$\beta_i = \beta \cdot \min(1, 0.5(e^{i/I \cdot \gamma} - 1))$，其中$i$和$I$分别为当前和总迭代次数）。此外，训练中还引入随机背景颜色（防止模型生成大量与固定白色背景匹配的冗余高斯）和不透明度正则化（防止生成过多透明高斯）。

数据方面，DD3G收集了120k高质量RGBA图像（80k合成图像 + 40k Objaverse渲染图像），经人工验证确保质量，涵盖日常物体、植物、动物及虚拟角色等多种类别。推理时仅需0.06秒即可将单张图像转换为3D高斯——用户随机选取相机姿态与高斯噪声拼接后，PEPD即可输出完整的3D高斯物体。

\paragraph{2D 扩散模型直接生成 3D 高斯}

Zero-1-to-G（2025.01）~\cite{2501_05427}的关键洞察是：高斯作为3D表示可以分解为编码不同属性（RGB颜色、缩放、旋转、不透明度、位置）的多视角属性图像（Splatter Images），每张属性图像都在预训练2D扩散模型的分布范围内。因此可以将3D生成任务重新框定为2D扩散框架内的多图像生成问题。具体实现上，方法将每个3D物体表示为 $N$ 个视角下的Splatter Images，通过变换将14通道高斯属性拆分为5张RGB兼容的属性图像。引入跨视角和跨属性注意力层，捕获复杂关联并强制生成高斯间的3D一致性。此外，方法微调LGM的重建模块生成更平滑的高斯泼溅网格作为训练数据（而非逐物体拟合的高频泼溅），因为后者经VAE编解码后会产生严重的外观变化。这是首个有效利用预训练2D扩散先验的直接图像到3D生成模型。

\textbf{Splatter Image分解与属性映射。}
每个3D物体的高斯泼溅可以被重新排列为 $H \times W$ 的规则网格（即Splatter Image），每个像素包含14个通道编码5种高斯属性：RGB颜色、缩放（scale）、旋转（rotation，四元数）、不透明度（opacity）和位置（position）。Zero-1-to-G通过精心设计的变换函数，将这14通道拆分为5张3通道属性图像，使每张图像都落入自然图像的分布范围内。这种分解使得可以直接利用预训练Stable Diffusion的VAE将属性图像编码到潜空间，从而在2D潜扩散框架中进行训练和推理。训练数据方面，方法在G-buffer Objaverse数据集（约262,000个物体）上进行训练，每个物体渲染38个视角，采用6个固定视角（输入视角+方位角30°、90°、180°、270°、330°）生成Splatter Images。

\textbf{联合多视角-多属性扩散建模。}
方法将3D物体的分布建模为多个视角和属性的联合分布。给定固定视角集合 $\{\pi_1, \pi_2, \ldots, \pi_K\}$ 和条件输入图像 $y$，联合分布为：

\begin{equation}
p(\mathbf{z}) = p\left(\mathbf{z}^{(1:K, 1:N)} \mid y\right) = p_{\text{pos, op, sc, rot, rgb}}\left(\mathbf{z}_{\text{pos}}^{1:K}, \mathbf{z}_{\text{op}}^{1:K}, \mathbf{z}_{\text{sc}}^{1:K}, \mathbf{z}_{\text{rot}}^{1:K}, \mathbf{z}_{\text{rgb}}^{1:K} \mid y\right)
\end{equation}

其中 $K$ 为视角数，$N=5$ 为属性数。为保持生成高斯间的3D一致性，方法在预训练扩散U-Net中插入两类额外的注意力层：（1）\textbf{跨视角注意力}——将同一属性下不同视角的token拼接后执行自注意力，实现视角间的信息交换，学习一致的视角分布；（2）\textbf{跨属性注意力}——将同一视角下所有属性的token组合后执行自注意力，建模属性间的相关性。训练时在每个注意力块中交替应用跨视角注意力和跨属性注意力，并参照Wonder3D使用属性切换器（attribute switcher）指定当前生成的属性类别。去噪损失函数为：

\begin{equation}
\mathcal{L} = \mathbb{E}_{t, \mathbf{z}_0^{(1:K, 1:N)}, k, n, \boldsymbol{\epsilon}^{(1:K, 1:N)}}\left[\left\|\boldsymbol{\epsilon}^{(k,n)} - \hat{\boldsymbol{\epsilon}}^{(k,n)}\left(\mathbf{z}_t^{(1:K, 1:N)}, t\right)\right\|^2\right]
\end{equation}

其中 $\boldsymbol{\epsilon}^{(k,n)}$ 是添加到第 $k$ 个视角、第 $n$ 个属性的噪声，$\hat{\boldsymbol{\epsilon}}^{(k,n)}$ 为模型对应预测。训练遵循两阶段策略：第一阶段仅训练跨视角注意力（4卡L40，13k迭代，约1天），第二阶段加入跨属性注意力层并联合微调（8卡L40，30k迭代，约2天）。

\textbf{VAE解码器微调。}
预训练Stable Diffusion的VAE在自然图像上训练，直接用于Splatter Image的编解码虽能产生视觉上可接受的属性图像，但渲染后会出现明显伪影。原因在于：（1）Splatter Image中每个像素对应一个高斯泼溅，像素值的微小变化会显著影响最终渲染；（2）Splatter Image包含VAE难以精确恢复的高频细节。为此，方法冻结VAE编码器，仅微调解码器，使用包含四项损失的渲染目标函数：

\begin{equation}
\mathcal{L}_{\text{decoder}} = \mathcal{L}_{\text{splatter}} + \mathcal{L}_{\text{normal}} + \mathcal{L}_{\text{rgb}} + \mathcal{L}_{\text{mask}}
\end{equation}

其中 $\mathcal{L}_{\text{splatter}}$ 为Splatter Image本身的重构损失，$\mathcal{L}_{\text{normal}}$ 为渲染法向量的余弦相似度损失，$\mathcal{L}_{\text{rgb}} = \mathcal{L}_{\text{MSE}} + \mathcal{L}_{\text{LPIPS}}$ 为渲染图像的MSE与感知损失之和，$\mathcal{L}_{\text{mask}}$ 为渲染掩膜的二元交叉熵损失。解码器微调在8卡L40上训练20k迭代。

\textbf{实验表现。}
在GSO数据集的30个物体评测集上，Zero-1-to-G在PSNR、SSIM、LPIPS和Chamfer Distance四项指标上均优于所有基线方法（包括TriplaneGaussian、TripoSR、LN3Diff、InstantMesh和LGM）。消融实验验证了三个关键设计的必要性：（1）不微调VAE解码器时，渲染质量显著下降；（2）移除跨属性注意力后，生成的高斯出现浮点伪影（floaters）和模糊纹理，因为不同属性无法正确对齐；（3）不使用预训练扩散先验（随机初始化UNet）时，模型在相同数据和训练迭代下无法收敛到有意义的结果。推理时，单个物体在单卡L40上仅需8.7秒（CFG scale $\nu=3.5$），总训练仅需3天/8卡L40，显著低于从头训练的直接3D生成方法和依赖多视角重建的两阶段方法。在真实野外图像（in-the-wild）测试中，Zero-1-to-G展现出良好的泛化性：两阶段方法（LGM、InstantMesh）常因上游多视角生成的不一致产生伪影纹理和几何缺陷，LN3Diff受限于训练资源产生过度平滑的几何，而Zero-1-to-G凭借预训练2D先验和单阶段直接生成，能够输出高保真、几何一致的3D渲染结果。

\subsubsection{其他3D生成模型}

本小节汇总不属于上述主要类别但在3DGS生成中具有独特贡献的方法。

\paragraph{视频扩散模型自蒸馏}

Lyra（2025.09）~\cite{2509_19296}提出从视频扩散模型中蒸馏隐式3D知识到显式3DGS表示，消除了对多视角训练数据的需求。方法基于GEN3C视频扩散模型构建师生框架：教师模型（视频扩散模型）生成相机控制的RGB视频，学生模型（3DGS解码器）在同一潜空间中操作，由教师的RGB输出进行监督。3DGS解码器将视频VAE的潜变量解码为3D高斯参数，整个训练仅需视频扩散模型生成的合成数据，无需任何真实3D或多视角标注。推理时，模型可从文本提示或单张图像合成3D场景并实时渲染。框架还扩展到从单目输入视频生成动态3D场景，展示时间维度的生成能力。

\textbf{自蒸馏架构。}
Lyra的核心思路是利用相机控制视频扩散模型 $V$ 的3D一致性来监督3DGS解码器 $D_s$，形成自蒸馏循环。教师模型 $V$ 基于GEN3C构建：它首先从输入图像构建时空3D缓存（将RGB图像通过深度估计反投影为彩色点云），然后将缓存从任意相机姿态渲染为RGB图像和遮挡掩码，作为结构化视觉引导。视频VAE将RGB视频 $x \in \mathbb{R}^{3 \times H \times W}$ 编码到潜空间 $z = \mathcal{E}(x) \in \mathbb{R}^{\tau \times C \times h \times w}$，其中潜通道维度 $C=16$，时间压缩因子 $\kappa_t=8$，空间压缩因子 $\kappa_s=8$。

给定输入图像 $I$ 和采样的相机轨迹 $\{Camera^{v}\}_{v=1}^{V}$，教师模型生成去噪后的视频潜变量 $z^v = V(I, \{Camera^{v}\}_{v=1}^{V})$。潜变量沿两条分支解码：预训练RGB解码器 $\mathcal{D}_r$ 生成RGB帧 $x^v_r = \mathcal{D}_r(z^v)$，3DGS解码器 $D_s$ 输出显式3D高斯 $\mathcal{G}$。学生模型的渲染视角 $\hat{x}^v_r = \text{Render}(\mathcal{G}, \{Camera^{v}\}_{v=1}^{V})$ 被监督以匹配教师的RGB帧 $x^v_r$，形成自蒸馏闭环。为扩大视角覆盖范围，每个输入图像采样 $V=6$ 条相机轨迹，每条轨迹包含 $L=121$ 个相机姿态，总计726个输入视角。

\textbf{基于潜空间的3DGS解码器。}
此前的前馈重建方法（如GS-LRM、AnySplat）在像素空间操作，受限于GPU内存仅能处理少量低分辨率视角。Lyra直接在视频潜空间操作，突破了这一瓶颈：3DGS解码器接收多视角视频潜变量 $z \in \mathbb{R}^{V \times L \times C \times h \times w}$ 和编码后的Pl\"ucker嵌入 $E$，通过空间 $2 \times 2$ 的patchification层转换为token，送入重建模块。重建模块采用Long-LRM的架构：每个block由1个Transformer层和7个Mamba-2层组成，重复两次共16层，隐层维度512。最后通过转置3D卷积输出14通道的高斯参数：3D位置 $(x,y,z)$、缩放 $(s_x,s_y,s_z)$、旋转四元数 $(q_w,q_x,q_y,q_z)$、不透明度 $\alpha$ 和RGB颜色 $(r,g,b)$。解码器形式化为：

\begin{equation}
    \mathcal{G} = D_s(z, E)
\end{equation}

其中Pl\"ucker嵌入 $E$ 由相机姿态和内参计算原始6维嵌入，利用预训练VAE编码器 $\mathcal{E}$ 分别编码光线方向和光线方向与原点叉积，拼接后通过轻量MLP降维至与潜变量 $z$ 相同的维度。

\textbf{损失函数。}
总损失函数为各损失项的加权和：

\begin{equation}
    \mathcal{L} = \lambda_{\text{mse}} \mathcal{L}_{\text{mse}} + \lambda_{\text{lpips}} \mathcal{L}_{\text{lpips}} + \lambda_{\text{depth}} \mathcal{L}_{\text{depth}} + \lambda_{\text{opacity}} \mathcal{L}_{\text{opacity}}
\end{equation}

其中 $\lambda_{\text{mse}}=1.0$、$\lambda_{\text{lpips}}=0.5$、$\lambda_{\text{depth}}=0.05$、$\lambda_{\text{opacity}}=0.1$。MSE和LPIPS损失构成图像重建监督；深度损失利用ViPE估计的视频深度进行尺度不变深度监督，防止解码器产生扁平几何；不透明度L1正则化将最低80\%不透明度的高斯剪枝，使表示更紧凑（剪枝后渲染704$\times$1280图像仅需18ms，加速1.67倍）。消融实验表明：仅用自蒸馏数据训练的模型甚至优于自蒸馏+真实数据联合训练，验证了视频扩散模型生成的合成数据具有足够的多样性和一致性。

\textbf{动态3D场景扩展。}
Lyra将静态3DGS解码器扩展为动态解码器 $D_d$，通过引入源时间嵌入 $T^{src}$ 和目标时间嵌入 $T^{tgt}$ 生成时间相关的3D高斯。时间嵌入通过2维正弦编码扩展通道后用VAE编码器编码至潜空间维度，直接加到潜变量和Pl\"ucker嵌入之和上。动态解码器形式化为 $G_t = D_d(z, E, T^{src}, T^{tgt})$，从预训练的静态解码器 $D_s$ 微调初始化。为解决动态场景中监督信号时间维度不平衡问题（早期帧靠近输入视角，后期帧远离），引入反向视频增强策略：将输入视频帧序反转后送入视频模型，生成6条额外的内缩轨迹，与原始6条外扩轨迹组合为每时间步12条监督视角。

\textbf{实验结果。}
Lyra构建了完全由视频扩散模型生成的合成数据集Lyra-Dataset：3D重建部分使用59,031张图像（354,186条视频），4D部分使用7,378个视频（44,268条视频），涵盖室内外环境、人物、动物及写实与想象内容。在RealEstate10K、DL3DV和Tanks and Temples三个基准上，Lyra在所有指标的PSNR、SSIM和LPIPS上均超过ZeroNVS、ViewCrafter、Wonderland和Bolt3D等先前方法。架构方面，联合Transformer/Mamba-2设计比纯Transformer快6.5倍（3213ms vs.\ 20922ms），而潜空间操作使其能处理726个视角，远超像素空间方法的12帧上限。

\paragraph{统一生成与编辑多视角整流流}

SplatFlow（2024.11）~\cite{2411_16443}提供了统一的3DGS生成和编辑框架。多视角整流流（Rectified Flow）模型在潜在空间同时生成多视角图像、深度和相机姿态，条件为文本提示。架构上共享Stable Diffusion 3的编码器，使整流流模型与SD3共用潜空间，充分利用其预训练知识。GSDecoder在SD3解码器基础上修改输入/输出通道和跨视角注意力，将潜在输出前馈转换为像素对齐的3DGS。通过训练无关的反演（基于SDEdit）和修复（基于RePaint）技术，SplatFlow在统一框架内支持对象编辑、新视角合成和相机姿态估计等多种3D任务。

\textbf{多视角整流流的联合建模。}SplatFlow的核心设计选择是将多视角图像、深度图和相机位姿建模为联合分布 $p(\{I_i\}_{i=1}^K, \{D_i\}_{i=1}^K, \{P_i\}_{i=1}^K \mid C)$，而非分别独立生成。这一联合建模策略有两方面动机：一是允许将多种3D任务（位姿估计、新视角合成等）统一表述为修复（inpainting）问题，只需将已知部分作为约束、未知部分作为待修复区域即可；二是真实世界场景需要自适应相机轨迹而非固定模板。具体实现上，多视角RF模型将图像潜变量 $E(I_i)$、深度潜变量 $E(D_i)$ 和Pl\"ucker射线坐标 $R_i$ 沿通道维度拼接为单个张量。每个Pl\"ucker射线表示为 $(d_i, m_i)$，其中方向向量 $d_i = R_i^\top K_i^{-1} x_i$，矩向量 $m_i = (-[d_i]_\times) t_i$，$x_i$ 为2D像素坐标。射线坐标被调整至与潜变量相同的空间分辨率 $\mathbb{R}^{n \times h \times w}$，从而拼接后每视角的输入维度为 $(2n + 6) \times h \times w$。RF模型基于Stable Diffusion 3微调，通过条件流匹配损失（conditional flow matching）在直线插值路径 $Y_t = tY_1 + (1-t)Y_0$ 上训练向量场：

\begin{equation}
\mathcal{L}_{\text{CFM}} = \mathbb{E}_{t, Y_t, Y_1} \left[ \| u_\theta(Y_t, 1-t) - (Y_1 - Y_0) \|^2 \right]
\end{equation}

其中 $u_\theta$ 为神经网络参数化的向量场，$Y_0$ 为高斯噪声，$Y_1$ 为数据分布样本。

\textbf{射线流形投影与SD3引导。}在采样过程中，直接求解ODE可能导致生成的射线偏离有效位姿流形，降低相机轨迹精度。为此，SplatFlow在每个采样步骤 $t_k$ 预测 $t=0$ 时的目标状态，从中恢复相机位姿并重新构建射线，再将重建的射线前向传播回 $t_k$，从而将解约束在有效射线流形上。此外，为提升多视角图像的生成质量，方法在特定采样步骤引入SD3的向量场进行引导——对随机选定的视角，用SD3的向量场替换该视角的图像潜变量更新，融合SD3的强泛化能力。

\textbf{GSDecoder的深度增强与对抗训练。}在生成流程中，GSDecoder从 $K$ 视角的潜变量 $\{(I_i, D_i, R_i)\}_{i=1}^K$ 预测像素对齐的3DGS参数 $\{(\mu_j, \alpha_j, \Sigma_j, c_j)\}_{j=1}^{HWK}$。由于共享编码器 $E$ 为保持与SD3的兼容性而冻结，可能丢失细粒度空间细节，SplatFlow引入两项改进：(1)~深度潜变量集成——使用DepthAnythingV2提取深度图并归一化至 $[-1, 1]$，将深度潜变量作为GSDecoder的额外输入以增强3D结构信息，实验表明这一改进同时提升了收敛速度和最终性能；(2)~视觉辅助损失（vision-aided loss）——在GSDecoder达到一定收敛阈值后启用对抗训练以提升视觉质量，避免训练早期的不稳定性。总训练损失由LPIPS、MSE和视觉辅助损失组成。

\textbf{训练无关的3D编辑与应用。}由于RF模型与SD3共享潜空间，SplatFlow可直接利用现有的2D编辑技术实现3D任务，无需额外训练模块。对于3DGS编辑，采用适配整流流的SDEdit反演：在选定时间步 $t_k$ 对输入多视角潜变量进行反演，然后以目标文本为条件重新采样，生成编辑后的多视角潜变量，再由GSDecoder解码为编辑后的3DGS。由于RF模型联合建模所有模态，编辑过程中多视角一致性自然得到保证，无需额外的跨视角注意力机制。对于位姿估计和新视角合成等应用，已知数据（部分视角的图像或相机参数）作为约束条件，未知部分通过RePaint修复技术预测。

\textbf{实验表现。}在MVImgNet和DL3DV-7K数据集上，SplatFlow以显著更小的训练集（仅使用MVImgNet子集和DL3DV-7K，未使用完整DL3DV-10K或LAION）超越了Director3D的生成质量。在MVImgNet上，SplatFlow的FID达到34.85（Director3D为39.55），CLIP分数为31.43（Director3D为30.48）；在DL3DV上，FID为79.91（Director3D为88.44）。加入SDS++精炼后两个模型的CLIP分数均进一步提升，表明SDS++对不同3DGS生成模型的普适增益。在3DGS编辑任务上（100场景的MVImgNet基准），SplatFlow在CLIPScore和方向相似性（CLIP D-sim）上均优于优化式方法DGE，且在定性比较中可完成DGE难以实现的完整物体替换。在相机位姿估计任务中，多视角RF模型在8视角设置下超越了RayDiffusion和RelPose++等专用方法。有趣的是，不使用DepthAnythingV2深度时位姿精度反而更高——分析发现这是由于联合生成深度和射线时，模型产生了比DepthAnythingV2更精细、更具上下文感知的深度估计。单场景生成耗时约20秒，SDS++精炼约5分钟。

\paragraph{单样本生成可控 3D 场景}

ExCellGen（2024.12）~\cite{2412_16253}从单个真实世界样本（手持视频或无人机拍摄）出发生成3D场景，解决"从单例学习生成"的独特问题。与需要大规模3D数据集训练的生成方法不同，ExCellGen将3D场景生成建模为基于样例的合成（exemplar-based synthesis）任务——灵感源自纹理合成的经典范式——仅需单个真实场景作为训练数据。方法采用两阶段管线：第一阶段将不规则高斯抽象为稀疏特征化体素网格，用轻量级生成模型在亚秒级合成完整体素；第二阶段通过基于块的重新映射从原始3D高斯中恢复高质量外观。整个管线实现了从单个真实样本到可泛化场景生成的飞跃。

\textbf{语义增强的3D高斯重建。}输入阶段首先用3DGS重建输入场景。为引入超越原始外观的高层语义信息以指导生成，方法借鉴Feature-3DGS等先前工作，为每个3D高斯附加语义特征向量 $f \in \mathbb{R}^d$。具体实现采用DINO特征：从多视角图像中提取768维2D特征图，通过PCA降维至8维，再通过将可微渲染方程中的颜色 $c^{(i)}$ 替换为特征 $f^{(i)}$ 将其提升至3D空间，经优化后每个高斯获得一致的语义特征。这些语义特征同时服务于两个目的：一是为体素化和生成提供语义引导，二是通过k-means聚类实现交互式场景区域选择——用户可在GUI中调整聚类数量并直接操纵3D高斯来选取感兴趣的区域，解决了真实场景中背景杂乱、物体边界模糊的分割难题。

\textbf{体素化与特征编码。}为将不规则的3D高斯集合转化为便于生成的规整表示，方法将其抽象为稀疏特征化体素网格 $\mathcal{V} = \{v_i\}_{i=1}^N$，其中 $N$ 为占据体素数，每个体素由位置索引 $p_i \in \mathbb{R}^3$ 和特征向量 $z_i$ 组成。体素化过程的关键参数如下：仅在不透明度超过阈值 $\alpha=0.1$ 的位置实例化体素；精细分辨率 $r_f = 64^3$ 用于生成表达力丰富的场景细节，粗糙分辨率 $r_c = 16^3$ 通过下采样获得，作为条件信号。为维持体素间的多样性，每个体素选取其所包含的不透明度最高的高斯的特征（而非简单平均），提取其球谐系数（外观，4维）和DINO语义特征（4维），各自独立PCA降维后重归一化，拼接为8维联合特征向量。这一稀疏体素表示在保留场景外观和语义信息的同时，将高斯的数量级压缩至数千个占据体素，大幅降低了生成任务的计算负担。

\textbf{生成元胞自动机（GCA）。}ExCellGen的核心生成引擎是生成元胞自动机（Generative Cellular Automata, GCA）的单样例适配版本。GCA将形状建模为稀疏体素网格 $(P, Z)$ 的联合分布，从初始状态 $S^0 = (P^0, Z^0)$ 出发，通过可学习转移核 $p_\theta$ 迭代采样中间状态直至最终状态 $S^T$。转移核实现为基于稀疏卷积的U-Net，仅更新已占据体素邻域内的单元——这种类似元胞自动机的局部更新规则利用了3D形状固有的连通性和稀疏性，在高分辨率网格中大幅缩减搜索空间。核将联合概率分解为占据概率（Bernoulli分布）和特征概率（Gaussian分布，其不确定性随时间递减）两个独立项，通过Infusion训练策略优化。

为适配单样例场景，ExCellGen对标准GCA进行了关键改造：（1）~采用浅层网络架构以大幅限制感受野（相比cGCA的深层U-Net），避免对单一输入的过拟合——模型参数仅约4.2M（cGCA为468M），训练速度约6分钟（cGCA需23分钟以上），推理速度约0.2秒（cGCA约0.33秒）；（2）~引入随机裁剪增强，在精细分辨率上随机裁取边长25--30的体素块（要求至少包含250个体素，通过拒绝采样保证），每个裁剪块在各维度填充2个体素作为上下文输入但不参与监督，防止模型学习裁剪边界；（3）~以粗糙体素网格 $r_c$ 同时作为初始状态和每步条件信号，训练时采用教师强制（teacher forcing），推理时可替换为任意有效体素输入实现用户控制。生成在精细分辨率 $r_f$ 上进行，特征从标准正态分布随机初始化。$r_c$ 和 $r_f$ 的分辨率选择在可控性与输出多样性之间取得平衡——粗糙分辨率越高控制越精确，但多样性降低。

\textbf{基于块的一致性场景组合。}GCA生成的体素仅包含8维低维特征，需恢复为具有丰富外观的3D高斯场景。朴素的逐体素最近邻匹配（将每个生成体素映射到特征最接近的样例体素并复制其高斯）因忽略邻域上下文而导致空间一致性差。为此，ExCellGen设计了稀疏的基于块的一致性精化步骤，将样例的块统计信息重新分配为连贯的空间排列。具体过程定义了结合占据和特征的块距离度量 $d(P_e, P_g) = (1-w) \cdot d_\text{occ}(P_e, P_g) + w \cdot d_\text{feat}(P_e, P_g)$，其中占据距离衡量块间几何匹配度，特征距离衡量外观和语义相似度。每轮迭代中，每个占据体素中心为其块在样例中寻找最佳匹配中心，然后通过占据加权平均聚合来自多个重叠块的高斯特征，并通过投票机制适度增减几何（投票阈值 $\tau = 0.5$）。整个过程使用块大小 $l = 5$，迭代 $K = 7$ 次，限制新增体素与初始生成体的距离不超过 $\delta_\text{patch} = 2$ 个体素。精化完成后，每个合成体素选取单一样例体素并将其中所有高斯直接复制至目标位置，无需进一步优化。这一步骤对于恢复GCA可能遗漏的几何细节至关重要——实验表明，移除GCA仅使用层次化块优化的替代方案会产生重复模式、缺乏空间和语义一致性，且倾向于在条件输入之外膨胀几何。

\textbf{实验结果与训练效率。}ExCellGen在来自LeRF数据集、YouTube和自采集的多种真实场景上进行了验证，涵盖室内外环境、植被等复杂几何。所有输入视频均为手机或无人机拍摄的随意真实视频，而非合成或精心准备的数据。方法展示了多样化生成能力：给定相同粗糙体素条件，可在数秒内产生多个视觉上合理且各不相同的场景。用户还可通过样例粗体素、输入网格或直接体素编辑来控制生成，并支持将外观迁移至新几何（来自网格或另一样例场景）。训练效率方面，每个场景的GCA生成核训练不到10分钟（即使对复杂输入），训练完成后单场景生成仅需0.5--2秒（包括GCA生成和基于块的一致性步骤），支持真正的交互式场景创作。浅层网络相比cGCA深层架构在模型大小上减少约111倍（4.2M vs.\ 468M参数），训练速度提升约3.7倍，同时避免了对输入形状和外观的过拟合。方法的局限性在于：仅从单个样例训练、无外部先验，难以泛化至与原始分布偏差显著的条件输入；对于粗体素在场景中多处重复出现的复杂场景，GCA无法观察完整生成过程，可能产生"位置不当"的生成结果；重建后不进行进一步优化，可能在块边界处出现瑕疵。

\paragraph{无姿态互联网照片 3D 一致视频}

KFC-W（2024.11）~\cite{2411_13549}（Keyframe-Conditioned Video Generation in-the-Wild）设计了一种自监督方法，利用视频的时间一致性和多视角互联网照片的视角多样性，训练可扩展的3D感知视频模型，无需任何3D标注。具体任务定义为：给定同一场景的2--5张无姿态互联网照片作为关键帧，生成一段模拟相机运动路径的视角一致视频——从第一个拍摄位置移动到第二个，再到第三个，以此类推。这一任务极具挑战性：现有帧插值方法（FILM、FLAVR、LDMVFI）因训练数据中相机运动极小而无法处理宽基线视角，只能产生模糊过渡；而商业级视频生成模型（如Luma Dream Machine）虽能生成高分辨率锐利视频，却倾向于创造性的变形效果而非真实相机运动——例如在Sacre Coeur场景中，Luma将人群和阶梯变换为一栋新建筑；在林肯纪念堂场景中，它生成了第二座雕像。

这说明单纯的大规模视频合成训练不足以理解3D场景布局，需要引入可扩展的3D感知目标。KFC-W的核心思路是联合训练两个自监督子任务：多视角修复（multiview inpainting，从MegaScenes的800万张无姿态互联网照片中学习3D先验）和视角插值（从RealEstate10K和DL3DV视频中学习时序一致性）。两个任务共享同一潜空间DiT骨干，且互为补充——多视角修复从宽基线无姿态图像中学习3D几何关系，视角插值则学习生成平滑一致的相机轨迹。尽管模型从未在训练中见过"互联网照片$\to$视频"这一输入输出配对，多任务学习带来了涌现能力，使其能够处理真实互联网照片。验证表明该方法在几何和外观一致性上均优于所有基线，且其输出可服务于3DGS等下游应用。

\textbf{方法架构。}KFC-W的目标是：给定同一场景的2--5张无姿态互联网照片，生成一段视角一致的视频序列。核心挑战在于缺乏监督训练对——不存在同一场景的互联网照片与干净视频配对。为此，方法利用非配对的互联网照片和视频语料库，用同一模型联合训练两个子任务。整个架构基于潜空间Diffusion Transformer（DiT），所有图像先经预训练VAE编码至潜空间，再经patchify和扩散前向过程添加噪声。对同一图像的所有patch，添加帧索引嵌入（frame index embedding）以标识图像序号，因为Transformer本身不感知顺序。方法可扩展至更长的视频或更密集的插值帧，代价是Transformer序列长度的增加。训练视角插值任务时，模型接收17、33、49或65帧的序列，分别对应2、3、4、5个关键帧，每对相邻关键帧之间生成15帧中间帧作为计算与视觉质量的平衡点。

一个有趣的观察是，模型能根据输入格式自动推断应执行哪个子任务，无需显式标志——因为两个子任务对模型而言足够相似：条件图像传达场景身份，CLIP嵌入提供照明信息，唯一区别是输出格式和长度（可由帧索引确定）。模型从一个内部预训练的 text-to-video 模型进行 finetune（其性能可与 CogVideoX-5B 等开源模型相比），帧索引嵌入和 CLIP 嵌入等额外条件信息通过简单加法注入图像 patch token。实验中也尝试了调制（modulation）方式注入条件信息，但未观察到差异。训练在32块 NVIDIA A100 80G GPU 上进行3天，使用 MegaScenes、RealEstate10K 和 DL3DV 三个数据集。推理时运行50步 DDIM 采样，仅解码中间帧，展示生成视频时不包含输入关键帧（因为关键帧可能包含遮挡物）。

\textbf{多视角修复。}该子任务从无3D标注的非结构化图像集合中学习3D先验。模型接收从场景随机、宽基线视角拍摄的若干条件图像，并对目标图像施加80\%遮罩进行修复。模型必须从条件图像中提取结构信息和场景身份，从目标图像剩余20\%区域中推断光照和场景布局。这一过程迫使模型学习对称性等3D先验知识，以合理补全目标中未直接观测到的内容。该方法的灵感来源于CroCo的跨视角掩码图像建模，后者观察到从成对图像的交叉视角重建可以教授强3D先验。

但KFC-W在几个关键方面不同于CroCo：（1）CroCo仅处理两张图像（为每张图像创建独立的Transformer解码器，多图像时内存开销剧增），而KFC-W利用自注意力机制处理任意数量的输入；（2）CroCo依赖相机姿态和启发式规则获取训练配对，而KFC-W仅需场景标签（如Wikimedia Commons的场景标签），数据标注需求极低；（3）CroCo使用确定性的MAE目标（不允许概率性生成），而KFC-W基于扩散模型可产生多样化的输出。训练数据仅要求图像来自同一场景，无需相机姿态。值得注意的是，MegaScenes数据集包含430K场景的8M张互联网照片，但其中仅约80K场景（1/5）能被SfM成功重建——这恰好说明了绕过3D标注的必要性。实现上，条件图像的patch和目标图像80\%的patch保持清洁，仅对剩余20\%施加扩散前向过程并计算损失。此外，使用冻结的语义分割模型忽略行人和车辆等瞬态物体。

该方法可轻松扩展到其他数据源，如YouTube视频、ScanNet等合成数据、Waymo等自动驾驶数据集，进一步扩大训练规模。

\textbf{视角插值。}该子任务教会模型在起始和终止条件图像之间生成平滑一致的中间帧。从视频中随机采样 $16(n-1)+1$ 帧连续序列（$n$ 为条件图像数，$2 \le n \le 5$），每隔16帧取一张作为条件，中间15帧作为需去噪的目标。所有patch同样添加帧索引嵌入（$0, 1, 2, \ldots, (n{-}1)\times 16$ 经线性层映射）。实验中也尝试了帧索引归一化和可学习位置编码，但未发现性能差异，因此采用最简单的线性映射方案。仅靠这一目标即可教会模型在输入图像间插值，但消融实验表明，单独的视角插值不足以处理宽基线和真实互联网照片场景。推理时，给定无姿态输入图像集合 $\{x_1, \ldots, x_n\}$，模型在每对相邻输入 $(x_i, x_{i+1})$ 之间生成15帧，最终拼接为完整视频序列 $y = \{y_1, y_2, \ldots, y_k\}$，其中 $k = 15(n{-}1)$。输入的顺序决定了视频序列的相机运动路径。此外，训练时对条件图像随机遮罩部分区域，模拟语义分割的效果，因为推理时输入关键帧也需要进行瞬态物体分割。

\textbf{照明控制。}为处理互联网照片中变化剧烈的光照条件，KFC-W设计了基于CLIP的照明控制机制。训练时对每张图像条件化其CLIP嵌入：将冻结CLIP编码器的全局特征图reshape至与图像patch相同的空间尺寸，再与patch相加。然而仅靠CLIP嵌入不足以迫使模型利用该信息（模型可直接从条件图像中提取光照），因此对条件图像施加极端颜色抖动（ColorJitter），使模型必须从CLIP嵌入中提取照明特征来去噪中间帧。推理时，用户可指定一张图像的CLIP嵌入控制整个视频序列的光照风格——默认使用第一张输入图像的嵌入。实验表明，当指定某张图像的CLIP嵌入时，生成的中间帧能反映该图像的光照风格（如整体色调），且在整个序列中保持一致；而不使用CLIP嵌入时，生成的帧无法维持一致的外观。值得注意的是，CLIP嵌入可能仅包含粗粒度的照明信息（因为其训练标题仅包含"cloudy"、"sunny"等描述词，而非太阳角度等物理属性），因此该方法目前只能控制粗粒度照明（如整体色调），细粒度的照明控制仍有待未来工作。

\textbf{下游应用验证。}KFC-W在Phototourism（21个地标、每场景12组稀疏视角、共252组）和RealEstate10K（50个场景、每场景12组稀疏视角、共600组）上进行了全面评估。每组测试数据从每个场景随机采样2、3、4、5个视角各3组。

用户研究中，与帧插值基线（FILM、FLAVR、LDMVFI）的比较在10名用户中收集了250组投票（Phototourism 150组、Re10k 100组），结果表明这些方法在Phototourism上完全不具竞争力——它们训练于小基线或固定相机的数据，旨在建模动态而非相机运动；在Re10k上，仅FILM能在相机运动较小时产生合理结果。与KFC-W消融模型（Video-only，仅训练视角插值）和Luma Dream Machine的比较在42名用户中收集了927组投票（Phototourism 587组、Re10k 340组）。Video-only基线在Re10k上表现良好（因为该数据集视角较近，模型在训练中见过类似图像域），但在Phototourism的宽基线场景上严重失败，即使训练时使用了颜色抖动和随机遮罩等增强技术——这证明互联网照片训练对处理真实场景至关重要。

KFC-W在所有维度上均优于Luma Dream Machine商业模型：Luma在关键帧过渡时出现明显的光照跳变（"一致性"维度差），倾向于幻视新的建筑和结构而非理解场景布局（"相机路径"维度差），如在Sacre Coeur场景中将人群和阶梯变形为新建筑，在林肯纪念堂场景中生成第二座雕像。有趣的是，Luma在Re10k室内场景上的表现反而比Phototourism更差（物体消失等问题更明显），且当输入视角非常接近时产生随机抖动以填充时间。尽管Luma作为商业模型在采样密度、分辨率和训练数据规模上远超KFC-W，其缺乏3D感知训练目标导致无法理解场景布局，只能产生变形效果来匹配条件图像。

在SfM重建应用中，仅用原始互联网照片时COLMAP仅能成功重建45\%的场景（882张图像中仅43\%注册成功），因为宽基线视角间难以找到足够的匹配特征；加入KFC-W完整模型生成帧后成功率提升至93\%，图像注册率从43\%提升至84\%，接近翻倍——生成帧提供了可靠的跨视角特征对应关系。作为对比，加入Video-only模型生成帧后成功率仅为71\%、图像注册率67\%，远低于完整模型的93\%和84\%。

在3DGS应用中，使用InstantSplat分别对原始稀疏视角和KFC-W生成帧训练3DGS。评估方法为：对每个场景训练后，采样10张未参与训练的图像，将其姿态注册到同一坐标系，用训练好的模型渲染并报告PSNR、SSIM和LPIPS指标。Phototourism的互联网照片具有宽基线、严重遮挡和变化剧烈的光照，导致直接用原始照片训练3DGS极为困难；KFC-W生成帧因更密集的视角覆盖和更一致的光照，在PSNR、SSIM和LPIPS上均显著优于直接用原始互联网照片训练的结果。即使Re10k数据集帧间光照条件和基线距离较接近，使用KFC-W生成帧训练仍在所有指标上有所提升。需要注意的是，Phototourism的指标值显著低于Re10k，因为其测试真值本身也是互联网照片，渲染损失天然较高。

消融实验中一个重要发现是"Long-Video"基线：仅用视角插值训练但在更宽基线的视频帧上作为条件（最多5倍宽基线），该模型在输入图像重叠极小的情况下仍然失败，而KFC-W完整模型仍能理解场景布局。这证明互联网照片中包含的极端多样化视角（如极端旋转、缩放）是视频数据无法替代的——多视角修复不仅改善真实场景的表现，还提升了通用的视角插值能力。

\paragraph{四面体泼溅表示}

TeT-Splatting（2024.06）~\cite{2406_01579}提出了替代3DGS的四面体泼溅（Tetrahedron Splatting）表示，兼顾优化收敛性、精确网格提取和实时渲染。标准3DGS的各向异性高斯在几何重建中存在局限性，TeT-Splatting在结构化四面体网格中集成基于表面的体渲染，同时保留精确网格提取能力。Eikonal正则化约束 SDF 梯度范数为1，法线一致性正则化确保表面光滑，两者共同改善生成质量和稳定性。TeT-Splatting可直接集成到现有3D生成管线中，在收敛速度、渲染效率和网格质量间取得优越的平衡。

\textbf{可变形四面体网格与不透明度计算。}TeT-Splatting的几何基元是可变形四面体网格 $(V_T, T)$，其中 $V_T = \{v_n\}_{n=1}^{N}$ 为顶点位置，$T = \{t_k\}_{k=1}^{K}$ 中每个四面体 $t_k$ 由四个顶点索引 $(a_k, b_k, c_k, d_k)$ 构成。每个顶点关联一个 SDF 值 $f_n$，通过插值建立四面体内的符号距离场。与3DGS将3D高斯投影为2D泼溅不同，TeT-Splatting将四面体的四个顶点投影到图像平面，形成由四个重叠三角形构成的2D四面体泼溅。不透明度采用 NeuS 方式计算：对于光线上与四面体的两个交点（SDF值分别为 $f_{\text{prev}}$ 和 $f_{\text{next}}$），其不透明度为
\begin{equation}
\alpha_t = \max\left(\frac{\sigma_s(f_{\text{prev}}) - \sigma_s(f_{\text{next}})}{\sigma_s(f_{\text{prev}})}, \; 0\right), \quad \sigma_s(x) = (1 + e^{-sx})^{-1}
\end{equation}

其中 $s$ 控制SDF到不透明度的转换陡度，按 $s = i/s_{\max} + s_0$ 逐迭代更新。随着 $s$ 增大，不透明度行为趋近于 Marching Tetrahedra——仅穿过零等值面的四面体可见，实现从体渲染到表面渲染的自然过渡。最终的法线图、深度图和不透明度图通过前到后的 $\alpha$-混合 顺序累积：$\mathbf{C} = \sum_i \alpha_i T_i \mathbf{c}_i$，其中 $T_i = \prod_{j=1}^{i-1}(1 - \alpha_j)$ 为透射率。

\textbf{快速可微四面体光栅化器。}借鉴3DGS的分块光栅化架构，TeT-Splatting实现了基于 CUDA 的自定义四面体光栅化器。屏幕被划分为瓦片（tile），剔除视锥外的四面体后，按瓦片ID和平均深度排序。为弥补逐瓦片排序与逐像素排序的差距，对每个像素维护大小为 $N_w$ 的短重排序窗口——得益于四面体网格的结构化特性，在网格分辨率256下仅需窗口大小5即可几乎消除排序误差。此外，预过滤机制通过计算四面体不透明度上界 $\alpha_{\max}$（将公式中 $f_{\text{prev}}$、$f_{\text{next}}$ 替换为四个顶点的最大、最小SDF值），剔除 $\alpha_{\max} < T_f$（阈值 $T_f = 1/255$）的近乎透明四面体，大幅减少参与泼溅的四面体数量。

\textbf{正则化与 SDF 质量。}Eikonal 损失约束每个四面体内的 SDF 梯度范数为1：$L_{\text{eik}} = \sum_k (\|\mathbf{g}_k\|_2 - 1)^2$，防止 SDF 值快速达到极端水平并陷入局部最优。法线一致性损失将每个四面体的法线投影回顶点，对共享边的相邻顶点的法线余弦相似度进行正则化：$L_{\text{nc}} = \sum_i (1 - \langle \mathbf{n}_{e_{i,1}}, \mathbf{n}_{e_{i,2}} \rangle)$，充当平滑先验以防止表面碎裂。两者共同确保 SDF 场的质量，使 Marching Tetrahedra 提取的网格更加紧凑光滑。

\textbf{两阶段3D生成管线与实验。}TeT-Splatting采用两阶段管线：第一阶段用 TeT-Splatting 进行几何优化（hash grid 编码 SDF 和顶点变形），第二阶段通过 Marching Tetrahedra 转换为多边形网格，使用 PBR 管线进行纹理优化。与 Magic3D（Instant-NGP→DMTet 的两阶段切换常导致质量退化）和 DreamGaussian（3DGS 直接提取网格常产生可见孔洞）相比，TeT-Splatting 在网格提取后质量退化极小。配备 RichDreamer 的丰富扩散先验时，几何优化仅需40分钟（RichDreamer 原版需70分钟），在 Geometry CLIP Score 和 Appearance CLIP Score 上均优于 RichDreamer，生成时间显著缩短。消融实验证实：移除 Eikonal 损失导致模型陷入非期望形状（SDF值快速达极端水平），移除法线一致性损失导致表面碎裂，提高四面体网格分辨率（128$\to$256）则带来更精细的几何细节。

\paragraph{跨模态对齐文本-图像-高斯}

TIGaussian（2026.01）~\cite{2601_19247}面向视觉-语言-3D跨模态对齐，提出多分支3DGS标记器将高斯的内在属性解耦为结构化潜在嵌入 $F_G^I$：分别处理几何属性（位置、缩放、旋转）和外观属性（不透明度、颜色），而非将所有属性混合编码。配合双向跨模态对齐策略：(1) 图像-3D对齐利用扩散先验的多视角特征融合解决视角模糊问题，生成3D感知视觉特征；(2) 文本-3D对齐通过自适应投影模块将空间特征 $F_G^I$ 映射到文本嵌入空间 $F_G^T$，使用InfoNCE对比损失联合优化。在跨模态检索、零样本分类和场景识别等多个任务上达到最优性能。

\textbf{问题动机与现有方法局限。}在TIGaussian之前，3DGS多模态对齐的代表性工作UniGS将每个高斯的所有属性（位置、颜色、缩放、旋转、不透明度）直接拼接后统一编码，存在两个关键缺陷：(1) \emph{属性纠缠编码}——异质属性（如位置与颜色）具有截然不同的数值范围和表征语义，混合编码导致信息瓶颈下各属性的特异性信号被压缩，无法充分挖掘高斯属性的分布模式和几何意义；(2) \emph{退化3D感知}——强制从单视角图像进行对齐无法捕获全局上下文，3D特征表征的跨视角一致性在训练过程中未被充分保持，导致其他视角的匹配性能劣化。TIGaussian从属性解耦和多视角融合两个维度同时解决这两个问题。

\textbf{多分支3DGS标记器设计。}给定由高斯表示的3D对象，首先通过最远点采样（FPS）下采样至1024个高斯，再用k近邻（kNN）分组形成固定大小的局部高斯面片（local Gaussian patch），作为后续ViT编码的处理单元。不同于UniGS的拼接式编码，TIGaussian为每个属性分支设计独立的编码器 $E_\mu, E_\alpha, E_c, E_s, E_q$：空间分支 $E_\mu$ 借鉴点云网络架构，加入可学习位置嵌入模块提取高维空间上下文，通过max-pooling聚合逐点特征为全局描述符；外观分支 $E_\alpha, E_c$ 采用sigmoid激活函数建模外观空间的非线性变化并约束值域；形态分支 $E_s, E_q$ 使用归一化层标准化输出。各分支提取的特征拼接后经两层MLP融合为统一的3DGS token，整体捕获属性间的相互依赖关系。在训练策略上，3DGS token通过预训练点云模型Uni3D的交叉注意力引导进行知识迁移——将位置和颜色属性编码为点云特征，通过ViT块的交叉注意力将预训练知识整合进3DGS token，最终经全连接层投影至 $d=512$ 维得到紧凑的高维3D特征 $F_G^I \in \mathbb{R}^d$。

\textbf{扩散增强的多视角图像-3D对齐。}传统方法将随机选取的单视角图像直接与3D特征对齐，视角信息受限。TIGaussian提出3D感知对齐方案：对单视角输入图像 $I \in \mathbb{R}^{H \times W}$，利用预训练多视角扩散模型（Hunyuan3D）生成 $N$ 个多视角图像 $\{I_0, I_1, \cdots, I_N\}$，各自经CLIP编码后，通过视角感知交叉注意力模块融合为3D感知特征：

\begin{equation}
F_I^{\text{mv}} = \text{LayerNorm}\!\left(F_I + \text{Softmax}\!\left(\frac{QK^\top}{\sqrt{d}}\right)V\right)
\end{equation}

其中 $Q = F_I$ 为查询，$K = V = [F_{I_0} | F_{I_1} | \cdots | F_{I_N}] + \text{PE}(\phi)$ 为所有视角特征的拼接加上正弦位置编码 $\text{PE}(\phi)$。该机制使3DGS特征获得多视角感知能力，增强3D空间中的编码能力。

\textbf{3D-文本投影模块。}图像对齐后的特征 $F_G^I$ 与文本模态之间仍存在显著差距。TIGaussian借鉴Q-Transformer架构，使用一组可学习查询 $F_q \in \mathbb{R}^{N_q \times d}$ 作为软提示，通过 $L=6$ 层Transformer迭代精炼3D特征。每层依次执行自注意力（查询精炼）、交叉注意力（整合3D特征上下文，以 $F_G^I$ 为键值）和MLP前馈网络，各模块通过残差连接和LayerNorm串联。最终精炼后的查询经池化层生成紧凑嵌入 $F_G^T \in \mathbb{R}^d$ 用于文本对齐。整体对比学习损失函数为：

\begin{equation}
\mathcal{L} = \lambda_T \cdot L(F_G^T, F_T) + \lambda_I \cdot L(F_G^I, F_I^{\text{mv}})
\end{equation}

其中 $\lambda_T = \lambda_I = 0.5$，$L(\cdot,\cdot)$ 为InfoNCE对比损失。由于现有方法已有效解决文本-图像间的多模态对齐，该框架不显式计算文本特征与融合图像特征之间的对比损失，聚焦于3D模态的双向对齐。

\textbf{实验结果。}TIGaussian在三个公开数据集上进行了全面评估：Objaverse（146k对象）、ABO（7.9k对象）和SUN RGB-D（6.1k真实室内对象）。在零样本分类任务上，TIGaussian在Objaverse-LVIS（318类）和ABO（23类）的Top-1/Top-3/Top-5准确率上均优于所有基线方法，尤其超越同为3DGS基础的UniGS和多视角方法Duoduo-CLIP。在文本-3D检索上，于Objaverse和ABO的1000对象测试集上取得最优Top-1/Top-5/Top-10检索精度，验证了文本投影模块与多分支标记器的有效性。在图像-3D检索上同样大幅超越现有基线，证明扩散增强多视角融合有效克服了视角问题。在SUN RGB-D真实场景识别上，TIGaussian达到76.46\%的平均准确率，显著优于UniGS的68.92\%。

在少样本线性探测上（Objaverse-LVIS，每类1至16个训练样本），同等数据量下一致优于UniGS。消融实验证实：(1) 多分支标记器在所有任务上持续提升准确率，移除后性能显著退化，验证了属性解耦对3D上下文提取的核心作用；(2) 多视角特征融合相比单视角对齐大幅改善图像检索和分类，移除融合模块导致明显退化；(3) 3D-文本投影模块有效桥接了3DGS特征与文本表征之间的模态差距；三个模块联合使用时达到最优性能。扩展性分析显示，将训练数据从100k扩展至800k对象后性能持续提升，且使用更大骨干（Uni3D-S $\to$ Uni3D-B $\to$ Uni3D-L）进一步增强效果。

\textbf{训练配置与模型复杂度。}TIGaussian使用Open-CLIP ViT-B-16作为预对齐文本-图像模型，Uni3D-S点云模型用于3DGS标记引导。训练分两阶段：先在Objaverse上以AdamW优化器、学习率 $1\times10^{-4}$ 训练15个epoch（batch size 24），再在ABO和SUN RGB-D上微调20个epoch以适应各数据集特性，所有特征维度统一为 $d=512$。全部训练在4块A100 GPU上完成，推理仅需单块A100。模型复杂度方面，3DGS Tokenizer占主导开销约88.2M GFLOPs（占总计算量约80\%），图像融合模块的交叉注意力仅1M GFLOPs，3D-文本投影模块为24M GFLOPs（6层Transformer）。推理时最大GPU内存占用为32GB，可在单块RTX 4090上运行。与多视角编码方法Duoduo-CLIP相比，TIGaussian在推理时间上具有显著优势，因为后者需要处理多张图像的全部特征，而TIGaussian仅需对单个3DGS对象进行前向传播。

\textbf{多视角数量选择与过度对齐现象。}消融实验系统评估了参与图像-3D对齐的多视角数量（从2到12）对性能的影响。结果表明，增加视角数量并不一定带来持续改善——6个视角时达到最佳平衡。论文分析认为，过多视角导致过度对齐（over-alignment）问题：融合后的视觉特征变得过于特定于某些视角细节，反而难以与更抽象的文本表征对齐。此外，随着视角数增加，对多视角扩散模型生成质量的要求也相应提高，同时模型参数量成倍增长。因此6视角是性能、计算开销和生成质量之间的最佳折中。在扩散模型选择方面，对比Hunyuan3D-v1-lite（轻量版）与v1-std（标准版）生成的多视角图像以及Zero-1-to-3提供的真实渲染图像，三种来源的性能差异不大，表明框架对扩散先验质量具有一定鲁棒性。

\textbf{预训练引导的必要性与泛化性验证。}移除Uni3D预训练点云模型的交叉注意力引导后，TIGaussian性能显著退化，证实预训练知识迁移对3DGS标记器有效学习至关重要——仅靠有限规模的3DGS数据难以从零学到高质量的3D特征表征。在泛化性方面，论文使用前馈重建方法AnySplat（基于VGGT架构）从Objaverse-LVIS对象的随机3个视角生成3DGS表示，然后在这些生成的高斯对象上测试TIGaussian的嵌入相似度热图。结果表明TIGaussian对3DGS对象的数据来源不敏感，在重建高斯上仍保持良好的语义区分能力，验证了其作为通用3DGS编码器的泛化潜力。在更广泛的基线对比中（包括OpenShape、TAMM、MixCon3D、ReCon/ReCon++等方法），TIGaussian在仅使用100k训练实例（每对象1024个高斯点）的条件下，仍然优于使用876k实例（每对象10000个点）训练的基线方法，进一步证明了3DGS表征与属性解耦编码策略的高效性。

\textbf{局限性与未来方向。}TIGaussian仍存在两方面局限：(1) \emph{泛化能力边界}——在遮挡严重的多物体场景或真实户外场景中，性能可能退化；(2) \emph{文本标签依赖}——文本-3D对齐质量本质上依赖训练标签，当前基准主要使用LLM自动生成的标注（如InternLM-XComposer），可能引入偏差。论文建议未来探索混合监督范式，结合LLM的高效标注与人工专家数据的精确性，以实现更好的细粒度文本检索。

\paragraph{去偏扩散先验 3D 注意力}

TD-Attn（2025.12）~\cite{2512_07345}从理论角度深入分析了文本到图像扩散模型中的先验视角偏差（prior view bias）。由于训练数据中视角分布严重偏斜（先验视角远多于其他视角），当主体词（如"PIKACHU"）缺乏视角标注时，交叉注意力会优先激活先验视角特征。偏差系数 $\beta$ 满足概率比 $R = p_D(v_{prior}|y_{obj})/(1-\beta) + p_D(v^*|y_{obj})$——当 $R > 1$ 时先验视角压倒目标视角，导致视角间外观冲突。基于此分析，方法提出3D感知注意力引导模块（3D-AAG）构建视角一致的3D注意力高斯，以及层次化注意力调制模块（HAM）通过语义引导树定位和调制高度响应视角条件的注意力层。TD-Attn可作为通用插件显著增强3D任务的多视角一致性。

Physically Aware 360（2025.12）~\cite{2512_10293} 将物理感知引入 360 度全景 3D 生成。方法的核心洞察是：现有全景3D生成方法（如Text2Light、PanoSwin等）仅关注视觉质量，而忽略了生成内容的物理合理性——光源方向、阴影投射和材质反射应与几何结构一致。方法设计了物理约束的生成管线，在扩散模型的去噪过程中注入物理先验（光照模型和阴影估计），使生成的全景内容不仅视觉逼真，而且在物理模拟和AR/VR场景中可直接使用。

\begin{takeaway}
3D 生成模型与 GAN 从逐场景优化走向前馈大规模生成。LRM 系列用 Transformer 实现秒级 3D 重建，3DGS 成为生成模型的理想输出表示。TD-Attn 揭示了扩散模型中先验视角偏差的理论根源，Physically Aware 360 将物理约束注入生成管线。核心趋势是从纯视觉生成走向物理感知生成。
\end{takeaway}

%% ============================================================


\section{其他3D生成方法}
\label{sec:08_5}

除上述主流技术路线外，还有一些方法从特殊应用场景或非标准输入出发进行3D高斯生成，涵盖了从点云到高斯的转换、大规模地球表面生成、程序化建筑生成、内部纹理生成、城市级别场景生成等方向。

至此，我们从文本/图像到3D、4D动态生成、3D生成模型三条主线系统地梳理了3DGS生成技术。本节将视野拓展到一些不那么主流但同样重要的方向。它们从特殊的输入形式（如无颜色点云）或特殊的应用规模（如城市级别）出发，展现了3DGS生成技术的多样性和潜力。

\subsection{点云转换与扩展应用}
\label{sec:08_5_1}

本节讨论3条技术路线：（1）点云到高斯转换；（2）大规模地理与城市场景生成；（3）特殊应用场景。

\subsubsection{点云到高斯转换}

本路线介绍从原始3D点云（尤其是无颜色点云）生成高保真3D高斯表示的方法，将最常见的3D数据形式转化为可渲染的3DGS资产。

3D点云是最常见的3D数据形式之一，但从原始无颜色点云生成高保真3D高斯仍是一个未解决的问题。

\paragraph{文本引导点云高斯化}

原始3D点云是最常见、最易获取的3D数据形式——LiDAR、深度相机等设备输出的就是点云。然而，这些点云既没有颜色属性，也缺乏拓扑连接信息，无法直接用于高质量渲染。传统做法是先将点云重建为网格、再进行UV展开和纹理贴图，但自动UV展开常导致纹理重叠、碎片化和扭曲，严重制约最终渲染质量。3DGS天然免去了UV参数化的需要，因此"将原始无颜色点云直接转换为高保真3D高斯"成为一个极具吸引力的目标。GAP（2025.08）~\cite{2508_05631}正是为此设计的方法：给定文本提示词$c$和点云$P = \{p_i\}_{i=1}^N$，目标是生成高保真3D高斯$G = \{g_i\}_{i=1}^M$。

\textbf{高斯初始化。}GAP的初始化策略充分利用了点云固有的几何信息。首先，高斯中心直接设为输入点云坐标，提供了对3D形状表面的粗略几何描述。在此基础上，方法利用CAP-UDF从点云学习一个神经无符号距离场（UDF）$f_u$，并通过梯度推断法线：$n_i = \nabla f_u(p_i) / \|\nabla f_u(p_i)\|$。值得注意的是，GAP选择2DGS（oriented disk）而非标准3DGS作为表示——2DGS将3D高斯椭球体替换为2D定向盘片，天然将法线编码为盘片方向，在表达精细局部几何时性能更优。每个盘片的旋转矩阵根据推断的法线初始化，确保盘片方向与局部表面对齐。

\textbf{渐进式多视角生成。}核心挑战在于：点云没有颜色，高斯的外观从何而来？GAP的答案是利用文本引导的深度感知扩散模型逐视角生成外观图像，作为高斯优化的监督信号。具体而言，对于预设的视角序列$\{v_j\}_{j=1}^K$，方法渐进地为每个视角生成外观：利用学习到的UDF场通过射线行进计算深度图$D_j$，将当前渲染图像$I_j$、深度图$D_j$、掩码$M_j$和文本提示$c$一起送入基于ControlNet的深度感知修复扩散模型，生成几何一致的外观图像。这里的关键设计是三区域掩码策略——掩码被划分为生成区$M_{\text{generate}}$（从未生成过的空白区域）、保留区$M_{\text{keep}}$（已处理且当前视角无更优观察角度的区域）和更新区$M_{\text{update}}$（当前视角提供更优观察角度的区域）。更新区的判定基于视角方向与法线的余弦相似度：$M_{\text{update}}^j = 1$当且仅当$M_{\text{similarity}}^j > M_{\text{similarity}}^{\text{others}}$。这一机制确保每个表面区域始终由最佳视角的外观主导，有效避免了邻接视角融合产生的伪影和不一致性。

\textbf{高斯优化与几何约束。}与标准3DGS在多个视角间反复迭代优化不同，GAP对每个视角仅执行单次优化——这让优化过程更加稳健，忠实呈现每个视角的高质量外观。优化时，方法仅修改沿视角射线可见的最近表面层高斯，而不触及背面的高斯，避免不必要的信息破坏。为维持几何精度，GAP引入了两个关键约束：(1)~表面锚定损失（Surface-Anchoring）将高斯中心约束在学习到的UDF零等值面上，$L_{\text{Distance}} = \|f_u(\mu_i)\|_2$，防止高斯在优化过程中飘离形状表面——飘离的高斯会在后续视角中产生错误的遮挡关系，导致掩码畸变并连锁降低外观生成质量；(2)~尺度损失约束每个高斯的最大尺度，$L_{\text{Scale}} = (\max(\max(s_i), \tau) - \max(s_i))^2$，防止单视角优化中产生过大的高斯。最终的优化目标为$L = L_{\text{Rendering}} + \lambda_{\text{Distance}} L_{\text{Distance}} + \lambda_{\text{Scale}} L_{\text{Scale}}$，其中渲染损失包含$L_1$项和D-SSIM项。

\textbf{基于扩散的3D空间修复。}即使经过多视角密集采样，物体仍存在难以观察的区域（如底部凹陷、内腔等）。GAP提出了直接在3D空间操作的扩散修复策略，利用可见高斯的空间关系和几何一致性推断不可见高斯的属性。对于每个不可见高斯$g_j^*$，方法定位其$L$个最近邻已优化高斯，通过加权扩散预测其属性：颜色权重综合了空间距离（防止远处高斯的颜色干扰）、法线一致性（$n_i \cdot n_j^* > 0.5$，即角度小于60度，保证颜色沿表面传播而非跨表面泄漏）和Opacity可靠性（高Opacity的邻居影响力更大）；尺度根据邻居距离的均值设定，确保稀疏区域有更大的高斯覆盖；Opacity则与局部密度成反比，防止高密度区域颜色过度累积。这一纯3D空间的修复策略无需额外网络推理，计算高效。

\textbf{实验评估。}GAP在多个层面进行了评估。在Objaverse子集（410个多类别物体）的文本驱动外观生成任务上，GAP在FID、KID和CLIP Score上均超越了依赖UV映射网格的纹理生成方法（TexTure、Text2Tex、Paint3D、SyncMVD），即使这些方法使用了艺术家创建的UV贴图作为对比基准。在Point-to-Gaussian生成任务上，GAP相比DreamGaussian、TriplaneGaussian和DiffGS均产生更高质量的结果——前者因SDS优化导致颜色过饱和，后两者受限于低分辨率三平面表示的细节捕捉能力。此外，GAP还能成功处理真实扫描的部分点云（SRB和DeepFashion3D数据集），表面锚定机制有效将分裂和克隆的高斯拉向缺失区域，同时保持几何一致性。更令人印象深刻的是，GAP可扩展至场景级别——在3D-FRONT室内场景和真实世界3D Scene数据集上，相比Paint3D（在场景数据上失败）和SceneTex（需要额外的VSD优化和LoRA训练），GAP仅需单次优化过程即可产生高质量结果。

\paragraph{从 3D 点云生长高斯}

GaussianGrow（2026.04）~\cite{2604_05721}从另一个角度解决点云到高斯的问题：以文本引导的多视角扩散模型为监督，从3D点云逐步"生长"出高斯。给定输入点云$P = \{p_i\}_{i=1}^N$和文本提示词$c$，目标是生成高保真3D高斯$G = \{g_i\}_{i=1}^M$。

\textbf{高斯初始化与几何信息图。}初始化策略与GAP类似——高斯中心直接设为点云坐标，利用CAP-UDF从点云学习无符号距离场（UDF）$f_u$，通过梯度推断法线$n_i = \nabla f_u(p_i) / \|\nabla f_u(p_i)\|$。GaussianGrow同样采用2DGS（oriented disk）表示，根据法线初始化旋转矩阵。与SDF需要封闭表面不同，UDF能够表示开放和复杂拓扑结构，更适合实际点云数据。为进一步利用点云的几何先验，方法从学习到的距离场中提取三种几何信息图：深度图$D_i$（通过射线行进在UDF中计算，用于引导Depth-Aware-ControlNet）、法线图$N_i$（通过推断UDF零等值面上的梯度获得）和位置图$C_i$（将每个像素映射到其对应点的精确XYZ坐标）。这三种信息图分别服务于外观生成的不同阶段——法线图和位置图为多视角扩散模型提供几何引导，深度图为单视角外观生成提供条件控制。值得注意的是，整个管线不需要显式的网格表示，所有几何信息均直接从点云及其距离场推导。

\textbf{重叠区域检测与额外视角优化。}多视角扩散模型的标准配置是生成六个基本视角的外观图像，但相邻视角的重叠区域经常出现外观不一致。为解决这一问题，GaussianGrow提出了密集视角生成框架，扩展到$K=10$个视角。具体而言，方法首先从每个视角$v_i$通过射线追踪检测可见高斯集合，然后计算相邻视角对$v_i$和$v_j$的可见高斯交集来确定重叠区域$R_{i,j}$——仅关注同时从两个视角可见的高斯。该过程采用了CUDA并行实现，将计算时间从分钟级降至秒级。对于每个重叠区域$R_{i,j}$，方法引入一个可学习的额外相机姿态$T_{i,j}$，通过优化使其对齐局部几何结构。优化目标为最小化相机射线方向与重叠区域高斯法线之间的对齐损失：

\begin{equation}
L_{\text{align}} = \sum_{g \in R_{i,j}} \left(1 - \frac{|\langle \mathbf{d}_{i,j}, \mathbf{n}_g \rangle|}{\|\mathbf{d}_{i,j}\| \cdot \|\mathbf{n}_g\|}\right),
\end{equation}

其中$\mathbf{d}_{i,j}$为从相机位置$T_{i,j}$发射的射线方向，$\mathbf{n}_g$为高斯$g$处的法线。可学习相机姿态被约束在以物体为中心的单位球面上，且始终朝向物体中心。最小化$L_{\text{align}}$确保额外相机被最优定位以与局部几何结构对齐，减少投影畸变并增强视角一致性。

\textbf{两阶段外观生成与高斯优化。}外观生成采用两阶段策略。第一阶段，利用Depth-Aware-ControlNet和Stable Diffusion根据深度图$D_i$和文本提示$c$生成高质量参考外观图像$I_r$，作为多视角一致性的锚点。然后以法线图$N_i$和位置图$C_i$为几何条件，结合参考外观，通过多视角扩散模型（实践中采用Hunyuan3D-Paint）在全部$K=10$个视角生成高保真外观图像$\{I_i\}_{i=1}^K$。第二阶段为高斯优化，同样采用两步策略：先处理六个基本视角$V = \{v_i\}_{i=1}^6$，再处理四个重叠区域额外视角。对于每个基本视角$v_i$，方法仅优化从该视角可见的正面高斯，不修改背面高斯——这种视角特异性优化方案确保已被良好优化的高斯不会在后续步骤中被不可见的视角错误修改。基本视角优化完成后，再在检测到的重叠区域$R_{i,j}$上利用额外相机位置$T_{i,j}$进行精调，仅优化重叠区域内的可见高斯$g_i \in G_{R_{i,j}}$。这种顺序策略确保基础外观先被建立，再处理过渡区域的潜在不一致。

\textbf{迭代式高斯修复。}尽管多视角外观生成覆盖了物体大部分表面，某些区域（如底部凹陷、内腔等）仍可能未被充分捕获。由于不同物体具有多样的几何结构，使用固定的密集视角集难以完全覆盖所有部分。为此，GaussianGrow提出基于可见性的迭代修复方案，逐步在不可见点云区域生长高斯。核心思想是分析高斯之间的遮挡关系——将高斯投影到当前视角的2D图像平面上，评估深度关系来量化遮挡。对于每对高斯$(g_i, g_j)$（其中$g_i$未优化，$g_j$已优化），遮挡损失定义为：

\begin{equation}
L_{\text{occ}} = \sum_{i,j} \sigma\!\left(\frac{((r_i + r_j)^2 - \|q_i - q_j\|^2)}{\tau}\right) \cdot \sigma(z_i - z_j),
\end{equation}

其中$q_i, q_j$为高斯在图像平面上的2D投影，$r_i, r_j$为各自的投影半径，$z_i, z_j$为距相机的深度，$\sigma$为sigmoid函数，$\tau$为控制比较锐度的温度参数。当$g_i$被$g_j$遮挡（即投影重叠且$z_i > z_j$）时损失值增大。这一可微公式使相机姿态可以通过梯度下降优化，系统地发现能揭示最大未观测区域的视角。

定位到最优相机位置$v_i$后，方法在2D空间检测未优化高斯生成掩码$M$，将掩码、深度图$D_i$、渲染遮挡图像$I_i$和文本提示$c$一起送入基于Stable Diffusion和ControlNet的深度感知修复扩散模型，生成完整的高质量图像$\hat{I}_i$。利用$\hat{I}_i$中修复部分的监督，优化不可见高斯。该迭代过程持续进行——每一步识别能观察到最大剩余未观测区域的相机姿态，直到所有高斯生成完毕。实验表明，对于大多数模型六次迭代即已足够。该方法的关键优势在于对复杂几何的自适应性：动态识别并优先处理遮挡区域，而非依赖预定义的视角模式，以更少视角实现更完整的覆盖。

\textbf{空间修复后处理。}由于原始点云数据中的噪声和密度不均匀，经过基于图像修复的高斯修复后仍可能存在难以观察的点。作为后处理步骤，GaussianGrow采用空间修复策略完成残余未优化高斯——从已优化的邻居高斯向附近的未优化高斯传播属性，无需额外网络推理。

\textbf{实验评估。}GaussianGrow在多个任务上进行了全面评估。在文本驱动外观生成任务上，以Objaverse子集（410个多类别3D模型）为基准，GaussianGrow在FID、KID和CLIP Score上均超越了依赖UV映射网格的方法（TexTure、Text2Tex、Paint3D、SyncMVD），包括在重建网格上运行的基线版本——这表明点云到网格的转换过程会平滑或扭曲几何细节，而GaussianGrow直接在3D空间优化高斯基元，完全规避了UV参数化问题。在文本到3D生成任务上，GaussianGrow支持两种几何设置：基于检索（使用Uni3D从G-Objaverse检索点云）和基于生成（使用LGM生成支架点云）。在T3Bench基准上，检索版本"Ours+Uni3D"在CLIP相似度、CLIP R-Precision和ImageReward三项指标上均取得最优成绩，生成版本"Ours+LGM"也达到与SOTA方法DiffSplat可比的性能。

在Point-to-Gaussian生成任务上，无论使用Objaverse合成点云还是DeepFashion3D真实扫描点云，GaussianGrow均优于DreamGaussian（SDS优化导致颜色过饱和且需昂贵参数敏感优化）和TriplaneGaussian（受限于三平面分辨率无法恢复精细外观和复杂几何）。消融实验验证了重叠区域策略和迭代修复机制的贡献：仅使用六个基本视角会导致所有指标明显退化，添加四个针对关键重叠区域的额外视角取得最佳性能，进一步增加视角数收益可忽略——说明相机姿态优化已捕获最关键的重叠区域。

\subsubsection{大规模地理与城市场景生成}

本路线介绍将3DGS生成扩展到城市、地理甚至地球表面等大规模场景的方法，核心挑战在于海量高斯的存储效率和跨尺度一致性。

当生成目标从单个物体扩展到城市或地球表面级别时，3D高斯的存储和计算面临量级挑战。

\paragraph{大规模地球表面生成}

EarthCrafter（2025.07）~\cite{2507_16535}致力于将3D生成扩展到数千平方公里的地理尺度。作为该领域的锚点工作，EarthCrafter在数据基础设施和模型架构两个层面进行了系统性创新，其整体生成管线将结构与纹理完全解耦：先由StructFM生成结构潜变量、StructVAE解码为高分辨率体素坐标，再由TexFM生成纹理潜变量、TexVAE解码为2DGS作为最终3D表示。

\textbf{Aerial-Earth3D数据集。}方法首先构建了迄今最大的3D航拍数据集Aerial-Earth3D，包含5万个人工筛选的600m$\times$600m场景，覆盖美国本土45M多视角Google Earth帧。每个场景提供位姿标注的多视角图像、深度图、法线图、语义分割和相机位姿，并经过显式质量控制以确保地形多样性。这一数据集为大规模地理生成提供了丰富的先验知识。

\textbf{双稀疏3D-VAE架构。}EarthCrafter的核心创新是将几何与纹理压缩完全解耦的双稀疏VAE设计，大幅缓解大规模地理数据带来的计算负担。

StructVAE负责压缩几何信息。它将全分辨率体素坐标$V_c$压缩为空间尺寸仅为原始1/256的结构潜变量$S_{La}$（32通道），通过4层Transformer加一层步长为2的稀疏3D卷积实现下采样。上采样路径则提出了创新的PSS（Pseudo-Sparse to Sparse）块：先用定制的像素洗牌层对稀疏表示进行上采样，产生"伪稀疏"体素（其中部分体素实际无效），再通过体素分类模块恢复有效稀疏体素，实现稀疏几何的精确上采样。实验表明StructVAE达到97.1\%的结构重建精度，为后续生成阶段节省了大量计算。

TexVAE负责压缩纹理信息。大规模场景下直接使用高维特征（如DINOv2的1024维特征，每个场景需约471M存储）在I/O上不可承受。为此，TexVAE采用分层FLUX-VAE特征（16通道）代替粗粒度大通道特征：将1/1、1/2、1/4三个尺度的特征拼接为48维，加上十字形RGB像素（15维）和法线特征（3维），构成66维的体素特征$f_{\text{feat}}$，每个场景仅占31M存储（DINOv2方案的6.4\%）。TexVAE编码器将体素特征编码为8通道纹理潜变量$T_{La}$，解码器将其恢复为2DGS表示（包含偏移、缩放、不透明度、旋转矩阵和球谐系数）。损失函数包含$L_1$、LPIPS（VGG+AlexNet双网络）和SSIM，兼顾视觉质量和像素精度。

\textbf{粗到细流匹配生成。}EarthCrafter分别训练结构流匹配模型StructFM和纹理流匹配模型TexFM，两者均支持灵活的条件输入（图像、语义分割或无条件生成）。

StructFM采用粗到细的两阶段流程。粗阶段由纯Transformer块构成，将随机初始化的噪声网格分类为激活/未激活体素（建模二项分布$p(G_{\text{class}}|G, C_{\text{cond}})$），其中条件网络基于Swin Transformer通过加法或交叉注意力机制将条件信号与体素特征对齐。图像条件需先估计单目深度，语义和无条件输入则使用固定深度值$z=128$。细阶段基于Swin注意力的U-Net，以上阶段阈值过滤后的结果作为稀疏输入，预测结构潜变量$S_{La}$；此外还利用潜变量特征进一步精化体素坐标——将无效体素的特征置零，保留超过50\%通道阈值$>0.3$的体素。为缓解两阶段间的域差异，方法提出体素膨胀增强策略强化细阶段训练。

TexFM在StructVAE解码得到的体素坐标$V_c$上生成纹理潜变量$T_{La}$。面对大规模场景中单场景高达22万体素的计算挑战，TexFM采用互补的U-Net架构：Swin注意力层擅长捕捉1/1、1/2、1/4高分辨率尺度下的局部特征，全注意力层则在1/8低分辨率尺度提供更广的感受野用于全局特征学习，在纹理质量和学习效率间取得良好平衡。

\textbf{实验评估。}EarthCrafter在极大规模生成任务上显著优于现有方法。框架支持多种灵活应用：语义引导的城市布局生成、图像条件的地形重建、无条件地形合成等，同时通过Aerial-Earth3D的丰富数据先验保持地理合理性。双稀疏解耦设计使方法能够处理数千平方公里的地球表面，突破了现有3D生成方法在空间尺度上的根本限制。

\paragraph{无界 3D 城市生成}

当生成目标从单个物体扩展到城市级别时，3D高斯的存储和计算面临量级挑战——无界场景要求将高斯数量扩展到十亿级，VRAM需求可达数百GB。现有的NeRF方法（如InfiniCity和CityDreamer）沿射线以相同密度采样并聚合所有点的颜色值，导致推理效率低下且细节丢失；而基于3DGS的方法（如DreamGaussian、GaussianDreamer、GaussianEditor）仅生成有限尺度的物体或场景。GaussianCity（2024.06）~\cite{2406_06526}是首个面向无界3D城市生成的3DGS生成模型，其核心洞察是引入BEV-Point（鸟瞰点）作为高度紧凑的中间表示，将高斯属性解耦为位置相关属性和风格相关属性两部分。

\textbf{BEV-Point表示。}在BEV-Point表示中，只有当前帧可见的鸟瞰点被保留，因此VRAM使用量不随场景尺度增长。具体实现上，鸟瞰语义图 $S$ 和高度场 $H$ 定义BEV点坐标：

\begin{equation}
C_F = \{(x,y,z) \mid H_{(x,y)} = z \text{ and } D_{(x,y)} = 1\}
\end{equation}

其中二值密度图 $D$ 根据语义类别调整采样密度（道路/水面等简单纹理区域降低密度，建筑立面等复杂纹理区域增加密度）。通过射线求交获得二值可见性图 $V: (x,y,z) \to v$（$v \in \{0,1\}$）过滤不可见点后，可见BEV点坐标为：

\begin{equation}
C_A = \{(x,y,z) \mid (x,y,z) \in C_F \text{ and } V_{(x,y,z)} = 1\}
\end{equation}

射线求交策略是最高效的可见性过滤方法，相比区域选择和实例选择分别减少97.7\%和96\%的总点数。密度图同样至关重要——移除密度图将使点数增加20\%--50\%。经过两层过滤后，典型可见BEV点数量远少于原始的2000万。BEV-Point表示在训练过程中保持恒定的VRAM使用量，而标准3DGS随场景规模增长导致VRAM不可承受，同时显著减少了文件存储开销。

\textbf{BEV-Point特征生成。}BEV-Point表示中的特征分为三类：实例属性、BEV点属性和风格查找表。

实例属性包含语义图 $S$ 提供的语义标签以及实例图 $Q = \text{Inst}(S)$（通过对语义图检测连通分量实现实例化）提供的实例标签 $I$。每个实例 $l \in L = \{1, 2, \ldots, N_{\text{ins}}\}$ 具有3D包围盒尺寸 $S(l) \in \mathbb{R}^3$ 和中心坐标 $C(l) \in \mathbb{R}^3$。

BEV点属性包括绝对坐标 $C_A$（世界坐标系中心为原点）和相对坐标 $C_R$（以各实例中心为原点的归一化坐标）：

\begin{equation}
C_R^i = \frac{2(C_A^i - C(I^i))}{S(I^i)}
\end{equation}

此外，通过场景编码器 $E$ 从BEV图提取场景特征 $F_S \in \mathbb{R}^{N_{\text{pt}} \times C_{F_S}}$，并使用绝对坐标索引：$G = E(H, S)$，$F_S = G_{\{(x,y) | (x,y,z) \in C_A\}}$。

风格查找表将所有实例的外观编码为一组潜向量 $Z_T \in \mathbb{R}^{N_{\text{ins}} \times C_Z}$，其中每个实例的风格码 $z_T^i \sim \mathcal{N}(0,1)$ 通过查找表查询 $z_T^i = T(l)$ 获得。这一设计将外观参数从逐高斯存储压缩为逐实例共享，大幅降低了计算开销。

\textbf{BEV-Point解码器。}BEV-Point解码器从BEV点特征生成高斯属性 $A$，包含五个关键模块：位置编码器、点序列化器、点Transformer、调制MLP和高斯光栅化器。

\textbf{位置编码器}将每个点坐标和对应特征经过拼接操作后，变换为高维位置嵌入向量 $F_P = P(\text{Concat}(C_R, F_S))$，其中位置编码函数 $P(x) = \{\sin(2^i x), \cos(2^i x)\}_{i=0}^{N_{\text{PE}}-1}$ 对每个特征值独立应用正弦-余弦编码，将低维坐标映射到高维频率空间以增强空间位置的表达能力。

点序列化器（Point Serializer）是BEV-Point解码器的核心创新。与NeRF沿射线采样天然保持空间相关性不同，BEV点和3D高斯作为无序点云缺乏空间结构，直接用MLP处理无法捕获结构与上下文特征。点序列化器 $L: (x,y,z) \to o$（$o \in \mathbb{Z}$）将点坐标转换为反映空间顺序的整数索引：

\begin{equation}
L(x, y, z, g) = x \cdot g^2 + y \cdot g + z
\end{equation}

其中 $g$ 为控制离散空间尺度的网格大小。此排序确保相邻索引的点在空间上接近，从而将无结构点云转化为结构化格式。序列化后的特征 $F_P^S$ 按序列化顺序重排BEV点。实验表明，该方法与Hilbert曲线序列化效果相当，但计算复杂度显著更低。

点Transformer处理序列化后的特征，输出 $F_T = \mathcal{F}(F_P^S)$。调制MLP以风格码 $Z_P$ 和实例标签 $I$ 为条件，结合位置特征和Transformer特征生成高斯属性 $A = M(\text{Concat}(F_P, F_T), Z_P, I)$。值得注意的是，由于BEV-Point初始化保证了所有点均匀分布在物体表面上，高斯属性中除RGB外的其他属性（缩放 $s=1$、旋转四元数 $q=[1,0,0,0]$、不透明度 $\alpha=1$）均可使用默认值——实验表明，同时优化过多属性不仅增加网络收敛难度，反而显著降低生成质量。

\textbf{损失函数。}生成器采用混合训练目标，结合重建损失和对抗学习损失：

\begin{equation}
\mathcal{L} = \lambda_{\text{L1}} \mathcal{L}_{\text{L1}}(I_G, I_R) + \lambda_{\text{VGG}} \mathcal{L}_{\text{VGG}}(I_G, I_R) + \lambda_{\text{GAN}} \mathcal{L}_{\text{GAN}}(I_G, S_G)
\end{equation}

其中 $I_R$ 为真实图像，$S_G$ 为透视视图语义图。损失权重 $\lambda_{\text{L1}}$、$\lambda_{\text{VGG}}$ 和 $\lambda_{\text{GAN}}$ 初始设为10、10和0.5，$\lambda_{\text{L1}}$ 和 $\lambda_{\text{VGG}}$ 每2万轮迭代减半。

\textbf{训练细节。}生成器使用Adam优化器（$\beta = (0, 0.999)$，学习率 $10^{-4}$），判别器使用Adam优化器（$\beta = (0, 0.999)$，学习率 $5 \times 10^{-7}$），共训练20万次迭代，批量大小为8，在8块NVIDIA Tesla V100 GPU上进行。网格大小 $g$ 设为0.01，位置编码维度 $N_{\text{PE}}$ 设为64，特征通道 $C_Z$、$C_{F_S}$、$C_{F_P}$、$C_{F_T}$ 分别设为256、61、1280和256。图像随机裁剪为GoogleEarth数据集的 $448 \times 448$ 和KITTI-360数据集的 $448 \times 224$。

\textbf{实验评估。}GaussianCity在两个基准数据集上进行评估。GoogleEarth数据集包含400条曼哈顿和布鲁克林的环绕轨迹（每条60帧图像，环绕半径125--813米，海拔112--884米），提供语义分割和建筑实例标注；KITTI-360数据集涵盖约73.7公里驾驶距离的城市场景，具有丰富的光影变化和3D包围盒标注。在无人机视角（GoogleEarth）上，GaussianCity在FID和KID指标上显著优于SGAM、PersistentNature、SceneDreamer、InfiniCity和CityDreamer；CityDreamer因复合两个神经场在图像拼接处产生伪影，而GaussianCity基于灵活点云表示避免了此问题。在街景视角（KITTI-360）上，GaussianCity同样优于StyleGAN2、GSN、GIRAFFE和UrbanGIRAFFE。

GaussianCity的推理速度比CityDreamer快60倍（在NVIDIA Tesla A100上生成 $960 \times 540$ 图像）。用户研究（20名志愿者从感知质量、3D真实感、视图一致性三个维度1--5分评估）进一步验证了GaussianCity的显著优势。消融实验证实了BEV-Point表示（可见性过滤+密度图）、点序列化器和点Transformer各组件的不可或缺性。此外，GaussianCity支持建筑风格插值（通过改变风格码 $z$ 实现平滑过渡）、局部编辑（仅修改指定建筑的风格码）以及重光照（将3D高斯导入Unreal Engine实现逼真光照效果）。方法的主要局限在于：BEV-Point初始化基于曼哈顿假设无法表示中空结构，且仅预测RGB属性未充分利用3DGS的表达能力。

\paragraph{程序化建筑生成}

Proc-GS（2024.12）~\cite{2412_07660}从游戏和影视工业中成熟的程序化建模（Procedural Modeling）技术汲取灵感，提出了首个将程序化代码集成到3DGS框架中的管线。核心观察是：建筑具有天然的层级结构和重复模式——同一栋建筑在不同楼层重复使用相同的窗户、角柱、墙体等基础资产，而工业界长期依赖程序化规则来高效构建虚拟城市。Proc-GS将这一先验引入3DGS训练，使高斯表示同时获得高质量渲染和灵活的结构化控制能力。

\textbf{程序化代码定义。}建筑的层级结构被编码为一段程序化代码字符串（procedural code）。具体而言，建筑沿垂直方向被分解为若干层（Level），每层由一组基础资产（base asset）按特定排列方式组成。例如，第1层可能包含角柱 C\_E 和窗户 W1，而第2、3、5层可能使用完全相同的资产排列。该编码还支持可缩放资产（标记为 $*$）以适应不同建筑尺寸，以及可重复组合（用括号 $(\,)$ 括起）描述层内周期性排列。方法基于UE5 City Sample项目构建了MatrixBuilding数据集，包含17栋风格各异的标志性建筑（模仿芝加哥、旧金山和纽约建筑风格），每栋建筑配有稠密多视角图像、相机位姿和真实程序化代码。

\textbf{真实场景程序化代码提取。}对于真实世界建筑，Proc-GS设计了从多视角图像自动提取程序化代码的流程：首先用2DGS获取具有良好几何结构的点云和网格，然后自动检测建筑立面，对每个立面渲染一张正对视角并标注2D程序化代码，最后将其投影到网格表面获得3D程序化代码。

\textbf{资产获取阶段（Asset Acquisition）。}资产获取阶段的目标是在3DGS训练过程中自动提取基础资产，而非依赖艺术家手工制作。给定程序化代码，每个基础资产 $i$ 的参数包括：包围盒尺寸 $(x_e^i, y_e^i, z_e^i)$、枢轴位置 $(x_c^i, y_c^i, z_c^i)$，以及该资产 $K$ 次实例化的变换集合 $\{[R_1^i, T_1^i, S_1^i], \ldots, [R_K^i, T_K^i, S_K^i]\}$。高斯初始化时，第 $i$ 个基础资产分配到的粒子数为：

\begin{equation}
N^i = N \cdot \frac{V^i}{\sum_j V^j}, \quad V^i = x_e^i \cdot y_e^i \cdot z_e^i,
\end{equation}

即按包围盒体积比例从总粒子数 $N$ 中分配。渲染时，所有基础资产按程序化代码组装为完整建筑：第 $j$ 次实例化的高斯中心、旋转和缩放更新为 $\mu' = R_j^i S_j^i \mu + T_j^i$，$R' = R_j^i R$，$S' = S_j^i S$。关键设计在于，重复出现的基础资产共享同一组高斯参数，渲染时同一图像中多次出现的同一资产将梯度回传至共享参数，实现同步更新——这构成了一种隐式的数据增强。

\textbf{方差资产与边界自适应控制。}由于同一基础资产的不同实例在外观和几何上存在细微差异，Proc-GS为每个实例额外学习一个方差资产（variance asset），捕获外观差异和几何变化。此外，标准3DGS训练中高斯基元可能越过资产边界扩散，导致资产耦合、难以编辑。为此，方法在密度控制中引入Bbox自适应裁剪（Bbox Adaptive Clamp）：对超出包围盒的高斯，将其尺度减半并将中心拉回边界。裁剪操作在略大于真实包围盒的"软边界"上执行（扩展约20cm），每100次迭代执行一次。

\textbf{资产装配阶段（Asset Assembly）。}提取基础资产后，通过操纵程序化代码即可生成多样化建筑。具体流程包括两个子模块：\textbf{建筑生成器}分析原始建筑中基础资产的排列模式（层内和层间），通过GPT-4o将原始配置转换为含可重复和可缩放元素的正则化程序化代码，然后按指定尺寸生成新建筑——每个基础资产随机分配方差资产以增强多样性，且可跨场景混合不同建筑的基础资产创造全新结构；\textbf{城市布局生成器}让用户选择城市地图边界和主干道端点，系统自动划分街区、生成次级道路、确定建筑位置和尺寸，并在道路两侧放置装饰元素（路灯、垃圾桶、邮箱等，同样由3DGS采集）。生成的场景可直接导入Houdini进行进一步编辑和渲染。

\paragraph{实验评估}
综合来看，Proc-GS 在 MatrixBuilding 数据集上渲染质量（PSNR/SSIM/LPIPS）与标准 3DGS 基本持平，而模型体积减少约 4 倍，共享基础资产的设计是压缩的关键来源。在稀疏视角场景下，Proc-GS 展现出显著优势：当训练视角数量减少时，3DGS 渲染质量急剧下降，而 Proc-GS 保持鲁棒，因为重复基础资产的同步更新天然构成数据增强。在城市生成质量上，Proc-GS 在深度误差（DE）和相机误差（CE）指标上均优于 CityDreamer 等方法，验证了生成场景的 3D 一致性和几何准确性。该方法在三个真实无人机拍摄场景上同样有效，进一步验证了从真实建筑中提取资产并装配虚拟城市的可行性与跨场景泛化能力。

\subsubsection{特殊应用场景}

本路线介绍面向特殊应用场景的3DGS生成方法，包括3D物体内部纹理生成、生成式蒸馏中的自适应密度控制以及渐进式缩放生成等。

以下介绍 3 个面向特殊应用场景的 3DGS 生成方法：3D 物体内部纹理生成（FruitNinja）、生成式蒸馏中的自适应密度控制（CAdam）以及渐进式缩放生成（GaussianZoom）。

\paragraph{3D 物体内部纹理生成}

FruitNinja（2024.11）~\cite{2411_12089}提出了一个此前未被研究的问题：当3D物体被切片或切割时，内部纹理如何生成？例如，一颗西瓜被切开后应露出红瓤黑籽，一块木头被锯断后应呈现年轮纹理——这些内部视觉细节在传统3D建模中从未被系统性地处理过。核心困难在于，没有任何现有数据集能够捕获物体的完整内部结构（物体内部是不可见的），因此无法通过数据驱动方法学习内部纹理，生成方法成为自然选择。FruitNinja为此设计了一套完整的内部纹理生成管线，包含三个核心步骤。

\textbf{第一步：内部高斯基元初始化。}方法的起点是一个仅包含外表面的3DGS模型，物体内部完全为空。为了识别空腔区域，方法首先构建一个连续3D不透明度场：

\begin{equation}
d(\mathbf{x}) = \sum_p \alpha_p \exp\!\left(-\frac{1}{2}(\mathbf{x}-\mathbf{x}_p)^T\Sigma_p^{-1}(\mathbf{x}-\mathbf{x}_p)\right),
\end{equation}

即将所有高斯基元在空间点 $\mathbf{x}$ 处的贡献求和。该场被离散化到一个3D体素网格上，不透明度低于预设阈值 $\tau_{th}$ 的体素被标记为内部候选。对每个候选体素，方法沿六条主轴方向投射射线，检测从低不透明度到高不透明度的过渡，从而精确定位内壁边界。每个被确认为空腔的体素随后被填入一组预设数量的高斯基元——这些基元具有均匀的颜色和不透明度，协方差矩阵设为球形且尺度随机分配（上限为物体整体尺寸的 $10^{-3}$ 倍），确保初始填充既密集又紧凑。

\textbf{第二步：条件化截面修复与渐进纹理细化。}由于无法穷举所有切割角度的截面数据，FruitNinja利用用户指定的一组标准切割平面（定义为 $ax + by + cz + d = 0$）来高效生成内部纹理。用户指定的截面被按物体内在特征分类（例如西瓜的水平切与竖直切），每类截面施加不同的文本条件引导扩散模型。

具体而言，截面修复采用两阶段优化策略。

\textbf{第一阶段}：对每个切割角度的渲染截面视角 $V_i$，独立进行SDS优化。由于初始阶段内部高斯未经训练，渲染视角缺乏显著特征，直接SDS效率极低。因此先利用深度条件化的Stable Diffusion模型（深度图由预训练深度估计器MiDaS生成）配合切割角度特定的文本提示（如"a horizontal cross-sectional view of a watermelon"），为每个截面独立生成参考图像 $I_{label}^p$。SDS损失为：

\begin{equation}
\mathcal{L}_{\mathrm{SDS}} = \mathbb{E}_{t,\epsilon}\left[w(t)\left(\hat{\epsilon}_\phi(z_t; t, e, d) - \epsilon\right)\right],
\end{equation}

其中 $z_t = \alpha_t I_{label}^p + \sigma_t \epsilon$，$\epsilon \sim \mathcal{N}(0,I)$ 为添加的高斯噪声，$\hat{\epsilon}_\phi$ 为扩散模型预测的噪声，$e$ 为文本嵌入，$d$ 为对应截面的深度图，$w(t)$ 为时间步加权函数。

\textbf{第二阶段}：利用第一阶段生成的参考图像通过重建损失更新3DGS参数：

\begin{equation}
\mathcal{L}_{\mathrm{recon}} = \lambda\,\mathcal{L}_{\mathrm{MSE}} + (1-\lambda)\,\mathcal{L}_{\mathrm{SSIM}},
\end{equation}

即在渲染图像 $I_{RGB}^p$ 与参考图像 $I_{REF}^p$ 之间同时优化像素级均方误差和结构相似性损失。为防止截面视角优化破坏物体外表面，方法在每轮训练中随机选取10--20个外表面视角与截面视角联合训练。此外，由于Stable Diffusion模型对截面图像这类稀缺数据的生成质量有限，方法可选地通过DreamBooth微调扩散模型——仅需1--6张截面图像即可显著提升生成质量。

然而，各截面独立优化会引入空间不一致性。例如，西瓜模型的竖直切面可能在某位置显示一颗黑籽，而通过同一位置的水平切面却只显示红色果肉——这是因为两个切面的参考图像是独立生成的，无法保证3D空间中的跨截面一致性。为解决这一问题，FruitNinja采用渐进式迭代细化：每轮迭代中先用当前3DGS模型渲染所有截面视角，再对每个参考截面图做若干步SDS优化，然后重新训练3DGS，直至所有截面的重建损失收敛到预设阈值 $\theta$ 以下。这一过程使3DGS模型与扩散模型在迭代中逐步达成空间一致。

\textbf{体素网格平滑。}由于输入的截面切片是离散的，部分高斯基元可能未被任何截面遮罩覆盖，处于未训练状态，暴露时降低视觉质量。为避免额外的优化开销，方法在3DGS上构建体素网格，按固定间隔（如每30--40次迭代）执行平滑操作：未训练高斯的颜色通过同一体素内已训练高斯的距离加权平均来填充：

\begin{equation}
C = \frac{\sum_i w_i C_i}{\sum_i w_i}, \quad w_i = \frac{1}{d_i},
\end{equation}

其中 $C_i$ 为已训练高斯的颜色，$d_i$ 为未训练高斯到已训练高斯的欧氏距离。网格分辨率按经验选取，确保每个体素包含的粒子数不超过总粒子数的1\%，在保持精细纹理细节的同时促进全局颜色一致性。

\textbf{第三步：OpaqueAtom GS——原子级不透明高斯策略。}标准3DGS在训练中倾向于优化出更大的高斯基元以快速降低损失，但在内部纹理场景中，过大的高斯会跨越多个截面区域，迫使模型在不同视角间折中表达，导致训练低效甚至不收敛。此外，大高斯缺乏表达精细纹理所需的空间分辨率，也无法支持精确的几何编辑（如切开一个巨大的高斯在物理上无意义）。

OpaqueAtom GS通过两个约束来解决这一问题。

\textbf{原子裁剪（Atomic Clipping）}：受AtomGS启发，将每个高斯的尺度上限限制为物体尺寸的极小比例（如1/3000），迫使模型用密集分布的小高斯表达纹理，确保空间分辨率和几何可编辑性。

\textbf{均匀不透明化（Uniform Opacification）}：将所有高斯基元的不透明度设为满值。标准3DGS的$\alpha$-混合机制会沿每条射线叠加所有高斯的贡献，包括表面后方被遮挡的粒子——这在静态渲染中可以接受，但在交互式切片场景中，几何编辑动态改变了高斯粒子的前后排序和可见性，后方粒子可能意外影响前景渲染，产生不真实的颜色分布。全不透明设置确保每个高斯粒子忠实反映其所代表材料的真实外观，使实时切片和渲染成为可能。

\paragraph{上下文自适应密度控制}

CAdam（2026.05）~\cite{2605_20872}专门分析生成蒸馏（如SDS优化）场景下的高斯密度控制问题。标准3DGS的基于梯度幅值累积的密度控制在随机生成引导下会失效——根本原因在于重建优化中梯度随时间衰减（$\|\nabla_t \mathcal{L}\| \to 0$），但生成蒸馏中伪目标 $\hat{x}_0^t$ 每步都在变化，产生"之字形"（zigzag）现象：梯度 = 几何漂移 $\delta_t$ + 随机噪声 $\epsilon_t$，噪声分量 $\epsilon_t$ 的期望为零但幅值不衰减，导致过度致密化或欠拟合的"致密化困境"。CAdam将密度控制重新解释为统计信号验证问题：利用梯度的一阶矩（指数移动平均）实现随机波动的相消干涉（$\mathbb{E}[\epsilon_t]=0$，多次平均后趋零）和几何漂移的相长干涉（$\delta_t$ 方向一致，多次平均后增强），将底层信号从生成噪声中分离。分位数感知和信噪比门控机制确保跨优化阶段的稳健适应。CAdam将高斯数量减少85\%--97\%，同时保持可比较的感知质量。

\paragraph{Sketch2MinSurf: 草图到可编辑极小曲面生成}

Sketch2MinSurf（2026.05）~\cite{2605_20733} 利用视觉-语言模型引导从草图生成可编辑极小曲面。方法将建筑手绘草图转换为3D极小曲面（平均曲率 $H=0$ 的曲面），包含三个核心组件：(1) 拓扑感知骨架编码器将极小曲面分解为鞍面区域（负高斯曲率）和柱面表面，通过实边/虚边双层机制描述拓扑连接；(2) 微调VLM检测草图并输出序列化骨架描述（节点坐标、维度、连接性）；(3) 解码器通过多管算子和曲率松弛重建3D流形几何。方法构建了3000样本的多模态数据集，在零样本设置下通用VLM（GPT-4.1、Claude Sonnet 4等）无法产生拓扑有效输出，验证了结构化编码和领域微调的必要性。生成的曲面可直接在Grasshopper等参数化设计环境中编辑，已成功应用于实际建筑装置的端到端建造。

\paragraph{FingerSplat: 3DGS指纹识别重建与合成}

FingerSplat（2025.09）~\cite{2509_15648} 首次将3DGS引入指纹识别领域，实现了非接触式指纹的3D重建、配准和合成。方法包含两阶段3D对齐：(1) 基于DUSt3R的成对局部对齐，直接在3D点云空间匹配指纹特征（而非传统2D空间配准）；(2) 全局优化将多视角对齐结果融合为完整3D指纹。重建后的3D指纹通过3DGS渲染生成高保真合成指纹，即使在未知手指姿态和稀疏输入图像条件下仍能有效运行。合成的指纹数据可用于训练指纹识别网络，有效提升非接触式指纹识别性能，缓解了生物识别领域训练数据不足的难题。

\begin{takeaway}
其他3D生成方法展示了3DGS在专业领域的生成应用。Sketch2MinSurf 将建筑草图转换为可编辑极小曲面，FingerSplat 将3DGS引入指纹识别实现非接触式3D重建。这些工作表明3DGS的生成能力不限于通用物体，在建筑设计和生物识别等专业领域同样具有应用价值。
\end{takeaway}

%% ============================================================

\section{近期研究（2026 年 8--11 月）}
\paragraph{Scene2Sound：将物体级音源绑定到可漫游的高斯场景}
2026.08~\cite{2608_00463}。Scene2Sound 面向已有的 3DGS 场景，先选取覆盖场景的观察位姿，使用视觉语言模型区分不需要定位的环境背景声与可定位的发声物体。模型将不同视图的物体掩码关联到高斯集合，再依据集合重合匹配三维实例，生成的音源交由对象式音频引擎按照听者位置渲染。

作者以 24 个生成的 SoundscapePLY 场景进行主要评测，并在由真实 360 度影像重建的 D-SAV360 场景检查迁移表现。19 位参与者以移动视点比较三种空间音频基线。论文的失败案例显示，如果分割阶段无法定位发声物体，系统可能只能合成无明确方位的背景声。

\paragraph{GS-Voxel：为大范围航拍场景构造结构化高斯潜变量}
2026.08~\cite{2608_17988}。GS-Voxel 将预优化且满足输入约定的高斯集合确定性组织为稀疏体素。因子化 VAE 分别编码体素占据和局部高斯属性，图像条件流模型在潜空间生成，带重叠的平铺推理将合成范围拓展到单块训练裁剪之外。这里避免的是逐场景向固定大小模板拟合，训练资料仍来自已重建的 3DGS 场景。

论文在各含 1.5 万张渲染图像的生成与留出参照集合上报告 FID 28.0、KID 0.020，并明确区分不同论文的相机协议和数据。当前实现只在 SH0 航拍场景开发和验证，改变卫星影像来源后的鲁棒性尚未评估。极密体素的容量限制和有限空间分辨率也可能分别造成原语丢弃与细结构退化。

\paragraph{SPAR3S：在稀疏体素潜空间自回归补全场景}
2026.09~\cite{2609_03931}。SPAR3S 将稀疏条件图像编码到与高斯解码器对齐的体素潜空间。掩码自回归模型在其中联合预测新增占据和潜变量，随后解码出可供新视点渲染的高斯集合。研究针对条件视图很少、彼此甚至没有重叠的室内场景重建。

论文在合成室内 3DFront 与真实视频 RealEstate10K 上评估，前者提供真值位姿并从渲染深度回投点图，后者由 MASt3R-SfM 估计初始几何。作者指出占据预测漏掉墙壁或物体体素时会留下空洞，有限体素分辨率和每体素的高斯密度也会引入局部模糊。

\paragraph{Filling the Unseen：分阶段视角外推与质量感知监督}
2026.09~\cite{2609_13262}。该方法先检查新相机视图的重叠和依赖，把可独立外推的相机视图分组并行处理，再补全依赖视图。质量感知掩码以视线方向相关的高斯渲染质量估计调节损失权重，使生成监督用于缺失或质量较低的区域，同时减少对可靠观测区域的干扰。

论文在 Mip-NeRF 360 的 4 个场景采用无真值代理指标，在 ScanNet++ 的 10 个场景将一半区域留作外推目标并以原影像评估。两套协议分别衡量感知代理和有图像参照的受控外推表现；在真实未知区域，代理指标不能证明生成内容与未观测事实吻合。

\paragraph{PhGS：对单视图前馈高斯先裁剪再细化}
2026.09~\cite{2609_20623}。PhGS 承接每像素预测多个原语的单视图前馈骨干。它先依重要性从预测的高斯集合裁剪，再用学习得到的循环模块直接更新保留原语，避免单纯删除原语带来的外观或几何损伤。混合不同保留比例训练让一个细化器可在推理时适配不同保留数目。

作者在 RealEstate10K 上训练和评测，并以 KITTI 和 DL3DV 做零样本跨域评测。在 Flash3D、Niagara、SHARP 三种骨干的域内实验中，主设置可舍弃一半预测高斯并保留有竞争力的渲染表现。细化模块仍需训练，跨域结论限于论文测试的骨干和数据。

\paragraph{GenNVS：把前景背景分别建模，再用三维先验约束视频外推}
2026.09~\cite{2609_34579}。单张图像外推新视图时，前景和背景不一定能由统一变形解释。GenNVS 先分别建立前景和背景，再对前景高斯作粗到细配准。从粗三维场景得到的重投影视图和有效掩码通过双流输入控制视频扩散，使它能区分有几何支撑的区域和需要补全的区域。

基线包含前馈与生成式新视图方法，原文在 DL3DV 和 Mip-NeRF 360 报告渲染与姿态一致性对照，并演示场景编辑。训练依赖配对视频资料和高算力，推理设置为 50 步去噪。补全区域的视觉稳定性仍不能作为场景真实结构的证据。

\paragraph{Edit4D：将静态背景与动态人物分开生成和编辑}
2026.10~\cite{2610_21668}。Edit4D 从文本分别生成高斯背景和动态人物，在统一场景中通过目标区域定位与接触约束完成组合。组件独立表示使替换人物外观或调整人物落点时可以保留已有背景。模型还可渲染改变相机位姿后的动态场景。

论文沿用 Free4D 的正面视频协议评估文本对齐、时序一致性和美学质量，并以人物外观编辑和场景重组对比必须全局重生成的 Free4D。作者记录了椅子被桌子遮挡时人物落点不准、与桌面几何相交的案例，因而放置的可靠性依赖目标可见程度。

\paragraph{TriFold：用规则三平面编码不规则高斯场景}
2026.10~\cite{2610_23050}。TriFold 将带有空间密度差异和可变原语数量的场景高斯集合编码为固定布局的三平面变分潜变量。编码采用单元查询及以原语为锚的跨平面注意力，解码器再按占据与局部原语置信预测分配高斯。规则潜空间也供条件流匹配生成使用。

论文在 InteriorGS 与 SceneSplat-49K 的 Aria Synthetic Environments 子集验证重建，按原场景划分训练与测试以避免裁剪泄漏。论文另对点云条件和单图全形状补全展示生成结果，并指出观察稀疏的场景边界仍可能较真值模糊。

\paragraph{Plant Orbiter：用相机轨迹已知的生成视频构造植物高斯}
2026.11~\cite{2611_12739}。Plant Orbiter 由单张植物图像合成 73 帧、围绕静止植物的短视频，借预定轨道估计相机位姿并拟合可对齐输入的高斯模型。前接文生图模型后可用同一流程生成文本指定的植物资产。

论文报告各测试集的最佳几何 F-score，并在由两株植物、17 个时间点组成的 GrowSplat 样本上用点集测量植株高度，结论报告平均高度误差 5.9\%。单视图看不到的背侧和遮挡部位由模型推断，固定高度单轨环绕也会造成顶视或底视伪影，限制其用于强自遮挡结构的表型测量。

\paragraph{VideoRNG：用标准化视频观测生成可重照明的神经高斯}
2026.11~\cite{2611_25338}。VideoRNG 在规定的转台相机与照明下从文本或图像合成观察序列，将 29 帧编码至 8 个视频潜帧，前景特征经重建网络一次预测高斯几何与神经外观。几何与外观的组织允许在新观察方向和新光照下渲染资产。

作者用约 2.2 万个筛选后的 TexVerse 纹理资产训练，并在留出的视点和灯光条件下评估重建。论文深度锚定对比的 PSNR 为 29.22 dB，联合深度监督为 25.39 dB，其含义限于该比较。生成视频的跨帧不一致可在重建中留下重影，尖锐镜面效果仍难复现。


\section*{本章小结}
\addcontentsline{toc}{section}{本章小结}

\begin{futurebox}{本章总结与展望}{ch08_summary}

本章系统梳理了基于3D高斯泼溅的文本/图像到3D生成技术，涵盖四条主要技术路线和多类辅助方法。

\textbf{分数蒸馏方法}（第~\ref{sec:08_1}~节）是最早将2D扩散先验引入3D生成的路线。SDS及其变体（VSD、ISM、CSD等）通过将2D扩散模型的梯度"蒸馏"到3D表示中，无需3D训练数据即可生成3D内容。3DGS作为优化目标，凭借其高效的梯度传播和快速的渲染反馈，将文本到3D的生成时间从数小时压缩到数分钟。然而，SDS类方法本质上是一种逐样本优化，每次生成都需要完整的优化循环，限制了其在大规模部署中的效率。

\textbf{图像到3D方法}（第~\ref{sec:08_2}~节）通过前馈网络将单张图像直接映射为3D高斯。大规模重建模型（LRM）利用Transformer架构和多视角扩散先验，在秒级时间内生成高质量的3DGS。这些方法从重建走向生成，在保持效率的同时逐步缩小了与优化方法的品质差距。

\textbf{3D生成模型}（第~\ref{sec:08_4}~节）代表了从"优化生成"到"采样生成"的范式转变。GAN方法将经典3D感知对抗训练迁移到高斯空间；扩散模型直接在3DGS或其潜在表示上做去噪；自回归方法通过多尺度预测逐步构建高斯层级；VAE方法将高斯压缩到潜在空间后做生成。这条路线的核心挑战在于3D高斯的离散无序特性——UVGS通过球面映射将其结构化，DiffGS通过函数化表示将其连续化，DeG通过可学习密度控制将其自适应化。当前这条路线正在快速逼近2D生成模型的质量和多样性。

\textbf{4D动态生成}（第~\ref{sec:08_3}~节）将静态3D生成扩展到时间维度，同时建模空间外观和时序运动。以DreamGaussian4D为代表的SDS+形变场路线通过HexPlane分解压缩4D表示，以STP4D为代表的直接生成路线通过扩散模型直接输出4DGS。运动迁移与动画生成方法（如CharacterShot）将静态3D资产与姿态序列结合，生成可控的角色动画。4D生成的核心挑战在于同时保证空间一致性和时间连贯性。

\textbf{其他方法}（第~\ref{sec:08_5}~节）覆盖了从点云到高斯的转换（GAP、GaussianGrow）、大规模场景生成（EarthCrafter、GaussianCity、Proc-GS）、内部纹理生成（FruitNinja）等特殊方向，展现了3DGS生成技术在多样化应用场景中的潜力。

\textbf{开放问题与未来方向}：

\begin{enumerate}
\item \textbf{生成质量与效率的平衡}：当前最先进的3DGS生成方法在视觉质量上仍落后于2D生成模型，而高保真方法（如SDS优化）在效率上无法满足实时应用需求。如何在不牺牲质量的前提下实现一步到位的快速生成，是核心挑战。

\item \textbf{场景级生成的全局一致性}：物体级3D生成已取得显著进展，但场景级生成（多物体交互、复杂布局、大范围导航）仍面临全局一致性难题。从全景图到可探索3D场景的管线中，累积误差导致远距离视角出现伪影。

\item \textbf{物理一致性与可控性}：生成结果的物理合理性（支撑关系、碰撞检测、材质属性）仍需更多关注。当前大多数方法只关注视觉外观，忽略了生成内容在物理模拟和交互场景中的可用性。

\item \textbf{3D训练数据的规模与质量}：3D生成模型的性能在很大程度上受限于3D训练数据的规模和标注质量。与2D领域的互联网规模数据相比，高质量3D资产数据集仍显不足，如何利用合成数据和自监督方法弥补这一差距是重要方向。

\item \textbf{统一生成框架}：当前不同3D生成任务（文本到3D、图像到3D、4D动态生成）往往需要专门的方法和管线。统一的生成框架能否覆盖所有任务，同时保持各任务上的最优性能，是一个值得探索的方向。
\end{enumerate}

\end{futurebox}
